% concatenated from main.tex
\documentclass[review,11pt]{article}

\usepackage[utf8]{inputenc}
\usepackage[english]{babel}
\usepackage{placeins}
\usepackage[a4paper, left=1in, right=1in, top=1in, bottom=1in]{geometry}

\usepackage{amsmath,amssymb}
\usepackage{amsthm}
\usepackage{graphicx}
\usepackage{multirow}
\usepackage{enumitem}
\usepackage{float}
\usepackage{color}
\usepackage{xcolor}
\usepackage{subcaption}
\usepackage{adjustbox}
\usepackage{cancel}
\usepackage{ulem}
\usepackage{mathtools}
\usepackage{comment}

% Important: aliascnt BEFORE hyperref/cleveref
\usepackage{aliascnt}
\usepackage{hyperref}
\usepackage{cleveref}

\usepackage{multicol}
\usepackage{booktabs}
\usepackage{titlesec}
\usepackage{algorithm}
%\usepackage{algorithmic}
\usepackage{algpseudocode}
\usepackage{authblk}
\usepackage{tikz}
\usepackage{stmaryrd}
\usepackage{mathpazo}
\usepackage{bm}
\usepackage{microtype}
\usepackage{setspace}
\usepackage{tcolorbox}
\usepackage{wrapfig}
\tcbuselibrary{breakable}
\usepackage{longtable}

\onehalfspacing

\usepackage{fancyhdr}
\pagestyle{fancy}
\fancyhf{}
\renewcommand{\headrulewidth}{0pt}
\fancyfoot[C]{\small \thepage}
\fancyhead[L]{\small \textit{Space-time iterative scheme and adaptivity}}
\fancyhead[R]{\small Javed, Mitra, Pop}

\usetikzlibrary{arrows.meta,positioning,shapes.geometric}

% Custom subsubsubsection numbering
\setcounter{secnumdepth}{4}
\titleclass{\subsubsubsection}{straight}[\subsection]
\newcounter{subsubsubsection}
\titleformat{\subsubsubsection}
  {\normalfont\normalsize\bfseries}
  {\thesubsubsubsection}
  {1em}
  {}
\titlespacing*{\subsubsubsection}
  {0pt}
  {3.25ex plus 1ex minus .2ex}
  {1.5ex plus .2ex}

\renewcommand{\thesubsubsubsection}
  {\thesubsection.\arabic{subsubsection}.\arabic{subsubsubsection}}

\newcounter{AssDD}
\stepcounter{AssDD}

\newcounter{AssM}
\stepcounter{AssM}

\numberwithin{equation}{section}
\providecommand{\MR}[1]{MR #1}

% Title and authors
\title{An adaptive, space-time discretized linear iterative scheme for doubly-degenerate parabolic problems}

\author[1]{A. Javed}
\author[2]{K. Mitra}
\author[1]{I.S. Pop}

\affil[1]{Hasselt University, Belgium}
\affil[2]{Eindhoven University of Technology, The Netherlands}

\date{\today}

% Short commands
\def \a  {\alpha}
\def \oa {\overline{\alpha}}
\def \b  {\beta}
\def \g  {\gamma}
\def \G  {\Gamma}
\def \d  {\delta}
\def \D  {\Delta}
\def \e  {\varepsilon}
\def \eps {\epsilon}
\def \f  {\varphi}
\def \j {\ell}
\def \vr  {\varrho}
\def \vs  {\omega}
\def \k  {\kappa}
\def \l  {\lambda}
\def \om {\omega}
\def \oma {\omega_a}

\newcommand{\tsl}[1]{\bar{#1}^{\mathrm{t}}_a}

\def \omt {\tsl{\om}}
\def \Om {\Omega}
\def \r  {\rho}
\def \s  {\sigma}
\def \t  {\tau}
\def \th {\theta}
\def \z  {\zeta}
\def \del {\nabla}
\def \p  {\partial}
\def \N  {{\mathbb{N}}}
\def \R  {{\mathbb{R}}}
\def \dd {\mathrm{d}}
\def \m {\mathrm{m}}
\def \x {{\bm{x}}}
\def \htt {{h\tau}}

\def \calA {\mathfrak{A}}
\def \calB {\mathfrak{B}}
\def \calE {\mathcal{E}}
\def \calF {\mathfrak{F}}
\def \calZ {\mathcal{Z}}
\def \calV {\mathcal{V}}
\def \calN {\mathcal{N}}
\def \calP {\mathcal{P}}
\def \calQ {\bm{\mathcal{Q}}}
\def \calR {\mathcal{R}}
\def \resH {{\calR}_{\mathrm{H}}}
\def \calS {\mathcal{S}}
\def \calT {\mathcal{T}}
\def \calW {\mathcal{W}}
\def \fli {\bm{\varsigma}^{i-1}_n}
\def \flx {\bm{\sigma}_{\htt}^i}
\def \flxa {\bm{\sigma}^a_{\htt}}
\def \xt {\underline{\bm{x}}_{\mathrm{t}}}
\def \nhat {\bm{n}_{\mathrm{x}}}
\def \Shti {\calS_\htt^i}
\def \Fhti {\bm{F}_\htt^i}

\def \Maxi {{\mathrm{M}}}
\def \Mini {{\mathrm{m}}}

\newcommand{\linErr}[1]{\calE_{\!_{\,{\rm lin},n}}^{#1}}
\newcommand{\linEst}[1]{\eta_{\!_{\,{\rm lin},n}}^{#1}}
\newcommand{\normErr}[1]{\tilde{\calE}_{\!_{\,{\rm lin},n}}^{\,#1}}
\newcommand{\normEst}[1]{{\eta}_{\!_{\,{\rm lin},n}}^{\,#1}}

\newcommand{\lp}[3]{\big(\!\big( #2,\,#3\big)\!\big)_{#1}}

\newcommand{\norm}[1]{{
  \left\vert\kern-0.25ex\left\vert\kern-0.25ex\left\vert #1
  \right\vert\kern-0.25ex\right\vert\kern-0.25ex\right\vert
}}

\newcommand{\efl}[1]{\bm{\sigma}_{n,\calT}^{#1}}

\newcommand{\snorm}[1]{{
  \left[\kern-0.25ex\left[ #1
  \right]\kern-0.25ex\right]
}}

% Theorem-like environments
% This fixes cleveref calling everything "Lemma".
\newtheorem{lemma}{Lemma}[section]

\newaliascnt{theorem}{lemma}
\newtheorem{theorem}[theorem]{Theorem}
\aliascntresetthe{theorem}

\newaliascnt{corollary}{lemma}
\newtheorem{corollary}[corollary]{Corollary}
\aliascntresetthe{corollary}

\newaliascnt{proposition}{lemma}
\newtheorem{proposition}[proposition]{Proposition}
\aliascntresetthe{proposition}

\newaliascnt{assumption}{lemma}
\newtheorem{assumption}[assumption]{Assumption}
\aliascntresetthe{assumption}
\renewcommand{\theassumption}{A.\arabic{assumption}}

\newaliascnt{definition}{lemma}
\newtheorem{definition}[definition]{Definition}
\aliascntresetthe{definition}

\newaliascnt{remark}{lemma}
\newtheorem{remark}[remark]{Remark}
\aliascntresetthe{remark}

\newaliascnt{example}{lemma}
\newtheorem{example}[example]{Example}
\aliascntresetthe{example}

\newaliascnt{conjecture}{lemma}
\newtheorem{conjecture}[conjecture]{Conjecture}
\aliascntresetthe{conjecture}

\newtheorem{problem}{Problem}

% Names for cleveref
\crefname{lemma}{Lemma}{Lemmas}
\Crefname{lemma}{Lemma}{Lemmas}

\crefname{theorem}{Theorem}{Theorems}
\Crefname{theorem}{Theorem}{Theorems}

\crefname{corollary}{Corollary}{Corollaries}
\Crefname{corollary}{Corollary}{Corollaries}

\crefname{proposition}{Proposition}{Propositions}
\Crefname{proposition}{Proposition}{Propositions}

\crefname{assumption}{Assumption}{Assumptions}
\Crefname{assumption}{Assumption}{Assumptions}

\crefname{definition}{Definition}{Definitions}
\Crefname{definition}{Definition}{Definitions}

\crefname{remark}{Remark}{Remarks}
\Crefname{remark}{Remark}{Remarks}

\crefname{example}{Example}{Examples}
\Crefname{example}{Example}{Examples}

\crefname{conjecture}{Conjecture}{Conjectures}
\Crefname{conjecture}{Conjecture}{Conjectures}

\crefname{problem}{Problem}{Problems}
\Crefname{problem}{Problem}{Problems}

\crefname{subsubsubsection}{section}{sections}
\Crefname{subsubsubsection}{Section}{Sections}

% Box
\newtcolorbox{mybox}[1]{
    breakable,
    colback=yellow!10!white,
    colframe=red!75!black,
    fonttitle=\bfseries,
    title={#1}
}

\newcounter{Ass}
\stepcounter{Ass}

\newcommand\KM[1]{%
  \protect\leavevmode
  \begingroup
        \color{red!55!yellow}%
        #1%
  \endgroup
}

\newcommand\AJ[1]{%
  \protect\leavevmode
  \begingroup
        \color{red!30!blue}%
        #1%
  \endgroup
}

\begin{document}

\maketitle
\begin{abstract}
Degenerate diffusion problems, where the governing parabolic equation can change type to either an ordinary differential equation or an elliptic equation, model many real life applications. Due to the presence of free-boundaries, accurate numerical simulation of such problems require extremely small mesh and time step sizes locally. To remediate this issue, in this work, we consider a space-time formulation of the problem based on an efficient splitting of the nonlinearities.  First, an iterative linearization scheme is proposed to resolve the nonlinearities that effectively reduces to solving a sequence of heat equations. Unconditional convergence of the scheme is proven even for double degenerate cases with linear convergence achieved if the problem is non-degenerate. Next,  the dual norm of the nonlinear residual is decomposed into a linearization error component and a discretization error component corresponding to the heat equation. This leads to reliable and fully computable a posteriori estimates for the problem that are robust with respect to the nonlinearities/degeneracies. These estimates are used then in a fully adaptive (discretization + linearization) space-time solver. Numerical experiments for multiple test cases (one and two dimensions in space) demonstrate that this solver efficiently allocates the computational resources in the space-time domain, resulting in a rapid decay of error in terms of total degrees of freedom spent.


% Using a fully space–time variational formulation, we develop a numerical approximation scheme for a class of nonlinear degenerate parabolic equations. In this approach, the spatial and temporal discretization are addressed simultaneously, allowing for a flexible mesh in space and time, but leading to large fully discrete systems. 
% %Additional challenges are encountered in cases when the diffusion is nonlinear, or even degenerate.
% The scheme relies on a splitting strategy, in which the nonlinear dependencies are considered as algebraic equations, at the expense of new unknowns. This is applied directly to the space–time formulation, leading to an equivalent system that is more suitable for handling degeneracies. Based on this reformulation, we employ a stabilized iterative linearization scheme and establish its convergence in the space–time setting.
% A central contribution of this work is the development of guaranteed a posteriori error estimators, based on locally equilibrated flux reconstructions on space–time vertex patches. The resulting estimators are fully computable and provide a clear separation between linearization and discretization errors. This enables a reliable adaptive strategy in which the mesh refinement is activated only when the discretization error dominates. In this way, the global number of unknowns is reduced significantly when compared to non-adaptive schemes. Numerical experiments demonstrate that the proposed adaptive space–time method efficiently concentrates the computational effort in regions exhibiting strong space–time variations, yielding accurate solutions at significantly reduced computational cost.
\end{abstract}
\section{Introduction} \label{introduction}
Degenerate parabolic equations arise in a wide range of applications, such as flow through porous media \cite{vazquez2007porous}, mathematical biology \cite{van2002mathematical}, phase transition \cite{Malgo}, traffic flow \cite{berthelin2008model}, and other nonlinear diffusion processes. 
%Typical examples include degenerate diffusion models \cite{pop2011regularization,ZouExistence,Malgo,van2002mathematical}, where the degeneracy appears both in the time derivative and the diffusion operator. 
The solution to such models features free boundaries (where the equation becomes an ordinary differential equation(ODE)) with sharp and moving interfaces, causing the solution to have low space regularity, along with degenerate regions (where the equation becomes elliptic) with low time regularity. These two "degeneracies" make it challenging to numerically approximate the solutions. They also render uniform discretizations (particularly, in time) inefficient since the sharp interfaces often require very fine meshes.

Let $\Om$ be a bounded open domain in $\R^{d}$ having a Lipschitz boundary $\p\Om$. For some final time $T>0$, let $Q:=(0,T]\times\Omega$ and $\Gamma:=(0,T]\times \p\Om$ be the space-time domain and interface respectively. Then the degenerate diffusion problem reads: find $u,\,w:Q\to \R$ satisfying
\begin{equation}\label{eq:parabolic_equation}
\left\{ \begin{array}{rcll}
\partial_t u &=& \Delta w + f, & \text{in } Q, \\
w &\in& \Phi(u), & \text{in } Q, \\
w &=& 0, & \text{on } \Gamma, \\
u(0) &=& u_0, & \text{in } \Omega,
\end{array}
\right.
\end{equation}
Here, $\Phi:[0,\omega)\to \R$ denotes a nonlinear, monotonically increasing mapping for either $\omega=1$ or $\infty$. In particular, $\Phi' $ might vanish at $0$, making the problem an ODE (see Figure~\ref{fig:RichVsBio_phi}). On the other hand, $\Phi'$ might blow up at $\om=1$ (for instance, for biofilm growth \cite{van2002mathematical} and traffic flow models \cite{berthelin2008model}) or at infinity, as in the porous medium equation (PME) case. It might even become multivalued at $\om=1$, like in the Richards equation case \cite{mitra2024posteriori}. These scenarios set off the elliptic degeneracy.

\begin{figure}[h]
    \centering
    \begin{subfigure}[b]{0.32\textwidth}
        \centering
    \includegraphics[width=\textwidth, height=4.1cm]{Figures/RichardsDegeneracy.eps}
        \caption*{$\Phi$ becomes multivalued at $u = 1$}
    \end{subfigure}
    \begin{subfigure}[b]{0.32\textwidth}
        \centering
        \includegraphics[width=\textwidth, height=4.1cm]{Figures/BiofilmDegeneracy.eps}
        \caption*{$\Phi$ becomes infinite at $u = 1$}
    \end{subfigure}
      \begin{subfigure}[b]{0.32\textwidth}
        \centering
        \includegraphics[width=\textwidth, height=4.1cm]{Figures/PME_Degeneracy.eps}
        \caption*{degeneracy occurs at $u = 0$}
    \end{subfigure}
    \caption{Examples of degenerate nonlinearities $\Phi$: (\textbf{Left}) Richards-type \cite{mitra2024posteriori}, (\textbf{Center}) biofilm-type \cite{van2002mathematical,berthelin2008model}, and (\textbf{Right}) porous media equation-type \cite{vazquez2007porous} degeneracies.}
\label{fig:RichVsBio_phi}
\end{figure}
To deal with the double degeneracy discussed above, we use the framework discussed in~\cite{brenner2017improving} and also see in~\cite{javed2025robust}. The system is reformulated in terms of a new variable $s$, and two increasing functions $b, B \in C^1(\mathbb{R})$ that satisfy
\begin{align}\label{eq:condition}
    \Phi =B \circ b^{-1}, \text{ and }\; 0\leq b',B'\leq 1, \text{ and }  b'+B'\geq 1.
\end{align}
With this choice, if $u = b(s)$ and $w=B(s)$, then one immediately gets $w \in \Phi(u)$.  In this way, problem \eqref{eq:parabolic_equation} becomes of finding $s:Q\to \R$ satisfying
\begin{equation}\label{eq:reformulation}
\left\{ \begin{array}{rcll}
\partial_t b(s) &= &\Delta B(s) + f, \qquad & \text{ in } Q, \\
B(s) &=& 0, & \text{ on } \Gamma, \\
b(s(0)) &=& u_0, &\text{ in } \Omega. \\
\end{array}
\right.
\end{equation}
The advantage of this reformulation is that the functions $b$ and $B$ are Lipschitz continuous on $\mathbb{R}$, with Lipschitz constants equal to~1.  In this setting, the two types of degeneracy arise when either $b' \searrow 0$ (corresponding to the fast diffusion (elliptic) regime in the original formulation) or $B' \searrow 0$ (corresponding to the slow diffusion (ODE) regime).
Such a decomposition is possible under relatively weak assumptions, and an explicit construction of the $b,\, B$ functions is presented in \Cref{lemma:b_B_construction}, see also \cite{javed2025robust}. The graphs of these functions are illustrated in \Cref{bB-graphs} for the profiles of \Cref{fig:RichVsBio_phi}.

\begin{figure}[h!]
    \centering
    \begin{subfigure}{0.32\textwidth}
        \centering
        \includegraphics[width=\linewidth]{Figures/RichardsReparametrization.eps}
    \end{subfigure}
    \hfill
    \begin{subfigure}{0.32\textwidth}
        \centering       \includegraphics[width=\linewidth]{Figures/BiofilmReparametrization.eps}
    \end{subfigure}
     \hfill
    \begin{subfigure}{0.32\textwidth}
        \centering       \includegraphics[width=\linewidth]{Figures/bB_PME_plot.eps}
    \end{subfigure}
\caption{The functions $b$ and $B$ for the mappings $\Phi$ appearing in different applications shown in \Cref{fig:RichVsBio_phi}, including their extensions to $s\le 0$. (Left) Richards case, (Center) biofilm case and (Right) porous medium equation case (PME).}
\label{bB-graphs}
\end{figure}


\noindent
Traditionally, the numerical treatment of such problems relies on semi-discretization techniques: either discretizing first in space \cite{MR2249024} (the method of lines), leading to a system of nonlinear ODEs in time, or discretizing first in time \cite{kavcur2006method} (Rothe’s method), solving a sequence of elliptic problems at successive time steps (time-stepping). However, features like free boundaries and saturated regions are typically confined to limited regions of the space-time domain and cannot be efficiently resolved by adaptive time-stepping or space refinements. Since their location is not known a priori, a posteriori indicators need to be considered. These considerations motivate variational space–time discretization \cite{neumuller2013space} and iterative linearization methods combined with adaptive refinements prompted by rigorous a posteriori error estimation techniques.  


%We formulate \eqref{eq:parabolic_equation} in a fully space-time variational framework, where the solution is sought on the entire space-time domain. In this sense, the time derivative acts as a strong convection term , leading to a unified treatment of spatial and temporal features of the solution.
%In recent years, there has been a rapidly growing interest in such simultaneous space-time methods for solving linear parabolic evolution equations, mostly focused on discretization aspects,  
%e.g.,~\cite{neumuller2011refinement, Andreev2013, 
%steinbach2015space, LANGER2016342, 
%GanderNeumuller2016, devaud2018space, rekatsinas2019optimal, SteinbachZank2020,StevensonWesterdiep2021}. 

\textbf{Space-time discretization:} In recent years, there has been a rapidly growing interest in simultaneous space-time methods across a range of problems, including 
linear parabolic evolution equations. Various formulations and strategies have been proposed and analyzed, including conforming methods on general meshes \cite{steinbach2015space}, refinement strategies for flexible space-time meshes \cite{neumuller2011refinement}, simplex space-time meshes enabling local temporal refinement and unstructured meshing in both space and time 
\cite{Simplex_sp}, sparse Petrov-Galerkin discretizations \cite{Andreev2013}, space-time formulations in isogeometric analysis \cite{LANGER2016342, AntoniettiDede2025}, parallel multigrid algorithms with provable scalability \cite{GanderNeumuller2016}, space-time $hp$-approximations achieving exponential convergence under analytic regularity \cite{devaud2018space}, optimal adaptive wavelet solvers of linear complexity \cite{rekatsinas2019optimal}, coercive formulations yielding symmetric positive definite systems \cite{SteinbachZank2020}, and stable mixed variational formulations \cite{StevensonWesterdiep2021}. Space-time domain decomposition methods with non-matching spatial grids and asynchronous local time stepping have been developed for parabolic problems with solution features localized in space and time \cite{jayadharan:hal-03355088}. Closer to the nonlinear setting, \cite{beranek2025quasioptimal} reformulate parabolic equations with Lipschitz continuous, strongly monotone spatial operators by introducing an auxiliary variable, obtaining an equivalent system with a Lipschitz continuous operator and inverse, and establish quasi-optimality of the resulting space-time Galerkin discretizations. The discontinuous Galerkin method has been studied in \cite{Jan2012}, see also \cite{congreve2026residual,dolejvsi2024error,DOLEJSI2019276,DolejsiRoskovec2021}. Without these exceptions, methods are predominantly focused on \emph{linear} parabolic problems, and their extension to \emph{degenerate} parabolic problems in particular remains largely unexplored. Using the iterative space-time linearization introduced below, \emph{we reduce solving the reformulated nonlinear system \eqref{eq:reformulation} to solving a sequence of heat equations}. These are then discretized by space-time conforming finite element method \cite{steinbach2015space}, which is well-posed and consistent for all iterations.

\textbf{Space-time linearization:} A central challenge in the numerical solution of nonlinear degenerate parabolic equations is the design of robust and efficient linearization schemes. \emph{To the best of our knowledge, no linearization scheme with provable convergence has been developed for double degenerate problems within a fully global space-time variational framework}. For nonlinear elliptic problems, particularly those arising from time-discretizing parabolic equations, Newton's method is the classical choice \cite{bergamaschi1999mixed,lehmann1998comparison} due to its quadratic convergence. However, its convergence is guaranteed only when the initial guess is sufficiently close to the solution, which in the degenerate case translates into severe restrictions on the time step size or regularization \cite{radu2006newton,javed2025robust,fevotte2024adaptive}. To overcome this limitation, fixed-point iterative schemes have been proposed as robust alternatives. Among these, the L-scheme provides unconditional convergence regardless of the time step size, spatial discretization, or degree of degeneracy, albeit with a slow, linear convergence rate \cite{pop2004mixed, list2016study}. Several attempts have been made to reconcile with the trade-off between the speed of the Newton scheme and the stability of the L-scheme, including a modified Newton method \cite{mitra2019modified}, a modified Picard method \cite{celia1990general}, nested Newton iterations \cite{casulli2010nested},  L-scheme with fixed preconditioners \cite{kohle2025robust}, and an adaptive L/Newton scheme \cite{stokke2023adaptive}, to mention a few. For the Richards equation specifically, 
\cite{brenner2017improving} proposed the $b-B$ parametrization strategy that significantly improves the performance of Newton's method. Building on this, \cite{javed2025robust} extended these ideas to the L/M-schemes, yielding guaranteed convergence in double degenerate cases along with asymptotic quadratic convergence for the adaptive M-version. A related approach for singular degenerate evolution systems modelling biofilm growth was developed in \cite{SmeetsEtAl2024}. 


However, all these results are restricted to the sequential time-stepping framework.  Here, we fill this gap by adopting the splitting L-scheme of \cite{javed2025robust} to space-time formulations. For an iterate $s^i$ and a constant $L>1$, the linearization finds the next iterate $s^{i+1}$ by solving a heat equation:
\begin{equation}\label{eq:Lsch-strong}
\partial_t(s^{i+1}-s^i) - \Delta(s^{i+1}- s^i) 
= -  L^{-1}[\partial_t b(s^i)-\Delta B(s^i) - f].
\end{equation}
We will show that \emph{this iterative method is unconditionally convergent even for double degenerate cases, the convergence being linear if the problem is non-degenerate}. The simple heat equation structure is used to construct a fully adaptive solver further on.

\textbf{A posteriori error estimation 
and adaptivity:} 
The global space-time approach, while powerful, comes with the high computational cost of solving and storing the entire 
space-time profile, all at once. However, this can be made computationally affordable by adaptive discretization methods, which enable local time-step refinements, unlike time-stepping. This goal is 
achieved for linear parabolic problems by space-time adaptive wavelet methods at optimal complexity \cite{schwab2009space}, but whose extension to nonlinear degenerate problems is non-trivial. Achieving adaptivity relies crucially on a posteriori error estimation. Guaranteed, locally space-time efficient, and polynomial degree robust estimates
were established for linear parabolic problems using equilibrated flux reconstructions in
\cite{ErnSmearsVohralik2017, ErnSmearsVohralik2019}, however, for the time-stepping case.  A posteriori estimators for linear and semilinear parabolic problems have also been developed in 
\cite{eriksson1991adaptive, Eriksson1995, 
schmich2008adaptivity, FUHRER202127, 
RattiVerani2019, GeorgoulisLakkisWihler2021}. 
Space-time adaptivity has been successfully applied to parabolic optimization and 
control problems in \cite{Meidner2012, gong2012space, LangerSchafelner+2022+247+266, GunzburgerKunoth2011}. 

 For nonlinear problems, a posteriori estimates driving adaptive refinement in both space and time were developed within sequential time-stepping frameworks in \cite{BernardiElAlaouiMghazli2014, 
CLEMENT2021103897}, with \cite{Dyja2018} investigating parallel-in-time framework and \cite{BARON2017104} additionally 
separating linearization errors. For the degenerate Richards equation,  in \cite{mitra2024posteriori}, guaranteed and locally space-time efficient estimators were derived that are robust with respect to the final time, and separate every error component. However, it too required sequential time-stepping and involved complex weighted norms.
Using discontinuous Galerkin (DG) methods, an adaptive strategy was proposed in \cite{DOLEJSI2019276}; reliable and efficient estimators were derived in \cite{DolejsiRoskovec2021,dolejvsi2024error}, however, in mesh and solution dependent norms; and in \cite{congreve2026residual} for mesh and nonlinearity dependent constants.

In this work, we draw heavy influence from \cite{mitra2023guaranteed}, where it was shown that for general elliptic problems, the total error, measured in an iteration-dependent energy dual-norm, can be orthogonally decomposed to a fully-computable linearization component and a discretization component arising solely from the linearization step. This produced fully-computable locally-efficient estimates, robust with respect to nonlinearities, paving the way for later works such as \cite{harnist2023robust,stokke2023adaptive,javed2025robust}. The estimates are even on a fixed norm for the L-scheme. Using the same principle here for the L-scheme, yields the following decomposition in \Cref{subsec:decomposition}
\[
\underbrace{\text{Total Error}}_{\substack{\text{measured by a fixed }\\\text{dual energy-norm of residual}}}\approx \underbrace{\left(\text{Linearization error}\right)}_{\substack{
\text{measured by the energy norm}\\
\text{of consecutive iterates}}} + \underbrace{\left(\text{Discretization error of the heat equation}\right)}_{\substack{
\text{measured by the dual energy norm }\\
\text{of the residual of the heat equation}}}.
\]
The discretization error due to the heat equation can potentially be estimated locally efficiently in space-time (see \cite{ErnSmearsVohralik2017} for the time-stepping case). Thus, \emph{so can the total error, and, in fact, robustly with respect to nonlinearities/degeneracies}. We propose such a reliable estimator that is fully-computable in parallel, and easily implementable in \Cref{sec:a-posteriori}. Using this estimator, a fully adaptive discretization/linearization scheme is proposed, which runs L-scheme iterations until linearization error is below a fraction of the discretization error, and then performs a refinement step based on the a posteriori estimator. Numerical experiments demonstrate the efficacy of this adaptive strategy.

The paper is organized as follows. 
\Cref{sec:setting} sets up the functional framework, states the main assumptions, and derives the weak and residual formulations. \Cref{Linearization} introduces the space-time $L$-scheme, establishes its rigorous convergence, and illustrates the convergence behaviour numerically. \Cref{subsec:decomposition} establishes the quasi-orthogonal decomposition of the total error into linearization and discretization components. \Cref{sec:apost} builds on this decomposition to derive fully-computable guaranteed a posteriori error bounds and proposes the adaptive refinement algorithm driven by local error indicators. \Cref{sec:Numerical_results} contains numerical experiments to validate the performance of the adaptive refinement strategy. Our findings are summarized in \Cref{sec:conclusion}.



%To derive such adaptive mesh refinement strategies, appropriate a posteriori error estimators are required. For elliptic boundary value problems, a posteriori error estimation is well established, see, e.g., \cite{verfurth1996review, VERFURTH199467}.
%Recent progress for nonlinear and doubly degenerate \emph{elliptic} problems has been achieved in \cite{ahmed2024efficient}, where guaranteed estimators based on equilibrated flux reconstructions were developed, combining discretization, regularization, and linearization error control with adaptive L-scheme/Newton switching. More generally, a framework for guaranteed and locally efficient a posteriori estimates for nonlinear elliptic problems was established in \cite{mitra2023guaranteed}, yielding an orthogonal decomposition of the total error into linearization and discretization components with adaptive stopping criteria ensuring natural balance between the two. However, both \cite{ahmed2024efficient} and \cite{mitra2023guaranteed} are restricted to stationary elliptic problems and do not address the additional challenges arising in parabolic problems. The construction of estimators for time-dependent problems is considerably more challenging, in particular when aiming at fully space-time adaptive settings for nonlinear degenerate problems.

% Early contributions laid the foundations for adaptive space-time strategies for parabolic problems \cite{Meidner2007}, 
% including goal-oriented approaches \cite{ besier2012goal,Meidner2012}. 



%Early contributions laid the foundations for adaptive space-time strategies for parabolic problems \cite{Meidner2007}, including goal-oriented approaches based on the dual-weighted residual methodology \cite{Meidner2012, besier2012goal}. Guaranteed a posteriori error estimates that are locally efficient in space and time and robust with respect to the polynomial degree have been established for high-order space-time discretizations of linear parabolic problems using equilibrated flux reconstructions \cite{ErnSmearsVohralik2017, ErnSmearsVohralik2019}, however, for the time-stepping case without local time-refinement. More broadly, a posteriori estimators for linear and semilinear parabolic problems have also been developed in sequential settings \cite{eriksson1991adaptive, Eriksson1995, schmich2008adaptivity, FUHRER202127, RattiVerani2019, GeorgoulisLakkisWihler2021}. Earlier, adaptive space-time DG simulations of variably-saturated porous media flows, using hp-adaptive strategies balancing spatial, temporal, and nonlinear algebraic errors, were carried out in \cite{DOLEJSI2019276}. Building on these, a posteriori estimators for nonlinear convection-diffusion-reaction problems were derived in \cite{DolejsiRoskovec2021} for continuous and discontinuous Galerkin discretizations in mesh-dependent norms, and further developed for the Richards equation in \cite{dolejvsi2024error}. Recently, residual-based a posteriori error bounds have been derived for space-time DG discretizations of the Richards equation \cite{congreve2026residual}, however, not accounting for the full degeneracy of $\Phi$.


%For the degenerate Richards equation specifically, a posteriori estimates separating space, time, and linearization errors have been developed within sequential time-stepping frameworks \cite{BernardiElAlaouiMghazli2014, BARON2017104, CLEMENT2021103897}. These works combine residual-based estimators with adaptive refinement and, in some cases, flux reconstruction techniques. Fully computable locally space-time efficient estimates for the fully degenerate case were derived in \cite{mitra2024posteriori}, though within a sequential time-stepping framework rather than a global space-time one.

% However, none of these works provides a guaranteed, fully computable error estimator within a fully global monolithic 
% space-time framework with explicit separation of linearization and discretization errors simultaneously. The works \cite{DOLEJSI2019276, dolejvsi2024error, congreve2026residual}, while adopting a space-time formulation for the Richards equation, which can exhibit both \emph{parabolic--hyperbolic} and \emph{parabolic--elliptic} degeneracies, do not treat the entire space-time 
% domain as a unified whole, solving the problem sequentially rather than globally. Moreover, \cite{DOLEJSI2019276} eliminates degeneracy by regularization, while \cite{dolejvsi2024error, congreve2026residual} do not establish estimates valid in the fully degenerate case. The estimators of \cite{DolejsiRoskovec2021, dolejvsi2024error, 
% congreve2026residual} are measured in \emph{mesh-dependent} norms, in contrast to the \emph{fixed} dual norm of \cite{ErnSmearsVohralik2017, ErnSmearsVohralik2019}. While 
% \cite{mitra2024posteriori} achieves fully computable estimates in a fixed $H^1(H^{-1}) \cap L^2(H^1)$ norm with 
% explicit separation of space, time, and linearization errors for the fully degenerate Richards equation, this is restricted to a sequential time-stepping framework.





% %The works \cite{DOLEJSI2019276, dolejvsi2024error, congreve2026residual}, while adopting a space-time formulation for the Richards equation, a nonlinear parabolic problem that can exhibit both \emph{parabolic--hyperbolic} and \emph{parabolic--elliptic} degeneracies are implemented slab-wise, solving the problem sequentially over time intervals using locally adapted spatial meshes rather than treating the entire space-time domain as a unified whole. Moreover, none of the above works fully addresses the doubly degenerate regime: \cite{DOLEJSI2019276} eliminates degeneracy by regularizing the constitutive relations, while \cite{dolejvsi2024error, congreve2026residual} consider the variably saturated case without establishing estimates valid in the fully degenerate case. Only \cite{mitra2024posteriori} addresses the fully degenerate Richards equation, but within a sequential time-stepping framework. Furthermore, the estimators of \cite{DolejsiRoskovec2021, dolejvsi2024error, congreve2026residual} are measured in \emph{mesh-dependent} residual-based norms that depend on the discretization parameters and change with mesh refinement, in contrast to the \emph{fixed} dual norm, independent of the discretization, employed in \cite{ErnSmearsVohralik2017, ErnSmearsVohralik2019} for linear parabolic problems. While \cite{mitra2024posteriori} achieves fully computable estimates in a fixed norm with explicit separation of linearization and discretization errors for the fully degenerate Richards' equation, this is restricted to a sequential time-stepping framework rather than a fully global space-time variational framework.

% The present work addresses all these gaps simultaneously. Building on the orthogonal decomposition framework of \cite{mitra2023guaranteed}, developed for nonlinear and possibly degenerate elliptic problems, we extend these ideas to the fully global space-time setting for the doubly degenerate problem \eqref{eq:reformulation}, treating time 
% and space on an equal footing. We develop guaranteed equilibrated flux-based a posteriori error estimators in a \emph{fixed} dual norm, in contrast to existing approaches 
% \cite{DolejsiRoskovec2021, dolejvsi2024error, congreve2026residual} that rely on mesh-dependent norms. The total error is bounded above by $\eta_{\mathrm{lin}} +\eta_{\mathrm{disc}}$, and the discretization error is bounded above by the total error plus $\eta_{\mathrm{lin}}$, yielding a quasi-orthogonal decomposition that drives both the adaptive stopping criterion for the linearization iterations and the local space-time mesh refinements. To the best of our knowledge, this constitutes the first guaranteed a posteriori error analysis for nonlinear doubly degenerate parabolic problems in a global space-time setting.
% %In the present work, we adopt a fully global space-time variational approach, treating space and time as independent variables and seeking the solution over the entire space-time domain rather than relying on a sequential slab-wise formulation. This formulation provides a unified mathematical framework that captures the strong coupling between spatial and temporal dynamics inherent in nonlinear and degenerate parabolic problems and allows for a genuinely space-time treatment of evolving solution features. However, this advantage comes at a high computational cost.
% %Given the increasing demand for accurate simulations in highly complex nonlinear systems, space–time discretizations offer a promising alternative to traditional time-stepping approaches.
% %In particular, storing and solving the entire evolution of a solution over time without adaptivity may require terabytes of memory, making naive space–time approaches infeasible in practice. This so-called curse of dimensionality becomes even more dramatic for higher-dimensional problems.

% %\noindent
% %At the same time, the increasing demand for accurate simulations of highly complex nonlinear systems continues to motivate the development of space–time discretization techniques, highlighting the need for adaptive strategies that can make such approaches computationally viable. Space-time adaptive refinement strategies aim to concentrate computational effort only where and when it is needed, for instance, in regions where the solution exhibits complex dynamics such as steep fronts, sharp transitions, or localized nonlinear effects. In this way, adaptivity plays a key role in making fully space-time discretizations feasible for nonlinear degenerate parabolic problems.

% %\AJ{The ultimate aim of adaptive methods is to obtain an approximate solution with an error below a prescribed tolerance at the expense of, up to an absolute multiple, the minimal amount of computational time and storage. For linear parabolic problems, space-time adaptive wavelet methods have been shown to achieve this goal at optimal computational complexity, even overcoming the curse of dimensionality through the use of tensorized multilevel bases \cite{schwab2009space}. The extension of such results to strongly nonlinear and degenerate parabolic problems, however, remains highly challenging. Achieving this goal relies crucially on a posteriori error estimation, which provides computable and locally meaningful indicators that guide mesh refinement in both space and time. Adaptive space-time finite element methods driven by such estimators have been successfully applied to parabolic optimization and control problems \cite{Meidner2012,gong2012space,LangerSchafelner+2022+247+266,GunzburgerKunoth2011}, to problems with distributional sources \cite{LangerSchafelner+2022+247+266}, and to certain classes of nonlinear and quasi-linear parabolic equations, including parallel in space-time frameworks \cite{Dyja2018}. These works clearly demonstrate that adaptivity is essential for making space-time methods computationally viable.}


% %\noindent
% %\AJ{To derive such adaptive mesh refinement strategies, appropriate a posteriori error estimators are required. For elliptic boundary value problems, a posteriori error estimation is well established, see, e.g., \cite{verfurth1996review,VERFURTH199467}. Recent progress for nonlinear and doubly degenerate elliptic problems has been achieved using equilibrated-flux-based a posteriori error estimation. In \cite{ahmed2024efficient}, guaranteed and locally efficient estimators were developed for a class of doubly degenerate elliptic equations, including fast and slow diffusion models. Their framework combines equilibrated flux reconstructions with an adaptive nonlinear solver that balances discretization, regularization, and linearization errors, and dynamically switches between L-scheme and Newton iterations. However, these developments are restricted to stationary elliptic problems and do not address the additional challenges arising in fully space--time formulations of parabolic problems.}

% %In contrast, the construction of estimators for time-dependent problems is considerably more challenging, in particular when aiming at fully space-time adaptive settings. Although a posteriori error estimators for parabolic problems have been developed, see, e.g., \cite{eriksson1991adaptive, Eriksson1995,schmich2008adaptivity, FUHRER202127}\textcolor{red}{\cite{RattiVerani2019,GeorgoulisLakkisWihler2021}}, many existing approaches apply adaptivity either in space or in time, or rely on sequential time-stepping strategies \cite{BernardiElAlaouiMghazli2014, BARON2017104,CLEMENT2021103897, dolejvsi2024error,congreve2026residual}. These works combine residual-based estimators with adaptive refinement and, in some cases, flux reconstruction techniques. However, they do not provide a guaranteed, fully computable error estimator within a fully global monolithic space–time framework, nor do they explicitly separate linearization and discretization errors within the adaptive strategy. Notably, fully computable and locally space-time efficient a posteriori error estimates for the fully degenerate Richard's equation, with explicit separation of linearization and discretization errors, have been developed in \cite{mitra2024posteriori}, though within a sequential time-stepping framework rather than a fully global space-time formulation.

% %\noindent
% %To make fully space-time adaptive methods feasible, it is necessary to construct error estimators that decompose the residual into spatial, temporal, and nonlinear contributions.
% %Such estimators rely on local indicators that are reliable, efficient, and fully computable at the discrete level. Building on these principles, we introduce an estimator framework that enables this decomposition and drives the proposed adaptive refinement strategy.

% %\noindent
% %In contrast to existing residual-based and slab-wise approaches, the present work develops guaranteed equilibrated flux-based a posteriori error estimators within a fully global space–time variational framework. The proposed estimators provide a unified and fully computable control of the discretization error in both space and time. In addition, the framework clearly separates linearization and discretization errors, allowing mesh refinement to be triggered only when the discretization error dominates. This integration of nonlinear iteration control with fully space–time adaptive refinement constitutes a central novelty of the present work.
% The resulting adaptive strategy is described in the following subsection.
% \subsection{Adaptive space–time refinement strategy}\label{sec:adap_SP_refinement}
% The adaptive strategy is based on monitoring the local discretization error during the nonlinear iteration and refining the mesh only in regions where this error is most significant. To achieve this, we reconstruct the equilibrated flux $\boldsymbol{\sigma}^i_{h\tau}$ (see \Cref{def:equilibrated} and \Cref{mass_balance_property}) and evaluate the fully computable a posteriori estimators 
% $\eta^i_K$ on each element $K\in\mathcal{T}_{h\tau}$ (see \Cref{thm:apost-reliability-i}), together with the 
% global discretization estimator $\eta^i_{\mathrm{disc}}$ defined in \eqref{eq:Hminus1-R}. These estimators precisely quantify where the discretization error dominates the 
% linearization error, thereby guiding the refinement process. Whenever the linearization error $\eta^i_{\mathrm{lin}}$ is sufficiently small in comparison with $\eta^i_{\mathrm{disc}}$, the adaptive module is activated. Elements associated with the largest values of $\eta^i_K$ 
% are then marked for refinement using the D\"{o}rfler bulk criterion, and the current solution is projected onto the resulting locally enriched mesh. In this way, computational effort is concentrated only where it is needed, leading to an efficient adaptive space-time refinement strategy. The complete adaptive procedure is summarised in Algorithm~\ref{alg:Lscheme-adapt} and illustrated in Figure~\ref{fig:flowchart}.
% %This leads to a highly efficient strategy in which computational effort is concentrated only 
% %where necessary. The resulting adaptive procedure, summarised in Algorithm~\ref{alg:Lscheme-adapt} and illustrated in Figure~\ref{fig:flowchart}.
% %===============================================
% \begin{figure}[h!]
% %\vspace*{-1.5cm}     % adjust this to pull the figure upward if needed
% \centering
% \scalebox{0.78}{%
% \begin{tikzpicture}[
%   >=Stealth,
%   node distance=10mm and 15mm,
%   every node/.style={font=\small, inner sep=3pt},
%   block/.style   ={rectangle, rounded corners, draw, align=center,
%                    minimum width=4cm, minimum height=8mm, fill=green!10},
%   decision/.style={
%     diamond, aspect=2, draw, align=center,
%     inner sep=0.5pt, % tighter
%     fill=yellow!20
%   },
%   boxblue/.style ={rectangle, rounded corners, draw, align=center,
%                    minimum width=4cm, minimum height=8mm, fill=blue!10},
%   stopbox/.style ={rectangle, rounded corners, draw, align=center,
%                    minimum width=4.2cm, minimum height=12mm, fill=red!25, font=\normalsize\bfseries},
%   line/.style    ={draw, -{Stealth}, thick}
% ]

% %--------------------------------------------------
% % MAIN COLUMN
% %--------------------------------------------------
% \node[block] (init)
%   {\textbf{Start}\\Given initial mesh $\mathcal T_{h\tau}^0$, $\ell_{\max}$\\
%   initial guess $s^0$, parameters:\\
%   $\ i_{\max},\ \varepsilon_{\mathrm{sol}},\ \varepsilon_{\mathrm{enter}},\ 
%   \delta_{\mathrm{lin}}$};

% \node[block, below=10mm of init] (solve)
%   {Solve linearized step\\ compute $s^{i+1}$};

% \node[block, below=8mm of solve] (error)
%   {Compute error:\\ $\mathrm{Error}=\|\nabla (s^{i+1}-s^i)\|_{L^2}$};

% \node[decision, below=10mm of error] (enter)
%   {Error sufficiently \\small to refine?\\
% $\mathrm{Error}\le \varepsilon_{\mathrm{enter}}$?};

% \node[decision, below=12mm of enter] (stoperr)
%   {Has the iteration\\fully converged?\\
%    $\mathrm{Error}\le \varepsilon_{\mathrm{sol}}$?};

% \node[stopbox, below=22mm of stoperr] (stop)
%   {\textbf{STOP}};

% %--------------------------------------------------
% % RIGHT COLUMN (iteration continuation)
% %--------------------------------------------------
% \node[block, right=40mm of stoperr.south|-stoperr.south] (update)
%   {Update $s^i\!=\!s^{i+1}$\\ $i=i+1$};

% \node[decision, below=9mm of update] (imax)
%   {Reached max\\iterations?\\{$i\ge i_{\max}$?}};

% %--------------------------------------------------
% % LEFT COLUMN (adaptivity)
% %--------------------------------------------------
% \node[boxblue, left=25mm of enter] (recon)
%    {Flux reconstruction:\\ $\boldsymbol{\sigma}^i_{h\tau}$};


% \node[boxblue, below=8mm of recon] (est)
%  {Compute estimators:\\
% $\eta^i_K$, $\eta^i_{\mathrm{disc}}$};

% \node[decision, below=8mm of est] (gate)
%   {$\mathrm{Error}\le \delta_{\mathrm{lin}}\eta^i_{\mathrm{disc}}$?};

% % combined mark + refine box
% \node[boxblue, below=8mm of gate] (markref)
%   {Identify regions with largest error\\[1mm]
%    and refine the marked elements};

% % *** FIXED: updmesh is below markref (not refine) ***
% \node[boxblue, below=8mm of markref] (updmesh)
%   {Update mesh $\mathcal T_{h\tau}$ and\\project solution,\\$\ell\gets\ell+1$};

% % refinement-level check
% \node[decision, below=8mm of updmesh] (levelcheck)
%   {Reached max\\refinement level?\\$\ell\ge\ell_{\max}$?};

% %--------------------------------------------------
% % ARROWS
% %--------------------------------------------------

% % main downward flow
% \draw[line] (init)   -- (solve);
% \draw[line] (solve)  -- (error);
% \draw[line] (error)  -- (enter);

% % "enter" decision
% \draw[line] (enter.west)  -- node[above,pos=0.5]{\textbf{YES}} (recon.east);
% \draw[line] (enter.south) -- node[right,pos=0.4]{\textbf{NO}}  (stoperr.north);

% % convergence decision
% \draw[line] (stoperr.south) -- node[right,pos=0.45]{\textbf{YES}} (stop.north);

% % stoperr -> update (No branch, routed around)
% \draw[line] (stoperr.east) -- ++(10mm,0)
%            |- node[pos=0.8,above]{\textbf{NO}} (update.north);

% % update -> imax
% \draw[line] (update.south) -- (imax.north);

% % imax -> STOP (Yes)
% \draw[line] (imax.south) |- ++(0,-4mm) -|
%            node[pos=0.25,above]{\textbf{YES}} (stop.east);

% % imax -> solve (No)
% \draw[line] (imax.west) -- ++(-10mm,0) |- node[pos=0.3,right]{\textbf{NO}} (solve.east);

% % left column (adaptivity branch)
% \draw[line] (recon.south) -- (est.north);
% \draw[line] (est.south)   -- (gate.north);

% % gate YES -> combined mark+refine
% \draw[line] (gate.south)  -- node[right,pos=0.4]{\textbf{Yes}} (markref.north);

% % combined mark+refine -> update mesh
% \draw[line] (markref.south)  -- (updmesh.north);

% % update mesh -> level check
% \draw[line] (updmesh.south) -- (levelcheck.north);

% % *** NEW: levelcheck YES -> STOP with short vertical + horizontal + up ***
% \draw[line]
%   (levelcheck.south) -- ++(12mm,0) node[above,pos=4.5]{\textbf{YES}} % short down
%   -- ++(58mm,0)                                                      % right
%   -- (stop.south);                                                   % up into STOP

%  % levelcheck NO -> go to update (so solution is updated, then loop)
% \draw[draw, thick, -{Stealth[length=2mm]}]
%   (levelcheck.east) -- ++(8mm,0)
%   |- node[pos=0.03,left]{\textbf{NO}} (update.south);

% % gate -> update (No branch, routed under the center diamond)
% \draw[line] 
%   (gate.east) -- ++(12mm,0) 
%   |- node[pos=0.010,above]{\textbf{NO}} (update.south);
% \end{tikzpicture}}
% \caption{[\Cref{sec:adap_SP_refinement}] Flowchart of the adaptive space-time iteration and mesh refinement algorithm.  
% Two error thresholds are used to guide the algorithm: 
% (1) $\varepsilon_{\text{enter}}$ (linearization threshold) and 
% (2) $\varepsilon_{\text{sol}}$ (solution accuracy threshold). 
% Their roles are interpreted as follows: 
% if $\mathrm{Error} > \varepsilon_{\text{enter}}$, the error is dominated by the linearization process and further iterations can improve the result; 
% if $\varepsilon_{\text{sol}} < \mathrm{Error} \leq \varepsilon_{\text{enter}}$, 
% the linearization error is no longer dominant and the discretization error becomes significant; 
% therefore, adaptive mesh refinement is performed; if $\mathrm{Error} \leq \varepsilon_{\text{sol}}$, the solution is considered converged and sufficiently accurate.
% }
% \label{fig:flowchart}
% \end{figure}
% %===============================================

\section{Mathematical formulation}\label{sec:setting}
In this section, we introduce the mathematical framework used throughout the paper. We start from the nonlinear degenerate parabolic model presented in \Cref{introduction}, followed by a reparametrization strategy. We then present the functional framework, weak formulation and the linear iterative scheme.
%\subsection{Model problem}\label{subsec:model-problem}
%With $\Om$ being a bounded domain in $\R^{d}$ having a Lipschitz boundary $\p\Om$ and for some $T > 0$, set $Q:=(0,T]\times\Omega$. We consider the problem
%\begin{equation}\label{eq:parabolic_equation}
%\left\{ \begin{array}{rcll}
%\partial_t u &= &\Delta w + f, \qquad & \text{ in } Q, \\
%w &=& 0, & \text{ on } (0, T] \times \partial \Omega, \\
%u(0) &=& u_0, &\text{ in } \Omega, \\
%\end{array}
%\right.
%\end{equation}
%where $w\in \Phi(u)$, with double degenerate $\Phi$ as shown below in \Cref{fig:RichVsBio_phi}. 

%\begin{figure}[h!]
%    \centering
%    \begin{subfigure}[b]{0.35\textwidth}
%        \centering
%    \includegraphics[width=\textwidth, height=4.1cm]{Figures/RichardsDegeneracy.eps}
%        \caption*{$\Phi$ becomes multivalued at $u = 1$}
%    \end{subfigure}
%    \begin{subfigure}[b]{0.35\textwidth}
 %       \centering
 %       \includegraphics[width=\textwidth, height=4.1cm]{Figures/BiofilmDegeneracy.eps}
%        \caption*{$\Phi$ becomes infinite at $u = 1$}
%    \end{subfigure}
%    \caption{Examples of doubly degenerate $\Phi$: Richards vs. Biofilm type degeneracy.}
%\label{fig:RichVsBio_phi}
%\end{figure}
\subsection{Reparametrization}\label{subsec:reparametrization}
In view of the modeled processes mentioned before, one can see $u$ as a concentration or saturation. Then it is natural to think of an upper bound $\omega$ of $u$. Here, we either take $\omega = 1$ if $u$ is, for example, a concentration or saturation, or $\omega = \infty$ for unbounded solutions. We assume
%Let $\vs \in \{1,\infty\}$ be the physical upper-bound on the concentration-type variable $u$. For biofilm or Richards equation $\vs=1$ with $u$ being interpreted as a biofilm concentration or water saturation. In the porous medium equation case, $\vs=\infty$. We assume
\begin{enumerate}[label=(A.1\alph*)]
\item \label{ass:A1a}
The function $\Phi:(-\infty,\vs) \to [0,\infty)$ is continuous, differentiable, and strictly increasing in $(0,\vs)$. 
Moreover, $\Phi(s)=0$ for $s\leq 0$, and either $\lim_{s\searrow 0}\Phi'(0)>0$, or $\Phi$ is strictly convex in a right neighbourhood of 0. Further, the limit $\Phi_\Maxi:=\lim_{u\nearrow \vs}\Phi$ is either infinite, or, if $\Phi_\Maxi<\infty$, then we extend $\Phi$ to the set $[\Phi_\Maxi,\infty)$ at $u=\vs$.  
%Finally, assume that for some $p>0$,
% \begin{align}\label{eq:growth}
%     (u_1-u_2)(\Phi(u_1)- \Phi(u_2))\geq |u_1-u_2|^{p+1} \text{ for all } 0\leq u_1,\,u_2\leq \vs.
% \end{align}
\end{enumerate}

Assumption \ref{ass:A1a} is very general, and covers all the nonlinearities described in \Cref{fig:RichVsBio_phi}.

% \begin{remark}[Assumption \ref{ass:A1a}]
%     Assumption \ref{ass:A1a} is very general, and covers all the nonlinearities described in \Cref{fig:RichVsBio_phi}. %In particular, for porous media type growth, i.e. $\Phi(u)~u^p$, \eqref{eq:growth} is satisfied since 
%     % \[
%     % |u_1^p -u_2^p|\geq |u_1-u_2|^p.
%     % \]
% \end{remark}
\begin{enumerate}[label=(A.2\alph*)]
\item \label{ass:A1b}
There exist functions $b, B$  that are Lipschitz continuous, non-decreasing, and satisfy
\begin{equation}\label{eq:conditionbB1}
\left\{
\begin{aligned}
&B(s) = (\Phi \circ b)(s), \quad s \ge 0, \quad b(0)=0,\\
&|B(s) - B(r)| \le |s - r|, \quad |b(s) - b(r)| \le |s - r|, \forall\, s,r \in \mathbb{R},\\
&(b+B)(s) - (b+B)(r) \ge s - r, \quad \forall\, s \ge r .
\end{aligned}
\right.
\end{equation}
\end{enumerate}

\noindent
Note that by introducing the reparametrizations $b$ and $B$, the double degeneracy in \eqref{eq:parabolic_equation} due to $\Phi$ can be transformed into two degeneracies in \eqref{eq:reformulation}, involving two Lipschitz continuous functions satisfying \eqref{eq:conditionbB1}. For any convex $\Phi$, these functions can be found explicitly, as follows from  \cite{javed2025robust}, Lemma~2.5. This is summarized below.




%The $b$ and $B$ reparametrization of $\Phi$ was introduced first in \cite{brenner2017improving}. The idea is to ``divide and conquer'' the double degeneracy in \eqref{eq:parabolic_equation} due to $\Phi$ to two single degeneracies in \eqref{eq:reformulation} caused by $b$ and $B$ that are themselves Lipschitz. The functions $b$-$B$ can be found for any convex $\Phi$ by the following explicit expression:


\begin{lemma}[Construction of the $b$--$B$ decomposition]\label{lemma:b_B_construction}Let 
$\Phi$ have locally Lipschitz continuous derivatives in $(0,\vs)$. Assume that there exists $u^*\in (0,\vs)$ such that (after a possible rescaling of $\Phi$),\label{ass:Bphi}
\begin{equation*}
    \Phi'(u)\begin{cases}
        \leq 1 &\text{ for } u\in (0,u^*),\\
        =1&\text{ for } u=u^*,\\
        \geq 1 &\text{ for } u\in (u^*,\vs).
    \end{cases}
\end{equation*}    
Then, the functions $b, B : \mathbb{R} \to \mathbb{R}$,
\begin{subequations}\label{eq:bBexpression}
    \begin{align}
            &b(s):=\int_0^s \min\left\{1, \frac{1}{\Phi'(b(\rho))}\right\}d\rho=\begin{cases}
             s &\text{ if } s\leq u^*,\\
             \Phi^{-1}(\Phi(u^*) + s-u^*) &\text{ if } s\geq u^*.
         \end{cases}\\
         &B(s):=\int_0^s \min \left\{1, \Phi'(b(\rho))\right\}d\rho=\begin{cases}
             \Phi(s) &\text{ if } s\leq u^*,\\
             \Phi(u^*) + s-u^*&\text{ if } s\geq u^*,
         \end{cases}
    \end{align}   
\end{subequations}
satisfy Assumption \ref{ass:A1b}.

\end{lemma}


\noindent
By an abuse of notation, we take $\min\left\{1, \frac{1}{\Phi'(z)}\right\} = 1$ if $\Phi'(z)=0$. 
This means that $b(s)=s$ and $B(s)=0$ for any $s \le 0$. %The argument $u^*>0$ is found by solving $\Phi'(u^*) = 1$, which, under the assumptions of \Cref{lemma:b_B_construction}, has a solution. 
%In case of multiple solutions, we take the lowest one. 
%Observe that, by \ref{ass:A1a}, $u^*>0$. 
For subsequent analysis, for any Lipschitz continuous $\psi$, we use the notation:

%The proof of the statement above follows Lemma~2.5 in \cite{javed2025robust}. The parameter $u^*>0$ is defined implicitly by $\Phi^\prime (u^*) = 1$. For subsequent analysis, we define for $\psi \in\mathrm{Lip}(\R)$ the notation 

\begin{equation}\label{eq:def_errors_b}
    \psi[s,r]
    :=
    \begin{cases}
        \dfrac{\psi(s)-\psi(r)}{s-r}, &\text{if } s\neq r, \\[1em]
        \dfrac{1}{2}\left[\lim\limits_{z\nearrow s} \psi'(z) + \lim\limits_{z\searrow s} \psi'(z)\right], &\text{if }  s=r.
    \end{cases}
\end{equation}
%Then the following extensions of $b$--$B$ to negatives are imposed since the finite element solution can become negative.
Using \ref{ass:A1b}, one easily gets for any $s, r \in \mathbb{R}$,
\begin{align}\label{eq:condition}
0 \leq b[s,r], \quad B[s,r]\leq 1 \quad {\text{and}} \quad 
 b[s,r]+B[s,r] \geq 1.
\end{align} 
The graphs of $b$ and $B$ given by \eqref{eq:bBexpression} for various $\Phi$ function are presented in \Cref{bB-graphs}. 

\noindent
Finally, for the source function $f$ and initial condition $u_0$ we assume the following
\begin{enumerate}[label=(B.\arabic*)]
\item \label{ass:A2} $u_0\in L^2(\Om)$ is in $[0,\omega)$ a.e. The source function $f\in L^2(Q)$ admits a non-negative solution to Problem~\ref{eq:reformulated_weak_form}.
\end{enumerate}
%\begin{definition}[Extension of the $b$--$B$ decomposition]\label{assump_bB}
%The functions $b,\, B \in\mathrm{Lip}([0,\infty)])$ that satisfy \ref{ass:A1b} are extended to $\R$ by
 %   \begin{align}\label{eq:conditionbB1_neg}
%b(s)=0, \quad B(s)=s \quad \text{ for } \quad s<0.
%    \end{align}
%Then for all  $s,\,r\in \R$ we have
%    \begin{align}\label{eq:condition}
%    \quad 0 \leq b[s,r], \, B[s,r] \leq 1, \quad \text{and} \quad b[s,r] + B[s,r] \geq 1.
%\end{align}
%\end{definition}

% let $\Phi:(0,\omega)\to \R^+$ be an increasing function in 
% Equation \eqref{eq:parabolic_equation} is double degenerate since $\Phi'(u)$ vanishes as $u\to 0$ and blows-up as $u\t$

%\Cref{bB-graphs} shows $b$ and $B$ functions computed from $\Phi$ functions using \Cref{lemma:b_B_construction} also extended to the negatives.


% the function may become infinite or multivalued near saturation limits. 

%To resolve this issue, we follow the ideas in  \cite{brenner2017improving} and reformulate the system in terms of a new variable $s$, such that 
%\begin{align*}
%    w = B(s) \quad \text{and} \quad  u= b(s)
%\end{align*}
%where $b,\,B$ are strictly monotone but Lipschitz continuous functions. 
%The underlying idea is to split the challenges associated with $\Phi$, which can exhibit double degeneracy, to two simpler functions $b$ and $B$. Then, 
%\eqref{eq:parabolic_equation} is reformulated as the system
%\begin{equation}\label{eq:reformulation}
%\left\{ \begin{array}{rcll}
%\partial_t b(s) &= &\Delta B(s) + f, \qquad & \text{ in } Q, \\
%B(s) &=& 0, & \text{ on } (0, T] \times \partial \Omega, \\
%b(s(0)) &=& u_0, &\text{ in } \Omega. \\
%\end{array}
%\right.
%\end{equation}
% Below, we outline the more general assumptions on $b$ and $B$. Introducing for any piecewise differentiable function $\psi \in\mathrm{Lip}(\R)$ the notation 
% \begin{equation}\label{eq:def_errors_b}
%     \psi[s,r]
%     :=
%     \begin{cases}
%         \dfrac{\psi(s)-\psi(r)}{s-r}, &\text{if } s\neq r, \\[1em]
%         \dfrac{1}{2}\left[\lim\limits_{\underline{s}\nearrow s} \psi'(\underline{s}) + \lim\limits_{\underline{s}\searrow s} \psi'(\underline{s})\right], &\text{if }  s=r,
%     \end{cases}
% \end{equation}
%  the full properties that $b$, $B$ need to satisfy are as follows:
% \begin{assumption}[Properties of the $b$--$B$ decomposition]\label{assump_bB}
% The functions $b,\, B \in\mathrm{Lip}(\R)$ are piecewise differentiable, and satisfy
%     \begin{align}\label{eq:conditionbB1}
%     \begin{cases} B(s) = (\Phi \circ b)(s) &\text{ for } s> 0, \\
%     b(s)=0 \quad \text{and} \quad B(s)=s &\text{ for } s\leq 0.
%     \end{cases}
% \end{align}
% Moreover, for all $s,\,r\in \R$ we have
%     \begin{align}\label{eq:condition}
%     \quad 0 \leq b[s,r], \, B[s,r] \leq 1, \quad \text{and} \quad b[s,r] + B[s,r] \geq 1.



% \end{align}
% \end{assumption}




% \noindent
% The construction of the $b$-$B$ function follows Lemma~2.4 in \cite{javed2025robust}; see also \cite{brenner2017improving} for a similar formulation. Since the spatially discrete solutions may fail to be positive a.e., we adopt the extension to $s<0$ introduced in \cite{Geurts2025Biofilm} which guarantees convergence. By construction, the function \(B\) satisfies \(B(0)=0\), so \(w=0\) implies that \(s=0\) on \(\partial\Omega\).

% \begin{remark}[Construction of the $b$--$B$ decomposition]\label{lemma:b_B_construction}
% Assume that $\Phi\in C^1(\mathbb{R})$ is strictly increasing and satisfies the
% structural bounds \eqref{eq:condition}.
% Following Lemma~2.4 in \cite{javed2025robust} and the discussion in
% \cite{brenner2017improving}, one defines
% \begin{subequations}\label{eq:bBexpression}
%     \begin{align}
%             &b(s):=\int_0^s \min\left\{1, \frac{1}{\Phi'(b(\rho))}\right\}d\rho=\begin{cases}
%              s &\text{ if } s\leq u^*,\\
%              \Phi^{-1}(\Phi(u^*) + s-u^*) &\text{ if } s\geq u^*.
%          \end{cases}\\
%          &B(s):=\int_0^s \min \left\{1, \Phi'(b(\rho))\right\}d\rho=\begin{cases}
%              \Phi(s) &\text{ if } s\leq u^*,\\
%              \Phi(u^*) + s-u^*&\text{ if } s\geq u^*.
%          \end{cases}
%     \end{align}   
% \end{subequations}
% where $u^*>0$ is defined implicitly by $\Phi^\prime (u^*) = 1$.
% %Under additional regularity assumptions on $\Phi$, explicit formulas for the
% %splitting functions $b$ and $B$ can be derived see Lemma~2.4 in \cite{javed2025robust} and the discussion in
% %\cite{brenner2017improving}.
% This construction guarantees that $b$ and $B$ are monotone and Lipschitz continuous. The extension of $b$ and $B$ to $s<0$ used in this work is described in
% Assumption~\ref{assump_bB} %Since discrete approximations may fail to remain nonnegative,
% %the functions $b$ and $B$ are extended to $s< 0$ as specified in
% %Assumption \ref{assump_bB}, following \cite{Geurts2025Biofilm}.
% This extension preserves the bijectivity of $B$ and ensures that the
% divided--difference estimates used in the convergence analysis remain valid
% on the whole real line.
% \end{remark}
\subsection{Notations and functional spaces}\label{subsec:functional-spaces}
We briefly introduce the functional setting used throughout the paper. 
%Consider the bounded Lipschitz domain $\Omega\subset \R^d$, final-time $T>0$, space--time cylinder $Q := (0,T] \times \Omega$, and boundary $\Gamma:= (0,T] \times \p\Omega$.
Let $L^2(\Omega)$ be the space of measurable, square-integrable functions defined on $\Omega$. The space $H^1(\Omega)$ includes functions in $L^2(\Omega)$ having weak derivatives in $L^2(\Omega)$, while, $H^1_0(\Omega)$ contains functions in $H^1(\Omega)$ having vanishing trace on $\partial \Omega$. 
The dual space of $H^1_0(\Omega)$ is denoted by $H^{-1}(\Omega)$. The inner product and norm in $L^2(\Omega)$ are denoted by $(\cdot, \cdot)$ and $\|\cdot\|$, respectively, while the duality pairing between $H^{-1}(\Omega)$ and $H_0^1(\Omega)$ is denoted by $\langle \cdot, \cdot \rangle$. We also use the notation $a\lesssim b$ to indicate that there exists a generic constant $C>0$, independent of the discretization parameters, such that $a\leq Cb$.

%\vspace{0.3em}
%\noindent\textbf{Sobolev spaces.}
%We use the standard Sobolev-Hilbert spaces 
%\[
%H_0^1(\Omega) := \{ v \in H^1(\Omega) : v = 0 \text{ in trace sense on } \partial\Omega \}, 
%\qquad 
%H^{-1}(\Omega) := \big( H_0^1(\Omega) \big)^*.
%\]
We will use the Poincar\'e--Friedrichs inequality which states that there exists a constant $C_{\,\!_{\Omega}}^{\mathrm{P}} > 0$ (depending only on the shape of $\Omega$) such that, for any $v \in H_0^1(\Omega)$ one has
\begin{align}\label{eq:Poincare_Friedrichs inequality}
\|v\|\le C_{\,\!_{\Omega}}^{\mathrm{P}} \, h_\Om\,\|\nabla v\|.
\end{align}
\noindent
Here, $h_\Om$ is the diameter of $\Om$. By \eqref{eq:Poincare_Friedrichs inequality}, $\|\nabla v\|$ defines an equivalent norm on $H_0^1(\Omega)$, and a norm on $H^{-1}(\Omega)$ can be defined as
\begin{align}
    \|g\|_{H^{-1}(\Omega)} 
:= \sup_{\varphi \in H_0^1(\Omega)\setminus \{0\}} 
\frac{\langle g, \varphi \rangle}{\|\nabla \varphi\|} \leq C_{\,\!_{\Omega}}^{\mathrm{P}} \, h_\Om \|g\|,\label{eq:defHinv_norm}
\end{align}
the last inequality holding only if additionally $g\in L^2(\Om)$.
By the Riesz representation theorem, for each $g \in H^{-1}(\Omega)$ there exists a unique $G_g \in H_0^1(\Omega)$ satisfying
\begin{align}\label{eq:weak_poisson}
(\nabla G_g, \nabla \varphi) = \langle g, \varphi \rangle \text{ for all } \varphi \in H_0^1(\Omega), \text{ with } \|g\|_{H^{-1}(\Omega)} = \|\nabla G_g\|.
\end{align}
%evident after identifying the $\mathrm{arg}\sup$ in \eqref{eq:defHinv_norm} to be $\f=G_g$.

\noindent
Let $Q'\subseteq Q$ be a space-time subdomain of $Q$. For the fluxes we use the space
\[
H(\mathrm{div};Q')
:= \big\{ \mathbf{v} \in [L^2(Q')]^d : \text{ the weak (spatial) divergence } \nabla\!\cdot \mathbf{v} \in L^2(Q') \big\},
\]
equipped with the norm $\|\mathbf{v}\|_{H(\mathrm{div};Q'.)} 
:= \|\mathbf{v}\|_{L^2(Q')} + \|\nabla\!\cdot \mathbf{v}\|_{L^2(Q')}$.
%Specifically, if $t+\nhat\cdot \bm{x}=e$ denotes a hyperplane inside $Q'$ for some fixed $\nhat\in \R^d$, then $\bm{v}\cdot \nhat$ remains conforming across this  hyperplane for $\bm{v}\in H(\mathrm{div};Q')$. However, it does not include any time derivative, and thus, does not require conformity in the time-direction.

%\AJ{After performing a space-time discretization we will use the  local analogue, on a space-time vertex patche $\widetilde\omega_a\subset (0,T)\times\Omega$,}
% \[
% H(\mathrm{div};\tilde\omega_a;\R^d)
% :=\big\{\boldsymbol v\in L^2(\tilde\omega_a;\R^d): \nabla\!\cdot \boldsymbol v\in L^2(\tilde\omega_a)\big\}.
% \]
% \AJ{Note that, in this case, the divergence operator acts only on the spatial variable.}
\noindent
For any Banach space $X$, the Bochner space  $L^2(0,T;X)$ consists of all strongly measurable functions $v:(0,T)\to X$ such that
\[
\|v\|_{L^2(0,T;X)}
:=\left[\int_0^T \|v(t)\|_X^2\,\dd t \right]^{1/2}
<\infty.
\]
We will also use the Bochner-Sobolev space $H^1(0,T;X)
:=\{v\in L^2(0,T;X): \partial_t v\in L^2(0,T;X)\}$. Using these, we define the main spaces that are used in the subsequent analysis:
\begin{subequations}\label{eq:spaces}
 \begin{align}
\mathcal{S} &:= \Big\{ s \in L^2(Q) 
\;\big|\; 
b(s) \in H^1(0,T;H^{-1}(\Omega)), \;
B(s) \in L^2(0,T;H_0^1(\Omega))\Big\}\\
\mathcal{V} &:= L^2(0,T;H_0^1(\Omega)) \cap H^1(0,T;H^{-1}(\Omega)) \\
\mathcal{W} &:= L^2(0,T;H^2(\Omega)) \cap H^1(0,T;L^2(\Omega)).
\end{align}
\end{subequations}
We will seek weak solutions to \eqref{eq:reformulation} in the space $\calS$. This solution will be obtained through iterations, where the iterates will lie in $\mathcal{V}$. Observe that $\mathcal{W} \subset \mathcal{V}$ and $\mathcal{W}\subset \calS$, but $\calV\not\subset \calS$.
The space $\mathcal{V}$ is also compactly embedded in $L^2(Q)$ due to Aubin-Lions Lemma, and continuously embedded in $C([0,T];L^2(\Omega))$, see \cite{simon1986compact}. This embedding ensures that functions in $\mathcal{V}$ are well-defined for every $t \in [0, T]$. Therefore, one can equip  $\mathcal{V}$ with the norm
\begin{align}\label{eq:norm_V}
    \|v\|_{\mathcal V}^2
:=\int_0^T \!\big(\|\nabla v(t)\|^2+\|\partial_t v(t)\|_{H^{-1}(\Omega)}^2\big)\,\dd t
+\|v(T)\|^2.
\end{align}
\noindent

%=====================================================================================================================================================
%=================================================================================================================
\subsection{Weak formulation}
We can now introduce the concept of a weak solution to \eqref{eq:reformulation}:
%With the definitions of the spaces in place, we write down the weak formulation of \eqref{eq:reformulation}. We would assume:

\begin{problem}[{Weak solution of \eqref{eq:reformulation}}]\label{eq:reformulated_weak_form}
A weak solution of \eqref{eq:reformulation} is a function $s \in \cal S$ satisfying $b(s(0))=u_0$ in $L^2(\Omega)$, and for all $\varphi\in L^2(0,T;H_0^1(\Omega))$,
    \begin{align}\label{eq:weak form}
    \int_0^T\left\langle \partial_t b(s), \varphi\right\rangle+ \int_0^T\left(\nabla B(s), \nabla \varphi\right) = \int_0^T\left( f, \varphi\right).
\end{align}
\end{problem} 
Without being exhaustive, for results concerning the existence and uniqueness of solutions for doubly-degenerate equations we refer to \cite{alt1983quasilinear, ANDREIANOV20173633,CarilloEntropy,CARRILLO199993,DroniouEymardTalbot,LuJaeger, pop2011regularization,ZouExistence}. 
%\KM{[Some comments on existence results for this equation (Sorin?)]}

\subsection{Residual formulation}\label{subsec:Residual-formulation}
To prepare for the space-time convergence and error analysis, we introduce a residual notation for Problem \ref{eq:reformulated_weak_form}. More precisely, we associate a residual operator $\calR : \calS 
\;\longrightarrow\; L^2(0,T;H^{-1}(\Om))$ to \eqref{eq:weak form}. Specifically, given $s \in \mathcal{S}$, $\calR(s) \in L^2(0,T;H^{-1}(\Om))$ is defined as
\begin{align}\label{eq:residual}
   \calR(s)( \varphi) := \int_0^T\!\left\langle \partial_t b(s), \varphi \right\rangle 
  + \int_0^T\!(\nabla B(s), \nabla \varphi)
  - \int_0^T\!( f, \varphi )
\end{align}
for all $\varphi \in L^2(0,T;H_0^1(\Omega))$. Observe that $\calR(s)=0$ if and only if $s$ solves Problem \ref{eq:reformulated_weak_form}. For future reference, we also define the operator for the heat equation $\resH:\calV\to L^2(0,T;H^{-1}(\Om))$ as 
\begin{align}\label{eq:resH_def}
        \resH(s)(\varphi)   
  := \int_0^T\!\left\langle \partial_t s, \varphi \right\rangle 
  + \int_0^T\!(\nabla s, \nabla \varphi).
\end{align}
The following result [Theorem 2.1 of \cite{ErnSmearsVohralik2017}] will be used multiple times.
%in this work:
\begin{lemma}\label{lemma_2.2} For any $s\in \calV$,
\[\|s\|_{\calV}^2 =\int_0^T \|\resH(s)\|^2_{L^2(0,T;H^{-1}(\Omega))} + \|s(0)\|^2.\]\label{lemma:resH}
\end{lemma}
\vspace{-15mm}
\begin{proof}
 Let \( G_s \in L^2(0,T;H_0^1(\Omega)) \)  satisfy by Riesz representation theorem 
\begin{align}\label{eq:GF_test}
    \int_0^T \left(\nabla G_s,\nabla \psi\right) = \int_0^T \left(s, \psi\right).
\end{align}
for all $\psi \in L^2(0,T;H_0^1(\Omega))$.
Observe that, for a.e.\ $t$, $G_s(t)$ is the Green's function associated to the Laplace equation with homogeneous Dirichlet boundary condition, i.e., the solution to $-\Delta G_s(t) = s(t)$ in $\Omega$. Since $s\in H^1(0,T;H^{-1}(\Omega))$, one obtains $\partial_t G_s \in L^2(0,T;H_0^1(\Omega))$, and for all $\psi \in L^2(0,T;H_0^1(\Omega))$
\begin{align}\label{eq:nabla_GF}
\int_0^T (\nabla \partial_t G_s, \nabla \psi)
= \int_0^T \langle \partial_t s, \psi\rangle,\text{ implying }
\|\partial_t G_s\|_{L^2(0,T;H^1_0(\Omega))}
\overset{\eqref{eq:weak_poisson}}{=}
\|\partial_t s\|_{L^2(0,T;H^{-1}(\Omega))}.
\end{align}
Then, for the residual $\resH$ one has $\resH(s)(\psi)= \int_0^T \left( \nabla ( \partial_t G_s + s), \nabla \psi \right)$. 
Therefore,
\begin{align*}
&\| \resH(s) \|^2_{L^2(0, T;H^{-1}(\Om))}\overset{\eqref{eq:weak_poisson}}= \int_0^T \| \nabla ( \partial_t G_s + s) \|^2 = \int_0^T \left( \| \nabla \partial_t G_s \|^2 + \| \nabla s\|^2 + 2 ( \nabla \partial_t G_s, \nabla s ) \right)\\
&\overset{\eqref{eq:nabla_GF}}= \|\p_t s\|_{L^2(0,T;H^{-1}(\Om))}^2 + \|\nabla s\|_{L^2(Q)}^2 + 2 \int_0^T \langle \partial_t s, s \rangle \, \overset{\eqref{eq:norm_V}}= \|s\|_\calV^2 - \|s(0)\|^2. 
\end{align*}
In the last step, we have used the equality $2\int_0^T\langle \p_t s, s\rangle = \int_0^T \frac{d}{dt} \|s\|^2=\|s(T)\|^2 - \|s(0)\|^2$. 
%This concludes the proof.
\end{proof}
% Thus, \eqref{eq:linearization_gen} can be equivalently rewritten in incremental form as
% \begin{align}\label{eq:linearized_equation}
%     \int_0^T \!\left\langle \partial_t (s^{i+1} - s^i), \varphi \right\rangle   
%     + \int_0^T \!(\nabla (s^{i+1} - s^i), \nabla \varphi) 
%     = -\frac{1}{L} \int_0^T \!\langle \mathcal{R}(s^i), \varphi \rangle, 
% \end{align}
% for all $\varphi \in \mathcal{V}$.\\ 
% Note that, given $s^i \in \mathcal{S}$, Problem \ref{pb:2} is linear. 


% This residual form will later serve as the basis for the decomposition of errors (see~\cref{subsec:decomposition}).







% \begin{problem}[{Space-time residual formulation of Problem~\ref{pb:2}}]\label{pb:residual_formulation}
% Given $s^i\in \mathcal{V}$ be an approximation of $ \mathcal{S}$. We want to compute $s^{i+1} \in \mathcal{V}$ that solves the linearized problem \eqref{eq:linearized_equation}. Due to discretization, we compute only $s^{i+1}_{h\tau}\in\mathcal V$,
% %Given $s^i \in S$ from the previous iteration, find 
% %$s^{i+1}_{h\tau} \in V$ 
% such that, for all $\varphi \in L^2(0,T;H_0^1(\Omega))$, the following holds:
% \begin{align}\label{eq:R_H_problem}
% \int_0^T \!\Big\langle \partial_t(s^{i+1}_{h\tau}-s^i), \varphi \Big\rangle
% + \int_0^T \!\big(\nabla(s^{i+1}_{h\tau}-s^i), \nabla\varphi\big)
% = -\frac{1}{L}\!\langle \mathcal{R}(s^i), \varphi \rangle
% + \int_0^T \!\langle R_{\mathrm lin}(s^{i+1}_{h\tau}, s^i), \varphi \rangle.
% \end{align}
% Here, the operator $R_{\mathrm lin}(s^{i+1}_{h\tau}, s^i)$ denotes the space-time (heat) residual associated with the discrete approximation and is defined by
% \begin{align}\label{eq:R_H_def}
% \int_0^T \!\langle R_{\mathrm lin}(s^{i+1}_{h\tau}, s^i), \varphi \rangle
% :=
% \int_0^T \!\Big\langle \partial_t(s^{i+1}_{h\tau}-s^i), \varphi \Big\rangle
% + \int_0^T \!\big(\nabla(s^{i+1}_{h\tau}-s^i), \nabla\varphi\big)
% + \frac{1}{L} \!\langle \mathcal{R}(s^i), \varphi \rangle.
% \end{align}
% This formulation is completely equivalent to the incremental form 
% \eqref{eq:linearized_equation}. 
% \end{problem}
% \noindent
% \Cref{subsec:Residual-formulation} expresses the discrete linearized scheme~\eqref{eq:linearization_gen} in
% a compact residual form that will serve as the basis for the forthcoming space–time error decomposition.

\section{Linearization scheme}\label{Linearization}
Note that Problem \ref{eq:reformulated_weak_form} is nonlinear in both $b(s)$ and $B(s)$; we consider a linear iterative scheme to approximate its solution.
With $i \in \mathbb{N}$ denoting the iteration index, we define the sequence of approximations \(\{s^i\}_{i \in \mathbb{N} }\) as follows. Given $s^i \in \mathcal{S}$, the next iterate solves the following problem. 
\begin{problem}[{Space-time L-scheme for Problem \ref{eq:reformulated_weak_form}}]\label{pb:2}
Let $i\in \N$ be fixed, and $L>1$ be a constant. Assume that $s_0 \in L^2(\Omega)$ can be taken such that $b(s_0) = u_0$, and $s^i \in  (\mathcal{S}\cap \calV)$ is known, satisfying $s^i(0)=s_0$. Find $s^{i+1} \in \mathcal{V}$ such that $s^{i+1}(0)=s_0$ and for all $\varphi \in  L^2(0,T;H_0^1(\Omega))$, one has
%Let $s_0 \in H^1_0(\Om)$ satisfy $b(s_0)=u_0$ and $s^0\in \calW$ satisfy $s^0(0)=s_0$. Let $L>1$ be a constant. 
%Assume $\{s^j\}_{j\leq i} \subset \mathcal{S}$ is known for $i\in \N$. Find $s^{i+1} \in \mathcal{V}$ such that $s^{i+1}(0)=s_0$ and for all $\varphi \in  L^2(0,T;H_0^1(\Omega))$,
\begin{align}\label{eq:linearization_gen}
\int_0^T \!\left\langle \partial_t \left( L(s^{i+1} - s^i) + b(s^i) \right), \varphi \right\rangle \,  
+ \int_0^T \!\left( \nabla \left( L(s^{i+1} - s^i) + B(s^i) \right), \nabla \varphi \right) 
= \int_0^T \!( f, \varphi ).
\end{align}
\end{problem}
Observe that \eqref{eq:linearization_gen} is nothing but the linear heat equation in a weak form. This linearization approach is inspired by the $L$-scheme \cite{list2016study,pop2004mixed} developed for elliptic problems. This particular implementation with $b$-$B$ functions is inspired by the double splitting scheme in \cite{javed2025robust}. By \eqref{eq:condition}, $b$ and $B$ are Lipschitz continuous, both having $1$ as 
Lipschitz constant. To ensure convergence, one needs to take the stabilization parameter $L \geq 1$, and the proof of \Cref{theo:converegence} below gives the additional restriction $L < 2$. Therefore, we take $L \in [1, 2)$, and in practice use $L = 1$. We define the linearized residual
\[
\mathcal{R}_{\mathrm{lin}} :
\mathcal{V}\times(\mathcal{S}\cap\mathcal{V})
\to L^2(0,T;H^{-1}(\Omega)),
\]
\begin{align}\label{eq:defRlin}
\mathcal{R}_{\mathrm{lin}}(s,\bar{s})&:= L \mathcal{R}_{\mathrm{H}}(s-\bar{s}) + \mathcal{R}(\bar{s}),
\end{align}
with $\mathcal{R}_{\mathrm{H}}$ introduced in \eqref{eq:resH_def}.
\begin{remark}[Operator definition of the L-scheme] Given $s^i \in \mathcal{S}\cap\mathcal{V}$, Problem \ref{pb:2} can be interpreted as to find $s^{i+1} \in \mathcal{V}$ satisfying $s^{i+1}(0)=s_0$ such that
\begin{align}
&\mathcal{R}_{\mathrm{lin}}(s^{i+1},s^i)= 0.
\label{eq:l_b}
\end{align}
\end{remark}
%\KM{Can be removed }We start the iteration from some $s^0 \in \mathcal{W}$, in particular, if the exact solution $s$ belongs to $\mathcal{W}$, one may take $s^0 = s(0)$.
%Initially, one may take any function in $\mathcal{S}$, e.g. $s^0 = s(0)$.  

%\vskip 0.3cm
%\noindent
%Here, $L>0$ is a stabilization parameter whose choice influences the behavior of the iterative method.  
%In particular, since $b$ and $B$ are assumed to satisfy \eqref{eq:condition}, choosing $L>1$ yields an iteration that is
%%monotone and converges regardless of the initial guess, see
%similar in spirit to the classical $L$–scheme.
%\cite{list2016study,pop2004mixed}, and in particular, to the double splitting scheme considered in \cite{javed2025robust}.
The main results regarding well-posedness and the convergence of the scheme are summarized in the following theorems.
\begin{theorem}[Well-posedness of the L-scheme]
\label{Convergence_theorem}
Let $s^i \in \mathcal{S}$ be given. Then, for any stabilization parameter $L\geq1$, 
the linearized problem~\eqref{eq:linearization_gen} admits a unique solution $s^{i+1} \in \mathcal{V}$. 
If, in addition, $s^i \in \mathcal{W}$, then $s^{i+1} \in \mathcal{W}\subset (\calS\cap \calV)$, and the iteration preserves regularity.
\end{theorem}
\noindent
Given $s^i \in \mathcal{S}$, the existence and uniqueness of $s^{i+1} \in \mathcal{V}$ are well known (see, e.g., \cite[Chapter~7]{evans2022partial}). Moreover, if $s^i\in \calW$, then $\calR(s^i)$ can be extended to a linear functional on $L^2(Q)$. The embedding $\calW\subset C(0,T;H^1_0(\Om))$ gives $s^i(0)\in H^1_0(\Om)$. Using regularity results from \cite[Chapter 7]{evans2022partial}, one also gets $s^{i+1}\in \calW$. 


%The proof of the above statement follows from classical arguments. For $s^i\in \calS$, we have $\calR(s^i)\in L^2(0,T;H^{-1}(\Om))$ by definition of $\calS$ in \eqref{eq:spaces}. Then, the existence of $s^{i+1}\in \calV$ solving Problem \ref{pb:2} is evident from \cite[Chapter 7]{evans2022partial}.  Moreover, if $s^i\in \calW$, then $\calR(s^i)\in L^2(Q)$. Due to the embedding $\calW\subset C(0,T;H^1_0(\Om))$ it further holds that $s^i(0)\in H^1_0(\Om)$. This gives $s^{i+1}\in \calW$ again using regularity results from \cite[Chapter 7]{evans2022partial}.

\begin{theorem}[Convergence of the L-scheme]\label{theo:converegence}
Take $L \in (1,2)$ and let $s^0\in \calW$ be an initial guess. Further, let $\{s^i\}_{i\ge 0}\subset \calW$ be the sequence generated by the $L$-scheme, and $s\in \calS$ be the weak solution to Problem~\ref{eq:reformulated_weak_form}. For every $t\in(0,T]$, following convergence results hold as $i \to \infty$
\begin{subequations}\label{eq:L_conv_1}
    \begin{align}
\|s^{i+1}-s^{i}\| &\longrightarrow 0
&&\text{in } L^2(Q),\label{eq:L_conv_1a}\\
b(s^i) &\longrightarrow b(s)
&&\text{in } L^2(0,T;H^{-1}(\Omega)),\\
\int_0^t B(s^i) &\longrightarrow \int_0^t B(s)
&&\text{in } L^\infty(0,T;L^2(\Omega)),\\
\int_0^t [B(s^{i-1}) + L(s^i-s^{i-1})]& \longrightarrow \int_0^t B(s) &&\text{in } L^\infty(0,T;H^1_0(\Omega)).
\end{align}
\end{subequations}
\textbf{[Improved convergence]:} If, in addition, the inverse mapping $\Phi^{-1}$ is $\a$-H\"older continuous for some $\a\in (0,1)$, then one has
\begin{align}\label{eq:L_conv_2}
b(s^i) &\longrightarrow b(s)
\quad &&\text{in } L^{1 + 1/\a}(Q).
\end{align}

\textbf{[Contraction]:}
If \( \ell := \inf_{r\in\mathbb{R}} \min\{b'(r),\, B'(r)\} > 0 \), then the iteration error is contractive,
and the iterative scheme converges linearly in ${L^2(Q)}$. 
\end{theorem}

First observe that $\|s^{i+1}-s^{i}\| \longrightarrow 0$ in \eqref{eq:L_conv_1a} does not imply the convergence of $s^i$ to $s$. However, it is used in the proofs for the other convergence results; therefore, we give it here explicitly. Also, observe that the H\"older continuity of $\Phi^{-1}$ holds if $\Phi$ has a porous medium equation type growth, i.e., $\Phi(u)\sim u^m$ for some $m\geq 1$. In this case, $\a=1/m$ since 
\[
|u_1-u_2|^m \leq |u_1^m-u_2^m|\lesssim |\Phi(u_1)-\Phi(u_2)|\;\; \text{ or }\;\; |\Phi^{-1}(w_1)-\Phi^{-1}(w_2)|\lesssim |w_1-w_2|^{1/m}.
\]
This covers all the cases in \Cref{fig:RichVsBio_phi}.
We prove \Cref{theo:converegence} below.  

\subsection{Convergence Analysis}\label{sec:convergence-proof}

We give here the proof for \Cref{theo:converegence}. To begin, we introduce the errors at the $i$-th iterate
\begin{align}\label{eq:def_errors_a}
e_s^i := s^i - s, \quad e_b^i := b(s^i) - b(s), \quad e_B^i := B(s^i) - B(s).
\end{align}
Subtracting ~\eqref{eq:linearization_gen} from ~\eqref{eq:weak form}, and rearranging the resulting terms, yields
\begin{align}\label{eq:linearized_error_form_eq}
\int_0^T \left\langle \partial_t \left( L(e_s^{i+1} - e_s^i) + e_b^i \right), \varphi \right\rangle \,  
+ \int_0^T \left( \nabla \left( L\left(e_s^{i+1} - e_s^i\right) + e_B^i \right), \nabla \varphi \right) \, = 0.
\end{align}
Observe that $e^i_s(0) = e^i_b(0) = 0$ for any $i \in \mathbb{N}$. With $\chi_{(0,t)}$ denoting the characteristic function of $(0,t)$, for any $t \in (0,T]$ one can take in \eqref{eq:linearized_error_form_eq} $\varphi = \chi_{(0,t)}\psi$ with $\psi \in H^1_0(\Omega)$, using $(L(e^{i+1}_s -e^i_s) + e^i_b)(0)=0$ and $\int_0^t \p_t(L(e^{i+1}_s -e^i_s) + e^i_b)(\t)\,\dd \t= (L(e^{i+1}_s -e^i_s) + e^i_b)(t)$, to obtain
\begin{align}\label{eq:time-integrated-error}
\left\langle \left( L(e_s^{i+1} - e_s^i) + e_b^i \right)(t), \psi \right\rangle \,  
+ \left( \nabla \int_0^t \left( L\left(e_s^{i+1} - e_s^i\right) + e_B^i \right), \nabla \psi \right) \, = 0.
\end{align}
Therefore, the Riesz representative of $(L(e^{i+1}_s -e^i_s) + e^i_b)(t)$ is
\begin{align}\label{eq:green_def_new}
    \mathcal{G}^{i+1}(t) 
    := -\int_0^{t} 
    \bigl(L(e_s^{i+1} - e_s^i)(\tau) + e_B^i(\tau)\bigr)
    \,\dd\tau.
\end{align}
From \eqref{eq:weak_poisson}, one gets for a.e.\ $t\in (0,T]$
%Hence, combining \eqref{eq:time-integrated-error} with \eqref{eq:weak_poisson}, we get for all $t\in (0,T]$
%\begin{align}\label{eq:green_riesz}
%\left\{\begin{aligned}
%   &\bigl(\nabla\mathcal{G}^{i+1}(t), \nabla\psi\bigr) = 
%\bigl\langle L(e_s^{i+1}-e_s^i)(t)+e_b^i(t), \psi
%\bigr\rangle\qquad \forall\, \psi\in H_0^1(\Omega), \\
%&\text{ implying } \quad \|\nabla\mathcal{G}^{i+1}(t)\| = 
%    \|L(e_s^{i+1}-e_s^i)(t)+e_b^i(t)\|_{H^{-1}(\Omega)}. 
%\end{aligned}
%\right.
%\end{align}
\begin{align}\label{eq:green_riesz}
\bigl(\nabla\mathcal{G}^{i+1}(t), \nabla\psi\bigr) = 
\bigl\langle L(e_s^{i+1}-e_s^i)(t)+e_b^i(t), \psi
\bigr\rangle, 
\end{align}
for all $ \psi\in H_0^1(\Omega)$, and $\|\nabla\mathcal{G}^{i+1}(t)\| =  \|L(e_s^{i+1}-e_s^i)(t)+e_b^i(t)\|_{H^{-1}(\Omega)}$. Furthermore, since $L(e_s^{\,i+1}-e_s^{\,i}) + e_b^{\,i}\in H^1\big(0,T;H^{-1}(\Omega)\big)$, one has $\mathcal{G}^{\,i+1} \in H^1\big(0,T;H_0^1(\Omega)\big)$.

The remaining part of the proof is divided into three steps.

\textbf{(Step 1) An inequality involving successive iterations:} Inserting $\varphi =\chi_{(0,\tilde{t})} \mathcal{G}^{i+1}$ into~\eqref{eq:linearized_error_form_eq}  gives
\begin{align}\label{eq:Main_new}
\int_0^{\tilde{t}} \left\langle \partial_t \left( L(e_s^{i+1} - e_s^i) + e_b^i \right), \mathcal{G}^{i+1} \right\rangle \,  + \int_0^{\tilde{t}} \left( \nabla \left( L\left(e_s^{i+1} - e_s^i\right) + e_B^i \right), \nabla \mathcal{G}^{i+1} \right) = 0.
\end{align}
Denoting the terms on the left by $T_1$ and $T_2$ respectively, one has
\begin{align}
 T_1 &:= \int_0^{\tilde{t}} \left\langle \partial_t \left( L(e_s^{i+1} - e_s^i) + e_b^i \right), \mathcal{G}^{i+1} \right\rangle \overset{\eqref{eq:green_riesz}}{=} 
\int_0^{\tilde{t}} (\nabla \partial_t \mathcal{G}^{i+1}, \nabla \mathcal{G}^{i+1})= \frac{1}{2} \int_0^{\tilde{t}} \frac{d}{dt} \big\|\nabla \mathcal{G}^{i+1}(t)\big\|^2\,dt \nonumber\\
&= \frac{1}{2}\big\|\nabla \mathcal{G}^{i+1}({\tilde{t}})\big\|^2\overset{\eqref{eq:green_riesz}}{=}\frac{1}{2} \big\|(L(e_s^{i+1}-e_s^i)+e_b^i)({\tilde{t}})\big\|_{H^{-1}(\Omega)}^2.\label{T_1_energy_new}
\end{align}
In the above, we used  $\mathcal{G}^{i+1}(0)=0$ that follows from~\eqref{eq:green_def_new}.
We now consider the second term in \eqref{eq:Main_new}, expanded as
\begin{align*}
T_2 & := \int_0^{\tilde{t}} \left( \nabla \left( L\left(e_s^{i+1} 
- e_s^i\right) + e_B^i \right), \nabla \mathcal{G}^{i+1} \right) \overset{\eqref{eq:green_riesz}}{=} 
\int_0^{\tilde{t}} \Bigl( L\left(e_s^{i+1} - e_s^i\right) + e_b^i,\, 
L\left(e_s^{i+1} - e_s^i\right) + e_B^i \Bigr)\\
&= L^2 \int_0^{\tilde{t}} \| e_s^{i+1} - e_s^i \|^2
+L \int_0^{\tilde{t}} \left( e_B^i +e_b^i, \,
e_s^{i+1} - e_s^i\right) 
+ \int_0^{\tilde{t}} \left( e_B^i, e_b^i \right) \\
&\overset{\eqref{eq:def_errors_a},\eqref{eq:def_errors_b}}{=} 
L^2\int_0^{\tilde{t}} \|e_s^{i+1} - e_s^i\|^2 +L \int_0^{\tilde{t}} \bigl((b[s^i,s]+B[s^i,s])e_s^i,\,
e_s^{i+1} - e_s^i\bigr)
+\int_0^{\tilde{t}} \bigl(e_s^i b[s^i,s],\,e_s^i B[s^i,s]\bigr).
\end{align*}
Applying the elementary identity $a(a-b)=\tfrac{1}{2}(a^2-b^2+(a-b)^2)$, with $a=e_s^{i+1}$ and $b=e_s^{i}$ to the second term (mixed product), one gets after combining terms
\begin{align}\label{eq:T2_rewritten_new}
T_2
&= \frac{L}{2}\int_0^{\tilde{t}} \int_\Omega 
    (b[s^i,s] + B[s^i,s])\, |e_s^{i+1}|^2\,dx\,dt\nonumber\\ 
  &\quad      
  +\frac{L}{2}\int_0^{\tilde{t}} \int_\Omega 
    \big( 2L - (b[s^i,s] + B[s^i,s]) \big)\, 
    |e_s^{i+1} - e_s^i|^2\,dx\,dt       \nonumber\\ 
  &\quad 
  + \frac{1}{2}\int_0^{\tilde{t}} \int_\Omega 
    \big( 2\, b[s^i,s]\, B[s^i,s] 
    - L(b[s^i,s] + B[s^i,s]) \big)\,
    |e_s^i|^2\,dx\,dt.
\end{align}
Inserting~\eqref{T_1_energy_new}--\eqref{eq:T2_rewritten_new} 
into~\eqref{eq:Main_new} and multiplying by $2$ gives
\begin{align}\label{eq:Main_equality_new}
    &\left\| \nabla \mathcal{G}^{i+1}({\tilde{t}})\right\|^2
    + \int_0^{\tilde{t}} \int_\Omega L \left( b[s^i,s] + B[s^i,s] \right) \left| e_s^{i+1} \right|^2 \nonumber \\
    &+ \int_0^{\tilde{t}} \int_\Omega \left( 2L^2 - L\left( b[s^i,s] + B[s^i,s] \right) \right) 
    \left| e_s^{i+1} - e_s^i \right|^2 
    + (2-L)\int_0^{\tilde{t}} \int_\Omega  b[s^i,s] B[s^i,s]  \left| e_s^i \right|^2 \nonumber \\
    &= \int_0^{\tilde{t}} \int_\Omega L\left( b[s^i,s] + B[s^i,s]-  b[s^i,s] B[s^i,s] \right) \left| e_s^i \right|^2.
\end{align}
We identify three coefficient functions 
in~\eqref{eq:Main_equality_new} as
\begin{subequations}\label{eq:P-coeff_new}
\begin{align}
P_1^i
&:=  L \left( b[s^i,s] + B[s^i,s] \right)
\overset{\eqref{eq:def_errors_b}}{=}
L \frac{(b + B)(s^i) - (b + B)(s)}{s^i - s}
\overset{\eqref{eq:condition}}{\geq} L, \\
P_2^i
&:= 2L^2 - L\left( b[s^i,s] + B[s^i,s] \right)
\overset{\eqref{eq:condition}}{\geq} 2L(L-1), \\
P_3^i
&:= L\left( b[s^i,s] + B[s^i,s]-  b[s^i,s] B[s^i,s] \right) = L - L(1 - b[s^i, s])(1 - B[s^i, s]) \overset{\eqref{eq:condition}}\in [0,L]. 
\end{align}
\end{subequations}
Using~\eqref{eq:P-coeff_new} in~\eqref{eq:Main_equality_new}, one has the following inequality, which is fundamental for convergence
\begin{align}\label{eq:final_estimate_new}
    &L \int_0^{\tilde{t}} \|e_s^{i+1}\|^{2}
    +\left\| \nabla \mathcal{G}^{i+1}({\tilde{t}})\right\|^2+
    (2-L)\int_0^{\tilde{t}} \int_\Omega e^i_B e^i_b\,dx\,dt + L(L-1)\int_0^{\tilde{t}} \|s^{i+1} -s^{\,i}\|^{2}\,dt
    \nonumber \\ 
    &\le\;
    L\int_0^{\tilde{t}} \|e_s^{\,i}\|^{2}\,dt.
\end{align}
\textbf{(Step 2) Proof of \eqref{eq:L_conv_1}:} 
Adding ~\eqref{eq:final_estimate_new} for $i=0$ to $N$, one has
\begin{align}\label{telescoped_estimate}
&L \int_0^{\tilde{t}} \|e_s^{N+1}\|^2\,dt + 
\sum_{i=0}^{N} 
\left\| \nabla \mathcal{G}^{i+1}({\tilde{t}})\right\|^2
+ (2-L)\sum_{i=0}^{N} 
\int_0^{\tilde{t}} \int_\Omega e_B^i\, e_b^i\,dx\,dt + L(L-1)\sum_{i=0}^{N} 
\int_0^{\tilde{t}} \|s^{i+1}-s^i\|^2\,dt
\nonumber\\
&\le L \int_0^{\tilde{t}} \|e_s^0\|^2\,dt.
\end{align}
Since ${\tilde{t}} \in (0,T]$ was chosen arbitrarily, and $\mathcal{G}^{i+1} \in  H^1(0,T;H^1_0(\Omega))$, we immediately get that for every $i\in \mathbb{N}$ and ${\tilde{t}} \in [0,T]$,
\begin{align}
    \left\| \nabla \mathcal{G}^{i}(\tilde{t})\right\|^2 \le L \int_0^{T} \|e_s^0\|^2\,dt,
\end{align}
so the sequence $\{\mathcal{G}^{i}\}_{i \in \mathbb{N}}$ has an equibounded norm in $L^{\infty}(0,T;H^1_0(\Omega))$. This immediately implies the same for the norm in $L^{\infty}(0,T;H^{-1}(\Omega))$ of $\{L(s^{i+1}-s^i) + e_b^i\}_{i \in \mathbb{N}}$. 

Further~\eqref{telescoped_estimate} implies the absolute convergence of the series, yielding the  convergences
\begin{align}\label{eq:inq_new}
&\|\left(L(s^{i+1} - s^i) + e_b^i\right)({\tilde{t}})
\|^2_{H^{-1}(\Omega)} = \left\|\nabla \int_0^{\tilde{t}} (L(s^{i+1} - s^i) + e_B^i)(\t)\, d\t\right\|= \left\| \nabla \mathcal{G}^{i+1}({\tilde{t}})\right\|\to 0,
\nonumber\\
&\int_0^{\tilde{t}} \|s^{i+1}-s^{i}\|^2\,dt \to 0, \qquad\text{and}\qquad
\int_0^{\tilde{t}}\int_\Omega e_b^i\,e_B^i\,dx\,dt \to 0,
\end{align}
as $i\to\infty$. Since $\tilde{t} \in (0,T]$ was chosen arbitrarily, one gets the convergence
\begin{align}
\|s^{i+1} - s^i\|_{L^2(Q)} \to 0 \quad \text{as } i \to \infty.
\end{align}
To prove that $e_b^i$ vanishes as $i \to \infty$, we recall the equiboundedness in $L^{\infty}(0,T;H^{-1}(\Omega))$ of 
$L(s^{i+1}-s^i) + e_b^i$. Using that for every $\tilde{t}\in (0,T]$ one has $\left\|\left(L(s^{i+1}-s^i) + e_b^i\right)(\tilde{t})\right\|_{H^{-1}(\Omega)} \to 0 \quad \text{as } i \to \infty$, the dominated convergence theorem, together with the equiboundedness, guarantees that 
\begin{align}
\left\|L(s^{i+1}-s^i) + e_b^i\right\|_{L^2(0,T;H^{-1}(\Omega))} 
\to 0 \quad \text{as } i \to \infty.
\end{align}
With this, we apply the triangle inequality and \eqref{eq:defHinv_norm} to obtain
%\begin{align}
%&\|e_b^i\|_{L^2(0,T;H^{-1}(\Omega))}
%\le 
%\|L(s^{i+1}-s^i)+e_b^i\|_{L^2(0,T;H^{-1}(\Omega))}
%+L\,\|s^{i+1}-s^i\|_{L^2(0,T;H^{-1}(\Omega))} \nonumber\\
%&\quad \leq \sqrt{T}\sup_{0<t\leq T} \|\left(L(s^{i+1} - s^i) + e_b^i\right)(t)\|_{H^{-1}(\Om)} \overset{\eqref{eq:defHinv_norm}}+LC_{\,\!_{\Omega}}^{\mathrm{P}} \, h_\Om\,\|s^{i+1}-s^i\|_{L^2(Q)} .\label{eq:b_convergence}
%\end{align}
\begin{align}
&\|e_b^i\|_{L^2(0,T;H^{-1}(\Omega))}
\le 
\|L(s^{i+1}-s^i)+e_b^i\|_{L^2(0,T;H^{-1}(\Omega))}
+L\,\|s^{i+1}-s^i\|_{L^2(0,T;H^{-1}(\Omega))} \nonumber\\
&\quad \leq \|L(s^{i+1} - s^i) + e_b^i\|_{L^2(0,T;H^{-1}(\Omega))} 
\overset{\eqref{eq:defHinv_norm}}{+}
LC_{\,\!_{\Omega}}^{\mathrm{P}} \, h_\Omega\,
\|s^{i+1}-s^i\|_{L^2(Q)}.\label{eq:b_convergence}
\end{align}
Since both terms on the right are vanishing  as $i \to \infty$, one gets that% due to \eqref{eq:inq_new}, we have as $i\to \infty$,

\begin{align}
    \|e_b^i\|_{L^2(0,T;H^{-1}(\Omega))}  \;\to\; 0,
\qquad\text{i.e.,}\qquad
b(s^i)\to b(s) \quad \text{in } L^2(0,T;H^{-1}(\Omega)).\label{eq:bconv_new}
\end{align}
Next, by \eqref{eq:inq_new} one gets that for every $t\in(0,T]$,
\begin{align}
\left\|\nabla \int_0^t \left(L(s^{i+1}-s^i)+e_B^i\right)(\tau)\,d\tau\right\| \to 0 \quad \text{as } i \to \infty.
\end{align}
By Poincar\'e's inequality,
\begin{align}
\left\|\int_0^t \left(L(s^{i+1} - s^i) + e_B^i\right)(\tau)
\,d\tau\right\| \to 0 \quad \text{as } i \to \infty.
\end{align}
Further, since $\|s^{i+1}-s^i\|_{L^2(Q)} \to 0 \quad \text{as } i \to \infty$, one immediately gets that
\begin{align}
 \left\|\int_0^t (s^{i+1}-s^i)(\tau)\,d\tau\right\|_{L^2(\Omega)}\to 0 \quad \text{as } i \to \infty.  
\end{align}
 Using the triangle inequality again gives the convergence
\begin{align}
    \int_0^t B(s^i)\,d\tau \;\longrightarrow\; \int_0^t B(s)\,d\tau,
\end{align}
in $L^2(\Omega)$, for every $t\in (0,T]$. This completes the proof of \eqref{eq:L_conv_1}.





%in $H^1_0(\Om)$ due to \eqref{eq:inq_new} implies that
%\begin{align}
%    \int_0^T \bigl(L(s^{i+1} - s^i) + e_B^i\bigr)(\tau)\,d\tau 
%    \;\longrightarrow\; 0 
%    \quad \text{in } L^2(\Omega).
%\end{align}
%Since $\|s^{i+1} - s^i\|_{L^2(Q)} \to 0$ one has that 
%$\int_0^T (s^{i+1} - s^i)(t)\,dt  \to 0$ in $L^2(\Omega)$ since
%\begin{align*}
%    \left\|\int_0^T (s^{i+1} - s^i)(t)\,dt\right\|^2
%    &= \int_\Omega \left(\int_0^T (s^{i+1} - s^i)\,dt\right)^2 dx \leq \int_\Omega \left(\int_0^T dt\right)\left(\int_0^T (s^{i+1} - s^i)^2\,dt\right) dx \\
%    &\leq T \int_0^T \int_\Omega (s^{i+1} - s^i)^2\,dx\,dt 
 %   \;\leq\; T\,\|s^{i+1} - s^i\|_{L^2(Q)}^2.
%\end{align*}
%One concludes that $\int_0^T (e_s^{i+1} - e_s^i) \to 0$ in $L^2(\Omega)$.
%This also implies as $i\to \infty$,
%\begin{align}
 %   \int_0^T B(s^i)\,d\tau \;\longrightarrow\; \int_0^T B(s)\,d\tau 
%    \quad \text{ in }  L^2(\Omega).
%\end{align}
%Since $T>0$ is arbitrary, the above convergence result   along with \eqref{eq:inq_new} and \eqref{eq:bconv_new} hold for all $t\in (0,T]$.
%This proves \eqref{eq:L_conv_1}.

\textbf{(Step 3) Proving \eqref{eq:L_conv_2} and the contractive behaviour:} 
From  \eqref{telescoped_estimate} one also gets that
%The estimate \eqref{eq:final_estimate_new} further implies after summing telescopically that 
\[
\int_Q (B(s^i)-B(s))(b(s^i)-b(s)) 
= \int_Q e^i_B e^i_b \longrightarrow 0.
\]
However, due to the $\alpha$-H\"older continuity of $\Phi^{-1}$ one has
\begin{align*}
    &\int_Q |b(s^i)-b(s)|^{1+1/\alpha}
    \lesssim \int_Q |b(s^i)-b(s)|
    |\Phi^{-1}(b(s_1)) - \Phi^{-1}(b(s_2))|\\
    &= \int_Q (b(s^i)-b(s))(B(s^i)-B(s)) 
    \longrightarrow 0.
\end{align*}
This shows the strong convergence of $b(s^i)$ to $b(s)$ in $L^{1+1/\alpha}$. Finally, if $\ell:= \inf\{B',b'\}>0$, then we have 
\[
\int_Q (B(s^i)-B(s))(b(s^i)-b(s))
\geq \ell \int_Q \|e^i_s\|^2.
\]
Inserting this into~\eqref{eq:final_estimate_new} we have
\begin{align}
    &\int_0^T \|s^{i+1} -s\|^{2}
    +\frac{1}{L}\left\| \nabla \int_0^T\left (L\big(s^{i+1} - s^i\big) + B(s^i)-B(s)\right)\right\|^{2}\leq 
    \left (1-\ell \left(\frac{2-L}{L}\right)\right)\int_0^T \|s^{\,i}-s\|^{2}.
\end{align}
This shows that the iterative method is contracting the $L^2(Q)$ error in $s$ for every iteration step.

%is contractive.\qed

\subsubsection{Numerical study on L-scheme convergence}\label{sec:numerical_validation}
Our goal is to develop an adaptive space-time iterative scheme, based on a posteriori estimates and employing the $L$-type scheme introduced through Problem~\ref{pb:2}. Before proceeding with the adaptivity and a posteriori estimates, we assess the behaviour of this linearization by investigating the decay of the iteration errors. To this end, we compare it with the $N_{\mathrm{reg}}$-scheme
\cite{mitra2019modified,javed2025robust}, which can be seen as a regularized variant of the Newton scheme; see \eqref{eq:M-scheme}. For both schemes, the space-time discretization presented in \Cref{sec:apost} is used on a fixed mesh. The experiments are for a one-dimensional setting with Barenblatt data described in \Cref{sec:1DPME_1DST}.
%We investigate the decay of iteration errors for the \(L\)-scheme and a regularized Newton (\(M\)-)scheme for the space-time formulation introduced in \eqref{eq:M-scheme}. The experiments are for a one-dimensional setting with Barenblatt data described in  \Cref{sec:1DPME_1DST}. The discretization is space-time finite elements presented in \Cref{sec:apost} with a fixed mesh.
Figure~\ref{fig:iter-combined} illustrates the evolution of the iteration error with respect to the iteration number for different types of nonlinearities, including degenerate ($b'$, $B'$ vanishes) and non-degenerate cases (when either $b'$ or $B'$ is bounded away from 0). The first row reports the error measured in the norm $\|b(s^{i+1})-b(s^i)\|_{L^2(Q)}$, while the second row shows the error measured in the norm $\|\partial_x(B(s^{i+1})-B(s^i))\|_{L^2(Q)}$. The regularized Newton ($N_{\mathrm{reg}}$) scheme with regularization parameter $\e=0.1$ is employed for the degenerate and single-degenerate cases (left and center columns) since lower values of $\e$ result in the iterations diverging, whereas $\e=0$ is used in the non-degenerate case (right column).





%===============================================================
\begin{figure}[h]
\centering
%==============================================================================================

% ===== Top Row: Gradient of L2 norm of error =====
\begin{subfigure}[t]{0.34\textwidth}
  \centering
  \includegraphics[width=\linewidth]{Figures/Double_comparison_UniL_UniM_bbL2.eps.eps}
  \caption{Degenerate}
%\caption{Convergence of the regularized Newton (\(M\)-)scheme and \(L\)-scheme for the degenerate diffusion case \eqref{eq:bB_PME} on a fixed mesh.}
\end{subfigure}\hfill
\begin{subfigure}[t]{0.33\textwidth}
  \centering
  \includegraphics[width=\linewidth]{Figures/iter_singledeg_M2_H1_bbL2_NEWTON.eps.eps}
 \caption{\centering Single-degenerate.}
 \end{subfigure}\hfill
\begin{subfigure}[t]{0.33\textwidth}
  \centering
  \includegraphics[width=\linewidth]{Figures/nondegerate_comparison_UniL_UniM}
  \caption{\centering Non-degenerate}
\end{subfigure}
  %=======================================

%==============================================================================
% ===== Top Row: Gradient of L2 norm of error =====
\iffalse
\begin{subfigure}[t]{0.34\textwidth}
  \centering
  \includegraphics[width=\linewidth]{Figures/iter_regularizedBs_PME_M2_bbL2_regNewton.eps}
\caption{Degenerate}
\end{subfigure}\hfill
\begin{subfigure}[t]{0.33\textwidth}
  \centering
  \includegraphics[width=\linewidth]{Figures/iter_singledeg_M2_H1_bbL2_NEWTON.eps.eps}
 \caption{\centering Single-degenerate}
 \end{subfigure}\hfill
\begin{subfigure}[t]{0.33\textwidth}
  \centering
  \includegraphics[width=\linewidth]{Figures/iter_nondegenerate_M2_bbL2_Newton_scheme}
  \caption{\centering Non-degenerate}
\end{subfigure}
  %=======================================
 % \begin{subfigure}[t]{0.42\textwidth}
 %   \centering
 %   \includegraphics[width=\linewidth]{Figures/itererr-vs-iter_double-deg.eps}
  %  \caption{Double-degenerate}
   % \end{subfigure}\hfill
   
    % \begin{subfigure}[t]{0.32\textwidth}
    %\centering
    %\includegraphics[width=\linewidth]{Figures/iter_doubledeg_UexactH1.eps}
    %\caption{Double-degenerate}
    %\end{subfigure}\hfill
   %======================================= 
  \vspace{1em} 
  % ---- Middle Row:
  %L2 norm of error ----
  \begin{subfigure}[t]{0.32\textwidth}
    \centering
    \includegraphics[width=\linewidth]{Figures/iter_singledeg_PME_M2_L2}
    \caption{Degenerate}
  \end{subfigure}\hfill
  \begin{subfigure}[t]{0.32\textwidth}
    \centering
    \includegraphics[width=\linewidth]{Figures/iter_singledeg_M2_L2_bbL2_NEW.eps.eps}
 \caption{\centering Single-degenerate}
 \end{subfigure}\hfill
  \begin{subfigure}[t]{0.32\textwidth}
    \centering
    \includegraphics[width=\linewidth]{Figures/iter_nondegenerate_M2_L2}
  \caption{\centering Non-degenerate}
  \end{subfigure}
\vspace{1em} 
\fi
%===============================================
 % ---- Bottom Row:
  %L2 norm of error ----
  \begin{subfigure}[t]{0.34\textwidth}
    \centering
    \includegraphics[width=\linewidth]{Figures/iter_doubledeg_M2_BH1_NEWTON_Lscheme.eps.eps}
    \caption{Degenerate }
  \end{subfigure}\hfill
  \begin{subfigure}[t]{0.33\textwidth}
    \centering    \includegraphics[width=\linewidth]{Figures/single_comparison_UniL_UniM_BBH1_norm.eps.eps}
 \caption{\centering Single-degenerate}
 \end{subfigure}\hfill
  \begin{subfigure}[t]{0.33\textwidth}
    \centering    \includegraphics[width=\linewidth]{Figures/iter_nondegenerate_M2_bbL2_Newton_scheme.eps.eps}
  \caption{\centering Non-degenerate}
  \end{subfigure}
\caption{[\Cref{sec:numerical_validation}] Decay of the linearization error versus the number of iterations for the \(L\)-scheme, and the regularized Newton scheme \eqref{eq:M-scheme} for a fixed mesh. \textbf{Top row:} error measured as $\|b(s^{i+1})-b(s^i)\|_{L^2(Q)}$. \textbf{Bottom row:} error measured as $\|\partial_x(B(s^{i+1})-B(s^i))\|_{L^2(Q)}$. From left to right, the columns correspond to degenerate (porous medium equation, nonlinearity as in \eqref{eq:bB_PME}), single-degenerate ($B(s)=\frac{1}{2}s^2+\frac{1}{2}s,\; b(s)=s^2$ so that $\inf_{\mathbb{R}} B'>0$), and non-degenerate diffusion ($B(s)=\frac{1}{2}s^2+\frac{1}{2}s, \; b(s)=\frac{1}{6}s^3+\frac{1}{2}s$ so that $\inf_{\mathbb{R}} B',\, \inf_{\mathbb{R}} b'>0$) models. The regularized Newton scheme is used in the degenerate and single-degenerate cases with $\e=0.1$, while the classical Newton scheme ($\e=0$) is used in the non-degenerate case.}
\label{fig:iter-combined}
\end{figure}

For the single-degenerate case when $\inf_{\mathbb{R}} B'>0$ and non-degenerate cases, a significantly faster decay of the iteration error is observed for the regularized Newton ($N_{\mathrm{reg}}$) scheme compared to the \(L\)-scheme. In particular, for the non-degenerate case, the Newton scheme exhibits quadratic convergence (since $\e=0$). However, when \(B'(s)\) vanishes, the error decay for the regularized Newton scheme becomes irregular and eventually stagnates, even in the presence of the regularization parameter \(\e>0\), which appears insufficient to compensate for the vanishing diffusion. In contrast, for the degenerate diffusion case, the \(L\)-scheme exhibits a faster initial decay of the iteration error and reaches a much lower error level before eventually stagnating. Better error behaviour is seen for the error metric  $\|b(s^{i+1})-b(s^i)\|_{L^2(Q)}$ due to convergence \eqref{eq:L_conv_2} predicted in \Cref{theo:converegence}.

Nevertheless, the main strength of the L-scheme comes from the fact that it reduces the nonlinear degenerate problem to solving a set of heat equations. This allows us to combine it with adaptive mesh refinement with precise a posteriori error control obtained through the error-residual relationship derived below.


%\begin{figure}[htbp]
%\centering

% ===== Top Row: Gradient of L2 norm of error =====
%\begin{subfigure}[t]{0.24\textwidth}
%  \centering
%  \includegraphics[width=\linewidth]%{Figures/iter_doubledeg_UH1.eps}
%  \caption{Double-degenerate}
%\end{subfigure}\hfill
%\begin{subfigure}[t]{0.24\textwidth}
%  \centering
%  \includegraphics[width=\linewidth]{Figures/itersingledegBBu2_bbu_H1.eps}
% \caption{\centering Single-degenerate \\ $B(u)=u^2,\; b(u)=u$}

%\end{subfigure}\hfill
%\begin{subfigure}[t]{0.24\textwidth}
%  \centering
%  \includegraphics[width=\linewidth]{Figures/iter_singledeg_Bu2_bu.eps}
% \caption{\centering Single-degenerate \\ $B(u)=u,\; b(u)=u^2$}
% \end{subfigure}\hfill
%\begin{subfigure}[t]{0.24\textwidth}
%  \centering
%  \includegraphics[width=\linewidth]{Figures/itererr-vs-iter_non-degenerate.eps}
%  \caption{\centering Non-degenerate
%  $B(u)=u^2+0.1u, b(u)=u$}
%\end{subfigure}
  %=======================================
 % \begin{subfigure}[t]{0.42\textwidth}
 %   \centering
 %   \includegraphics[width=\linewidth]{Figures/itererr-vs-iter_double-deg.eps}
  %  \caption{Double-degenerate}
   % \end{subfigure}\hfill
   
    % \begin{subfigure}[t]{0.32\textwidth}
    %\centering
    %\includegraphics[width=\linewidth]{Figures/iter_doubledeg_UexactH1.eps}
    %\caption{Double-degenerate}
    %\end{subfigure}\hfill
   %======================================= 
 % \vspace{1em} 

  % ---- Bottom Row:
  %L2 norm of error ----
%  \begin{subfigure}[t]{0.24\textwidth}
%    \centering
%    \includegraphics[width=\linewidth]{Figures/iter_doubledeg_UL2..eps}
%    \caption{Double-degenerate}
%  \end{subfigure}\hfill
%  \begin{subfigure}[t]{0.24\textwidth}
%    \centering
 %   \includegraphics[width=\linewidth]{Figures/itersingledegBBu2_bbu.eps}
% \caption{\centering Single-degenerate \\ $B(u)=u^2,\; b(u)=u$}
%  \end{subfigure}\hfill
%  \begin{subfigure}[t]{0.24\textwidth}
%    \centering
 %   \includegraphics[width=\linewidth]%{Figures/iter_singledeg_bbuu2_Bu_L2.eps.eps}
% \caption{\centering Single-degenerate \\ $B(u)=u,\; b(u)=u^2$}
% \end{subfigure}\hfill
  %\begin{subfigure}[t]{0.32\textwidth}
   % \centering
   % \includegraphics[width=\linewidth]{Figures/iter_err-vs-iter_double-deg.eps}
   % \caption{Double-degenerate}
    %\end{subfigure}\hfill
    % \begin{subfigure}[t]{0.32\textwidth}
   % \centering
    %\includegraphics[width=\linewidth]{Figures/iter_doubledeg_UexactL2.eps}
    %\caption{Double-degenerate}
    %\end{subfigure}\hfill
%  \begin{subfigure}[t]{0.24\textwidth}
%    \centering
%    \includegraphics[width=\linewidth]{Figures/iter_err-vs-iter_non-degenerate.eps.eps}
%  \caption{\centering Non-degenerate
%  $B(u)=u^2+0.1u, b(u)=u$}  \end{subfigure}
  
%\KM{Make the error from $\|s^i_h -s^{i-1}_h\|$ to $\|u^i_h -u^{i-1}_h\|$ for the double-degenerate cases.}  
% \caption{
%Decay of the linearization error for the three model classes. 
%\textbf{Top row:} error measured in the $H^1(Q)$ norm. 
%\textbf{Bottom row:} error measured in the $L^2(Q)$ norm. 
%From left to right, the columns correspond to the double-degenerate model,  
%the single-degenerate model with $B(u)=u^2,\; b(u)=u$,  
%the single-degenerate model with $B(u)=u,\; b(u)=u^2$,  
%and the non-degenerate model with $B(u)=u^2+0.1u,\; b(u)=u$.}
%
%  \label{fig:iter-combined}
%\end{figure}

\section{Decomposition of error into discretization and linearization components}\label{subsec:decomposition}
The main idea in the a posteriori analysis is in the decomposition of the error into two components, linearization and discretization. To do so, with $s \in \cal S$ and $s^{i+1} \in \mathcal{V}$ solving Problems~\ref{eq:reformulated_weak_form} and \ref{pb:2}, respectively, we define $\calE_{\mathrm{lin}}$ and $\calE_{\mathrm{disc}}^{\rm H}$. 

Let $\calR$  and $\calR_{\rm lin}$ be the residuals associated with Problems~\ref{eq:reformulated_weak_form} and \ref{pb:2}, as defined in \eqref{eq:residual} and \eqref{eq:defRlin}, respectively. We define the total error of $s^i$ as a candidate solution to Problem \ref{eq:reformulated_weak_form} as
\begin{align}\label{eq:total_indicator}
\mathcal{E}(s^i) 
:=\left(\left\vert\!\left\vert \mathcal{R}(s^i) \right\vert\!\right\vert\!_{L^2(0,T;H^{-1}(\Omega))} ^2 + L^2 \|s^{i}(0)-s_0\|^2\right)^{\frac{1}{2}}.
\end{align}
Now, for a given $s^i$, the next iteration $s^{i+1}$ is a solution to Problem \ref{pb:2}. It is approximated by a fully discrete one $s^{i+1}_{h\tau}$. The corresponding  discretization error is defined as
%The discretization error corresponding to the approximation of $s^{i+1}$ by $s^{i+1}_{h\t}$ as a solution to the heat equation is defined as
\begin{align}\label{eq:disc_error}
\calE_{\mathrm{disc}}^{\rm H}(s^{i+1}_{h\t})
:= \left(\left\vert\!\left\vert \mathcal{R}_{\mathrm{lin}}(s_{h\tau}^{i+1}, s^i) \right\vert\!\right\vert\!_{L^2(0,T;H^{-1}(\Omega))} ^2 + L^2 \|s^{i+1}_{h\t}(0)-s_0\|^2\right)^{\frac{1}{2}}.
\end{align}
Note that, $\calE_{\mathrm{disc}}^{\rm H}$  vanishes if and only if $s^{i+1}_{h\t}$ solves Problem \ref{pb:2}. It can be seen as the discretization error since $s^{i+1}_{h\t}$ is typically a finite dimensional approximation of $s^{i+1}$. This association will become clearer in Proposition~\ref{lemma:SP_Iterative}. Furthermore, the linearization error is defined as 
\begin{align}\label{err:linearization}
\calE_{\mathrm{lin}}^i 
:= \|s^{i+1}_{h\t}- s^i\|_\mathcal{V}.
\end{align}
It inculcates the linearization component, vanishing if the iterations converge (even for a fixed discretization). These errors are related as follows from the following result.

\begin{theorem}[Decomposition of Error]\label{theo:decomp}
Let $L\in (1,2)$. Given $s^i\in (\calS \cap \calV)$, let $s^{i+1}_{h\t}\in \calV$ denote an approximation of $s^{i+1}\in \calV$, the solution to Problem \ref{pb:2}. With the linearization and discretization errors, $\mathcal{E}^i_{\mathrm{lin}}$ and $\mathcal{E}^{\mathrm{H}}_{\mathrm{disc}}$ defined in \eqref{err:linearization} and \eqref{eq:disc_error}, respectively, the total error defined in \eqref{eq:total_indicator} satisfies the reliability and efficiency bounds:
\begin{subequations}
    \begin{align}
\mathcal{E}(s^i)
&\leq \calE_{\mathrm{disc}}^{\mathrm{H}}(s^{i+1}_{h\t}) + L\calE_{\mathrm{lin}}^i
\quad &&\textup{(Reliability)}, \label{reliability_bound}\\
\calE_{\mathrm{disc}}^{\mathrm{H}}(s^{i+1}_{h\t})
&\leq \mathcal{E}(s^i) + L\calE_{\mathrm{lin}}^i
\quad &&\textup{(Efficiency)}.\label{eq:efficiency}
\end{align}
\end{subequations}
\end{theorem}


%\begin{theorem}[Decomposition of %Error]\label{theo:decomp}
%For a given $s^i\in (\calS \cap \calV)$, let %$s^{i+1}_{h\t}\in \calV$ denote an arbitrary %approximation of $s^{i+1}\in \calV$ to solve Problem %\ref{pb:2} for some $L\in (1,2)$.  Let the residuals %corresponding to Problems %\ref{eq:reformulated_weak_form} and \ref{pb:2} be %$\calR$  and $\calR_{\rm lin}$ defined in %\eqref{eq:residual} and \eqref{eq:defRlin}, respectively.
%Define the total error of $s^i$ as a candidate %solution to Problem \ref{eq:reformulated_weak_form} as
%\begin{align}\label{eq:total_indicator}
%\mathcal{E}(s^i) 
%:=\left(\int_0^T\left\vert\!\left\vert R(s^i) %\right\vert\!\right\vert\!_{H^{-1}(\Omega)}^2 + L^2 %\|s^{i}(0)-s_0\|^2\right)^{\frac{1}{2}}.
%\end{align}
%The discretization error corresponding to the %approximation of $s^{i+1}$ by $s^{i+1}_{h\t}$ as a %solution to the heat equation is defined as
%\begin{align}\label{eq:disc_error}
%\calE_{\mathrm{disc}}^{\rm H}(s^{i+1}_{h\t})
%:= \left(\int_0^T\left\vert\!\left\vert %R_{\mathrm{lin}}(s_{h\tau}^{i+1}, s^i) %\right\vert\!\right\vert\!_{H^{-1}(\Omega))}^2 + L^2 %\|s^{i+1}_{h\t}(0)-s_0\|^2\right)^{\frac{1}{2}}.
%\end{align}
%The linearization error is defined as the difference %between iterations in $\calV$-norm
%\begin{align}\label{err:linearization}
%\calE_{\mathrm{lin}}^i 
%:= \|s^{i+1}_{h\t}- s^i\|_\mathcal{V}.
%\end{align}
%Then, the total error satisfies the reliability and %efficiency bounds:
%\begin{subequations}
%    \begin{align}
%\mathcal{E}(s^i)
%&\leq \calE_{\mathrm{disc}}^{\mathrm{H}}%(s^{i+1}_{h\t}) + L\calE_{\mathrm{lin}}^i
%\quad &&\textup{(Reliability)}, %\label{reliability_bound}\\
%\calE_{\mathrm{disc}}^{\mathrm{H}}(s^{i+1}_{h\t})
%&\leq \mathcal{E}(s^i) + L\calE_{\mathrm{lin}}^i
%\quad &&\textup{(Efficiency)}.\label{eq:efficiency}
%\end{align}
%\end{subequations}
%\end{theorem}

 

\textit{What makes \Cref{theo:converegence} quite interesting is as follows}: Since $s^i$ and $s^{i+1}_{h\t}$ are known (unlike the infinite dimensional object $s^{i+1}$), $\calE_{\mathrm{lin}}^i$ is easily estimable using \eqref{err:linearization}. In the next section, we will call this estimation $\eta_{\mathrm{lin}}^i$. On the other hand,  $\calE_{\mathrm{disc}}^{\rm H}$ corresponds to estimating the residual of a heat equation. This makes it possible to provide highly accurate, locally-efficient estimates that are robust with respect to the mesh-size, time-step size, and the spatial and temporal polynomial degree; see \cite{ErnSmearsVohralik2017}. We will call this estimator $\eta_{\mathrm{disc}}^{\rm H}$. The details regarding this estimator will be given below.

Now, as the iterations converge, upon meeting an adaptive stopping criteria (i.e., when $\eta_{\mathrm{lin}}^i\ll \eta_{\mathrm{disc}}^{\rm H}$), we can provide \emph{guaranteed, efficient a posteriori error bounds for the double degenerate problem that are not only robust with respect to discretization but also with respect to the nonlinearities}. This is the underlying idea of the adaptive (space-time) linear iterative scheme proposed in this work.

\subsection{Proof of \Cref{theo:decomp}}
The proof relies on the results stated below, which can be obtained straightforwardly from \Cref{lemma_2.2} and \eqref{eq:defRlin}-\eqref{eq:l_b}. First, given $s^i\in (\calS \cap \calV)$, if $s^{i+1}$ solves Problem \ref{pb:2} one has

\[
\calR(s^i) =-L\resH(s^{i+1}-s^i),
\]
and
\[
\calR_{\mathrm{lin}}(s^{i+1}_{h\t},s^i)= L\resH(s^{i+1}_{h\t}-s^i) + \calR(s^i)= L\resH(s^{i+1}_{h\t}- s^{i+1}).
\]
The results below involve the errors introduced in   \eqref{eq:total_indicator}, \eqref{eq:disc_error} and \eqref{err:linearization}.
%Since $(s^{i+1}-s^i)(0)=s_0-s^i(0)$, using \Cref{lemma:resH} one immediately attains the following lemma
%This immediately gives using \Cref{lemma:resH} noting that $(s^{i+1}-s^i)(0)=s_0-s^i(0)$,
%\begin{proposition}[Equivalent expression of the dual norm of the residual]\label{lemma:SP_Iterative}
%One has the following results: Equivalent, 
%\begin{align}
%[\mathcal{E}(s^i)]^2 =L^2 \|s^{i+1}-s^i\|_{\calV}^2
%\end{align}
%\end{proposition}

%\begin{proposition}[Equivalent expression of the dual norm of the linearized residual]\label{}
%Equivalent,
%\begin{align}
%[\calE_{\mathrm{disc}}^{\rm H}(s^{i+1}_{h\t})]^2= L^2  \|s^{i+1} - s^{i+1}_{h\t}\|_{\calV}^2.
%\end{align}
%\end{proposition}
%These equalities are a direct consequence of \Cref{lemma_2.2}, and of the initial conditions $s^{i+1}(0)=s^i(0)=s_0$.

\begin{proposition}[Equivalent expressions of the dual norms of the residual and linearized residual]\label{lemma:SP_Iterative}
One has the following 
\begin{align}
[\mathcal{E}(s^i)]^2 &=L^2 \|s^{i+1}-s^i\|_{\calV}^2, \\
[\calE_{\mathrm{disc}}^{\rm H}(s^{i+1}_{h\t})]^2 &= L^2  \|s^{i+1} - s^{i+1}_{h\t}\|_{\calV}^2.
\end{align}
\end{proposition}
These equalities are a direct consequence of \Cref{lemma_2.2}, and of the initial conditions $s^{i+1}(0)=s^i(0)=s_0$.
With the proposition above, the proof of Theorem~\ref{theo:decomp} reduces to a straightforward application of the triangle inequality. First, by Proposition~\ref{lemma:SP_Iterative},
\[
\mathcal{E}(s^i)=L \|s^{i+1}-s^i\|_{\calV}\leq L \|s^{i+1} - s^{i+1}_{h\t}\|_{\calV} + L \|s^{i+1}_{h\t} - s^{i}\|_{\calV}=  \calE_{\mathrm{disc}}^{\rm H}(s^{i+1}_{h\t}) + L\calE_{\mathrm{lin}}^i.
\]
This is nothing but \eqref{reliability_bound}. The proof of \eqref{eq:efficiency} is similar, starting from $\calE_{\mathrm{disc}}^{\rm H}(s^{i+1}_{h\t})
= L\|s^{i+1}-s^{i+1}_{h\t}\|_{\calV}$. The triangle inequality applied to $L\|s^{i+1}_{h\t}-s^{i}\|_{\calV}$ reveals the other direction.\qed \\[2pt]

We use now the result in \Cref{subsec:decomposition} for establishing a posteriori estimates for a space time discretization of Problem~\ref{pb:2} and ultimately, to develop an adaptive, linear iterative scheme for Problem~\ref{eq:reformulated_weak_form}..

\section{Space-time Finite Elements, A Posteriori Analysis, and Adaptivity}\label{sec:apost}
For the finite element discretization of Problem~\ref{pb:2}, we consider a regular shape partition of $Q$ into  $\mathcal{T}_{h\tau}$, a family of simplicial space-time elements of ( $d+1$)-dimensions, that are uniformly shape-regular. We denote by $(\boldsymbol{x},t) \in \mathbb{R}^{d+1}$ a generic space-time point, where $\boldsymbol{x} \in \Omega$ is the spatial component and $t\in(0,T]$ the temporal component.  Let $\mathcal{N}_{h\tau}$ denote the set of vertices of the triangulation $\mathcal{T}_{h\tau}$, and $\mathfrak{F}_\htt$ the set of faces. For  $a\in \mathcal{N}_{h\t}$, let $\calT_\htt^a$ denote the set of elements that have $a$ as a vertex, $\phi_a$ denote the corresponding hat function (P1-Lagrange element basis), and $\oma$ the space-time patch, i.e., support of $\phi_a$. Let $h_K$ be the diameter of element $K$, and $h_\t:=\max_{K\in \calT_{h\t}} h_K>0$ denote the mesh-size of $\calT_\htt$. We define the broken polynomial space of degree at most $p \ge 0$ by $\mathcal P_p(\mathcal T_{h\tau}):=\{ v \in L^2(Q)|\; v|_K \in \mathcal P_p(K), \; \text{for all}\; K \in \mathcal T_{h\tau} \},$ where $\mathcal P_p(K)$ denotes the space of polynomials of total degree at most $p$ defined on $K$. Then, the conforming space-time finite element space is defined as
\begin{align}\label{eq:def_Vht}
    \calV_{h\t}:= \mathcal{P}_1(\calT_{h\t}) \cap H^1(Q)\subset \calS\cap \calV. 
\end{align}
The elementwise $L^2$-orthogonal projection $\Pi^p_{h\tau} : L^2(Q) \to \mathcal{P}_p(\mathcal{T}_{h\tau})$ is defined by
\begin{equation}\label{eq:L2_projection}
\int_Q (\Pi^p_{h\tau} v)\, q 
=
\int_Q v\, q 
\qquad \text{for all}\quad q \in \mathcal P_p(\calT_{h\t}).
\end{equation}
and $\bm{\Pi}^p_{h\t}: L^2(Q)^d\to \calP_p(\calT_{h\t})^d$ be the component-wise projection for vector-valued functions.
\subsection{Space-time finite elements}
%%%%%%%%%%%%%%%%%%%%%%%%%%%%%%%%%%%%%%%%%%%%%%%%%%%%%%%%%%%%%%
%\mathcal T_{h\tau}  % mesh
%\mathcal V_{h\tau}  % FE space
%\mathcal P_p        % polynomials
%\mathcal N_{h\tau}  % vertices
%\tilde{\omega}_a    % patch
%K^a                 % subelements
%%%%%%%%%%%%%%%%%%%%%%%%%%%%%%%%%%%%%%%%%%%%%%%%%%%%%%%%%%%%%%

\begin{theorem}[Space-time finite elements]\label{theo:lsp}
Let $s^0\in (\calS\cap \calV)$ be given and $\{s^{i}\}_{i\in\N}\subset \calV$ the sequence of solution to the $L$-scheme, obtained by solving Problem \ref{pb:2}.   Then, there exists a unique sequence of finite element solutions $\{s_{h\t}^{i}\}_{i\in \N}\subset \mathcal V_{h\tau}\subset (\calS\cap \calV)$ that satisfies $s^i_{h\t}(0)= \Pi^1_{h\t} s_0$, $s^{0}_{h\t}= \Pi^1_{h\t} s^0$, and
  \begin{align*}
      \int_0^T \!\left\langle \partial_t \left( L(s^{i+1}_{h\t} - s^i_{h\t}) + b(s^i_{h\t}) \right), \varphi_{h\t} \right\rangle \,  
+ \int_0^T \!\left( \nabla \left( L(s^{i+1}_{h\t} - s^i_{h\t}) + B(s^i_{h\t}) \right), \nabla \varphi_{h\t} \right) 
= \int_0^T \!( f, \varphi_{h\t} )
  \end{align*}
%\textcolor{red}
{for all $\f_{h\t}\in \mathcal V_{h\tau}$, or equivalently, in terms of the linear functional defined in \eqref{eq:defRlin},
\begin{align}\label{eq:lsp}
\calR_{\rm lin}(s^{i+1}_{h\t},s^i_{h\t}) = 0 .
\end{align}}
Further, let $i\in \N$ be fixed and $\bar{s}^{i+1}\in \calV$ solve $\calR_{\rm lin}(\bar{s}^{i+1},s^i_{h\t})=0$ with $\bar{s}^{i+1}(0)=s_0$. Then, 
\begin{equation}\label{eq:apriori_est_lsp}
    \|s^{i+1}_{h\t}-\bar{s}^{i+1}\|_{L^2(0,T;H^1_0(\Om))}\lesssim \inf_{\f_{h\t}\in \calV_{h\t}}\|\bar{s}^{i+1}-\f_{h\t}\|_\calV \to 0 \qquad \text{ as } \qquad h_\t \to 0.
\end{equation}
% \begin{align}
%     L\|s^{i+1}_{h\t}-s^{i+1}\|_{L^2(0,T;H^1_0(\Om))}&\leq L \|\bar{s}^{i+1}-s^{i+1}_{h\t}\|_{L^2(0,T;H^1_0(\Om))} + \left\|L(s^i-s^{i}_{h\t})-(b(s^i)-b(s^{i}_{h\t}))\right\|_{H^1(0,T;H^{-1}(\Omega))}\nonumber\\
%     & + \left\|L(s^i-s^{i}_{h\t})-(B(s^i)-B(s^{i}_{h\t}))\right\|_{L^2(0,T;H^{1}_0(\Omega))}. \label{eq:apriori_est_lsp}
% \end{align}
\end{theorem}
%\begin{remark}[A convergence result %\eqref{eq:apriori_est_lsp}]
%Since we do not expect any higher regularity of $\bar{s}^{i+1}\in \calV$, particularly in the degenerate cases, an explicit rate of convergence with respect to $h_\t$ cannot be provided instead of \eqref{eq:apriori_est_lsp}.  As such, \eqref{eq:apriori_est_lsp} shows consistency but is not practically useful. The focus for the discretization part of this work will be on a posteriori estimates and an adaptive algorithm.
%\end{remark}

\begin{proof}
From \eqref{eq:defRlin} and \eqref{eq:lsp} we get $ L\resH(s^{i+1}_{h\t})(\f_{h\t})= L\resH(s^{i}_{h\t})(\f_{h\t}) -\calR(s^{i}_{h\t})(\f_{h\t}) $ for all $\f_{h\t}\in \calV_{h\t}$. Observe that $s^{i}_{h\t}\in \calV_{h\t}\subset (\calS\cap \calV)$ implies that $L\resH(s^{i}_{h\t}) -\calR(s^{i}_{h\t})\in L^2(0,T;H^{-1}(\Om))$. Thus, from Theorem 3.2 of \cite{steinbach2015space}, given $s_{h\tau}^i\in \calV_{h\t}$, there exists a unique $s^{i+1}_{h\t}\in \calV_{h\t}\subset (\calS\cap \calV)$, such that \eqref{eq:lsp} holds true, and one proceeds with the inductive logic. The estimate \eqref{eq:apriori_est_lsp} follows from Theorem 3.2 of \cite{steinbach2015space}.
\end{proof}
\begin{remark}
Particularly in the degenerate cases, we do not expect any higher regularity of $\bar{s}^{i+1}\in \calV$. Therefore, an explicit order of convergence with respect to $h_\t$ cannot be provided. As such, \eqref{eq:apriori_est_lsp} shows consistency, but is less useful for the main goal of this work, namely the a posteriori analysis and the development of an adaptive scheme.

\end{remark}
% \begin{proof}
% Expanding \eqref{eq:lsp} we get $\int_0^T \langle L\resH(s^{i+1}_{h\t}),\f_{h\t}\rangle =\int_0^T \langle L\resH(s^{i}_{h\t}) -\calR(s^{i}_{h\t}),\f_{h\t}\rangle $ for all $\f_{h\t}\in \calV_{h\t}$.
%     Observe that $s^{i}_{h\t}\in \calV_{h\t}\subset \calS$ implies that $L\resH(s^{i}_{h\t}) -\calR(s^{i}_{h\t})\in L^2(0,T;H^{-1}(\Om))$. Thus, from Theorem 3.2 of \cite{steinbach2015space}, there exists a unique solution to the problem \eqref{eq:lsp}. Consequently, $s^{i+1}_{h\t}\in \calV_{h\t}\subset \calS$, and one proceeds with the inductive logic.

%     Now, observe from the definition of $\bar{s}^{i+1}\in \calV$ that $\int_0^T \langle \resH(s^{i+1}_{h\t}-\bar{s}^{i+1}),\f_{h\t}\rangle =0$ for all $\f_{h\t}\in \calV_{h\t}$. Using the triangle inequality and Theorem 3.2 of \cite{steinbach2015space}, one further obtains
% \begin{align}
%         &L\|s^{i+1}_{h\t}-s^{i+1}\|_{L^2(0,T;H^1_0(\Om))}\leq         L\|s^{i+1}_{h\t}-\bar{s}^{i+1}\|_{L^2(0,T;H^1_0(\Om))} + L\|s^{i+1}-\bar{s}^{i+1}\|_{L^2(0,T;H^1_0(\Om))}\nonumber\\
%         &\lesssim \inf_{\f\in \calV_{h\t}} \|\bar{s}^{i+1}-\f_{h\t}\|_\calV \overset{\Cref{lemma:resH}}+ L\|  \resH(s^{i+1}-\bar{s}^{i+1})\|_{L^2(0,T;H^{-1}(\Om))}.\label{eq:TrianglSht}
% \end{align}
% For the last term, using the definition of $s^{i+1}$ and $\bar{s}^{i+1}$, we write
% \begin{align*}
%     &L\int_0^T \langle \resH(s^{i+1}-\bar{s}^{i+1}),\f \rangle  = \int_0^T \langle L\resH(s^{i}-s^{i}_{h\t}),\f \rangle -  \int_0^T \langle \calR(s^{i})-\calR(s^{i}_{h\t}),\f \rangle\\
%     &=\int_0^T \langle \p_t\left(L(s^i-s^{i}_{h\t})-(b(s^i)-b(s^{i}_{h\t})\right),\f \rangle +  \int_0^T \left (\nabla\left( L(s^i-s^{i}_{h\t})-(B(s^i)-B(s^{i}_{h\t}))\right),\nabla \f\right).
% \end{align*}
%     Using the definition of the $L^2(0,T;H^{-1}(\Om))$ norm in \eqref{eq:TrianglSht}, we get the result.
% \end{proof}

\subsection{A posteriori estimates}\label{sec:a-posteriori}
There are several works developing reliable and efficient estimators for nonlinear diffusion equations. We refer to \cite{dolejvsi2024error,congreve2026residual,DolejsiRoskovec2021} for discontinuous Galerkin space-time methods (on either mesh-dependent norms, or efficiency dependent on the nonlinearities), and \cite{ErnSmearsVohralik2017,ErnSmearsVohralik2019} for conforming finite elements. In these works, the local adaptation on time is not addressed. However, for problems with local, evolving features such as free boundaries, an efficient discretization scheme should employ space-time adaptivity. Here, we introduce an \emph{easy-to-implement, fully and parallelly computable a posteriori estimator that works in combination with linear iterations, specifically the L-scheme or the regularized Newton and is robust with respect to nonlinearities/degeneracies}. We will also show the reliability of this estimator, which will mainly be used in the adaptive algorithm. For this purpose, an equilibrated flux $\bm{\sigma}_h$ needs to be constructed, which is mass-conservative. 

\subsubsection{Space of fluxes}
Concerning fluxes, we start by noting that, given $\nhat\in \R^d$ and $e\in\mathbb{R}$, then $t+\nhat\cdot \bm{x}=e$ is the equation of a hyperplane inside $Q'$. For any $\bm{v}\in H(\mathrm{div};Q')$, $\bm{v}\cdot \nhat$ remains conforming accross $\mathcal{H}$. However, it does not include any time component, and thus, does not require conformity in the time-direction.

\begin{definition}[Raviart-Thomas spaces for space-time problems]\label{def:Qhtspace}
For a triangulation $\calT_\htt$, its vertex $a\in \calN_\htt$ with the corresponding space-time patch $\omega_a$, and a polynomial degree $p\geq 2$, we define the spaces
\[
\calQ_\htt:= H(\mathrm{div},Q)\cap \calP_p(\calT_\htt),\qquad  \calQ^a_\htt:=\{\boldsymbol{v}^a_h \in H(\mathrm{div},\oma)\cap \calP_p(\calT_\htt^a) \;|\; \bm{v}^a_h\cdot \nhat|_{\p\oma \setminus \Gamma_T}=0 \}.
\]
In the above, $\Gamma_T:=\p \Om \times [0,T]$ and $\nhat$ denotes the restriction to an outward normal of $\p\oma$ in the space direction. As above, the equation of a face $E$ that is part of $\p\oma$ is represented by $t+ \nhat\cdot \x=e$, where $t$ is the time component of $a$ and $e\in \mathbb{R}$ suitably chosen.
\end{definition}

The non-emptyness of the spaces will follow from \Cref{mass_balance_property}.

% \begin{lemma}[Well-posedness of the spaces $\calQ_\htt$ and $\calQ_\htt^a$]
%     The spaces $\calQ_\htt$ and $\calQ_\htt^a$ of \Cref{def:Qhtspace} are not-empty Hilbert subspaces of $H(\mathrm{div},Q)$ and $H(\mathrm{div},\oma)$ respectively.
% \end{lemma}
% \begin{proof}
%     It is direct to see that they are Hilbert subspaces if they are non-empty. We show that such a space is well-posed with a degree of freedom associated with each face $E\in \mathfrak{F}_\htt$. For the $(d+1)$ dimensional simplical element $K\in \calT_\htt$, let $\{E_j\}_{1\leq j\leq d+2}$  denote the set of faces. Choose one face $E_i$ where $1\leq i\leq d+2$, and define the vector function $\bm{v}^K_i:= \theta_i \x + \bm{b} t + \bm{c}$ on $\bar{K}$. We impose that $\bm{v}^K_i\cdot (\nhat)_j=0$ for all $j\not=i$, i.e., normal space-flux through all other faces in $\bar{K}$ is zero. Using the relation $t+ (\nhat)_j\cdot \x=e_j$ for each of the faces, we find the set of $2(d+1)$ equations
%     \[
%     \bm{b}_i\cdot (\nhat)_j=\theta_i, \text{ and } \bm{c}_i\cdot (\nhat)+ \theta_ie_j=0 \quad \forall\, i\not=j.
%     \]
%     Assuming, $(\nhat)_j$ are linearly independent (which should be possible since there are $(d+1)$ of them not including $i$), $\bm{b}_i$ and $\bm{c}_i$ can be solved if $\theta_i$ is prescribed. Thus, each $\bm{v}^K_i$ has exactly one free parameter. This can be used to prescribe a total mass flow between the elements $K$ and $K'$ that share $E_i$. Linear combination of all such mass-balancing flux components $\bm{v}^K_i$ from all $E_i\in \mathfrak{F}_\htt$ gives us  $\calQ_\htt$. Similarly, the non-emptiness of $\calQ_\htt^a$ is shown.
% \end{proof}


\subsubsection{Equilibrated flux}\label{sec:equilibrated_fluxes}
For the construction of the equilibrated flux, we rewrite \eqref{eq:lsp} as 
\begin{align}\label{eq:source_flx}
 \int_0^T ( \Shti, \varphi_{h\tau} ) \, dt
 + \int_0^T ( \Fhti, \nabla \varphi_{h\tau} ) \, dt = 0,
 \quad \;\text{for all}\quad\; \varphi_{h\tau} \in \mathcal{V}_{h\tau},
\end{align}
where
\begin{align}\label{eq:source_flx_Lscheme}
    \Shti:= \partial_t \left( L(s^{i+1}_{h\t} - s^i_{h\t}) + b(s^i_{h\t}) \right), \qquad\;\text{and}\; \qquad \Fhti:= \nabla \left( L(s^{i+1}_{h\t} - s^i_{h\t}) + B(s^i_{h\t}) \right).
\end{align}
Then, we can define the equilibrated flux.
\begin{definition}[Equilibrated flux $\flx$]\label{def:equilibrated}
    Let \eqref{eq:source_flx} hold. For a vertex $a\in \calN_\htt$, let the local flux $\flxa\in \calQ_\htt^a$ be defined by
    \begin{align}\label{eq:optimization}
        \flxa:= \underset{\substack{\bm{v}^a_h\in\calQ_\htt^a\\
       \nabla\cdot \bm{v}^a_h = -\Pi_h^0 (\phi_a \Shti) - \nabla\phi_a \cdot \bm{\Pi_h^0} \Fhti}}{\mathrm{argmin}} \left\| \bm{v}^a_h + \phi_a \Fhti\right\|_{L^2(\oma)}
    \end{align}
    Here $\phi_a$ is the $P_1$-lagrange element basis satisfying $\phi_a(a)=1$. After extending $\flxa$ by zero to $Q\setminus \oma$, the equilibrated flux $\flx\in\calQ_\htt$ is defined as
    \[
    \flx:=\sum_{a\in \calN_\htt} \flxa.
    \]
\end{definition}
\begin{proposition}[Well-posedness and mass balance property of $\flx$]\label{mass_balance_property}
    The equilibrated flux $\flx \in \calQ_\htt$ is well defined and satisfies 
    \[
    \nabla\cdot \flx= -\Pi_h^0(\Shti) \qquad \text{ in all } K\in \calT_\htt.
    \]
\end{proposition}
\begin{proof}
First, by counting degrees of freedom for a given $a\in \calN_\htt$ we show that there are indeed $\bm{v}_h^a\in \calQ_\htt^a $ such that $\nabla\cdot \bm{v}^a_h = -\Pi_h^0 (\phi_a \Shti) - \nabla\phi_a \cdot \bm{\Pi_h^0} \Fhti\in \calP_0(\calT_\htt)$. \textcolor{black}{For a given vertex $a \in \mathcal{N}_{h\tau}$ such that $a \notin \Gamma_T$, let $N_a \in \mathbb{N}$ denote the number of elements in $\omega_a$. This gives $N_a$ constraints for satisfying the divergence conditions. Observe that $\partial\omega_a$ contains $N_a$ faces. Since the elements in $\mathcal{T}_{h\tau}$ are $(d+1)$-simplices, each of the remaining $(d+1)N_a$ faces is shared by two adjacent elements, giving in total $((d+1)N_a/2)$ internal faces}. 
%Moreover, $\p\oma$ contains $N_a$ faces since every element has exactly one face not including $a$, and thus lying on the boundary. The remaining $(d+1)N_a$ faces of elements (since the elements are $(d+1)$-dimensional simplices),  are all shared between two elements, thus having $((d+1)N_a/2)$ internal faces. 
For $\bm{v}_h^a\in \calQ_\htt^a$, the trace $\bm{v}^a_h\cdot \nhat$ on each face $E$ gives a polynomial of degree $p$ with $d$ variables (since $t=e-\nhat\cdot \x$ can be substituted as the equation of the face). This implies at most $C_p^{p+d}$ unknown coefficients per  face, thus $((d+1)/2 +1) N_a C_p^{p+d}$ coefficients in total. Now, the total degrees of freedom of $\bm{v}_h^a$ (which is $d$-dimensional) per element are $dC^{p+d+1}_p$. The total number of degrees of freedom needs to exceed the number demanded by constraints, which yields the relation
\[
dC^{p+d+1}_p N_a \geq \left[\left(\frac{d+1}{2} +1\right) C_p^{p+d} + 1\right]N_a,
\]
or, after cancelling terms,
\[
\frac{d(p+d+1)}{d+1}\geq \frac{d+3}{2} + \frac{1}{C_p^{p+d}}.
\]
Rearranging, we get the condition \[p\geq \frac{1}{2}(3-d)(d+1) + \frac{d+1}{d C_p^{p+d}}\]
which is satisfied for all $p\geq 2$ if $d\geq 2$. For $d=1$, this restriction is actually an overcount since $\bm{v}_h^a$ is scalar, and thus $\bm{v}^a_h\cdot \nhat=0$ implies $\bm{v}_h^a=0$. This condition has to be enforced in the boundary nodes of $\oma$, see \Cref{fig:P2nodes}. The value at the internal nodes, $(N_a+1)$ in number, can be freely chosen to satisfy $N_a$ divergence constraints (actually one less due to the compatibility condition). When $a\in \calN_\htt$ is on the boundary $\Gamma_T$, then the restrictions are even less, since no conditions are required on $\p\oma \cap \Gamma$. Thus, for $p\geq 2$ the space $\calQ_\htt^a$  has the required degrees of freedom to satisfy conditions for $\bm{v}_h^a$ in \eqref{eq:optimization}.

    \begin{figure}[h]
                \centering
                    \includegraphics[width=.4\linewidth]{Figures/P2-nodes-patch.png}
                    \caption{For one space dimension ($d=1$), the  $P2$ nodal points of $\oma$.}
        \label{fig:P2nodes}
    \end{figure}



Next, we verify the compatibility criteria for $\bm{v}^a_h\in \calQ^a_\htt$. Using the divergence theorem on the field $(\bm{v}^a_h,0)$, we have
\begin{align*}
        0&=\int_{\p\oma} \bm{v}^a_h\cdot \nhat= \int_{\oma}  \nabla \cdot \bm{v}^a_h = - \left[\int_{\oma} \Pi_h^0 (\phi_a \Shti) + \nabla\phi_a \cdot \bm{\Pi_h^0} \Fhti\right]\\
        &= -\left[\int_0^T (\Shti, \phi_a ) \, dt
 + \int_0^T ( \Fhti, \nabla \phi_a ) \, dt\right]=0.
\end{align*}
The last two steps follow from the property of the projector $\Pi^0_h$, and from \eqref{eq:source_flx} by inserting $\f_\htt=\phi_a$. The Euler-Lagrange equation corresponding to the constrained minimization problem \eqref{eq:optimization} then becomes: find $(\flxa,\l^a_\htt)\in \calQ_\htt^a \times \calP_1(\calT_a)$ such that 
\begin{equation*}
    \left\{\begin{aligned}
        &(\flxa + \phi_a \Fhti, \bm{v}^a_h)_{L^2(\oma)} + (\l^a_\htt, \nabla \cdot \bm{v}^a_h)_{L^2(\oma)}= 0,\\
        &(l_h^a, \nabla \cdot \flxa)_{L^2(\oma)}= -(l_h^a\phi_a, \Shti)_{L^2(\oma)}- (l_h^a \nabla \phi_a, \Fhti)_{L^2(\oma)},
    \end{aligned}\right.  
\end{equation*}
for all $(\bm{v}^a_h,l_h^a)\in \calQ_\htt^a \times \calP_1(\calT_a)$. The existence-uniqueness of a solution to this system follows from the standard inf-sup condition since it is a saddle-point problem \cite[Chapter 4]{boffi2013mixed}.

Finally, we get by the partition of unity property $\sum_{a\in \calN_{h\t}} \phi_a=1$ that
\begin{align*}
    &\nabla \cdot \flx = \sum_{a\in \calN_{h\t}} \nabla \cdot \flxa= -\sum_{a\in \calN_{h\t}} \left[\,\Pi_h^0 (\phi_a \Shti) + \nabla\phi_a \cdot \bm{\Pi_h^0} \Fhti\right]\\
    &= - \left[\Pi_h^0 \left((\Sigma_{a\in \calN_{h\t}} \phi_a) \Shti\right) + \nabla(\Sigma_{a\in \calN_\htt}\phi_a ) \cdot \bm{\Pi_h^0} \Fhti\right]= - \Pi_h^0 \Shti.
\end{align*}
%This concludes the proof.
\end{proof}

\subsubsection{Guaranteed reliability estimate}
The equilibrated fluxes constructed in \Cref{sec:equilibrated_fluxes} are used to provide local estimates for the error. For this, we use the residuals $\calR$  and $\calR_{\rm lin}$ defined in \eqref{eq:residual} and \eqref{eq:defRlin}, respectively. For $\calV_\htt$ defined in \eqref{eq:def_Vht}, let $\{s_{h\t}^{i}\}_{i\in \N}\subset \mathcal V_{h\tau}\subset (\calS\cap \calV)$ be the sequence of finite element solutions introduced in \Cref{theo:lsp}. With the equilibrated flux $\flx\in \calQ_\htt$ given in \Cref{def:equilibrated} with $\calQ_\htt$, $\Shti$, $\Fhti$ from \Cref{def:Qhtspace} and \eqref{eq:source_flx_Lscheme} respectively, we introduce the flux and data oscillation estimators for $Q'\subseteq Q$ as,

%Recall the definitions of total error $\calE$ \eqref{eq:total_indicator}, discretization error $\calE_{\mathrm{disc}}^{\rm H}$ \eqref{eq:disc_error}, and linearization error $\calE_{\mathrm{lin}}^{i}$ \eqref{err:linearization}.
\begin{subequations}
   \begin{align}\label{flux_oscillation_estimator}
        \eta_{\mathrm{F},Q'}^i:= \|\flx + \Fhti\|_{L^2(Q')}, \qquad 
    \eta_{\mathrm{osc},Q'}^i:= C_{\,\!_{\Omega}}^{\mathrm{P}} \,h_\Om\,\|(\mathbb{I}-\Pi_h^0) (\p_t b(s^i_\htt) -f)\|_{L^2(Q')} 
\end{align}
along with the discretization and linearization estimators
 \begin{align}
     &\eta_{\mathrm{disc}}^{i}:= \left((\eta_{\mathrm{F},Q}^i + \eta_{\mathrm{osc},Q}^i)^2 + L^2\|s^i_{h\t}(0)-s_0\|^2\right)^{\frac{1}{2}},\\
    &\eta_{\mathrm{lin}}^{i}:= \left([C_{\,\!_{\Omega}}^{\mathrm{P}}\,h_\Om]^2\,\|\p_t (s^{i+1}_\htt - s^i_\htt)\|_{L^2(Q)}^2 + \|\nabla (s^{i+1}_\htt - s^i_\htt)\|_{L^2(Q)}^2 + \|(s^{i+1}_\htt - s^i_{h\t})(T)\|^2\right)^{\frac{1}{2}}.
 \end{align}
 \end{subequations}
 Let $ \eta^i:= \eta_{\mathrm{disc}}^{i} + L\eta_{\mathrm{lin}}^{i}$ be the total estimator. Then we have the following.
\begin{theorem}[Guaranteed a posteriori upper bound on residual]\label{thm:apost-reliability-i}
The errors $\calE$, $\calE_{\mathrm{disc}}^{\rm H}$, and $\calE_{\mathrm{lin}}^{i}$ defined in \eqref{eq:total_indicator}--\eqref{err:linearization} satisfy the reliability estimates
 \begin{align}\label{eq:reliability_estimate}
\mathcal{E}_{\mathrm{disc}}^{\rm H}(s^{i+1}_{_\htt}) \leq \eta_{\mathrm{disc}}^{i}, \qquad \mathcal{E}^i_{\rm lin} \leq \eta^i_{\rm lin},\;\qquad \text{and} \;\qquad \mathcal{E}(s^i_{_\htt}) \leq \eta^i.
 \end{align}
\end{theorem}
We stipulate that proving the efficiency of this estimator (lower bound of $\calE(s^i_{h\t})$ by $\eta_{\mathrm{disc}}^{i}$) might also be possible with minor modifications, see \Cref{remark:efficiency}. However, since our main target is adaptivity, we skip this for brevity.
\begin{proof}
Let $\varphi \in L^2(0,T;H^1_0(\Omega))$ with $\|\nabla\varphi\|_{L^2(Q)} = 1$.
We start from the residual functional \eqref{eq:source_flx},
\[
 \mathcal{R}_{\rm lin}(s^{i+1}_{h\tau}, s^i_{h\tau})(\varphi)
= \int_0^T\!\Big[(\mathcal{S}^i_{h\tau},\varphi) + (\mathbf{F}^i_{h\tau},\nabla\varphi)\Big]\,dt.
\]
The following term will be added to the above. Note that this term vanishes. More precisely, with $\hat{\underline{\bm{n}}}_{t,x}$ being the unit normal on $\p Q$, by the divergence theorem we have
%($\hat{\underline{\bm{n}}}_{t,x}\in \mathbb{S}^{d+1}$ is the unit-normal on $\p Q$)
\begin{align*}
    &\int_{Q} \left[\f\nabla \cdot \flx + \nabla \f\,\cdot \flx\right]= \int_{Q}\nabla \cdot (\f\flx) = \int_{\p Q} \f (\flx,0)\cdot \hat{\underline{\bm{n}}}_{t,x}\\
    &= \int_{\p \Om\times [0,T]} \f \,\flx\cdot \hat{\bm{n}}_{x} +  \int_{\Om\times \{0\}} \f \,\flx\cdot (\bm{0},-1) +  \int_{\Om\times \{T\}} \f\, \flx\cdot (\bm{0},1)=0.
\end{align*}
The first term on the right vanishes since  $\varphi\in L^2(0,T;H^1_0(\Om))$. Adding this term yields
\begin{align*}
 \mathcal{R}_{\rm lin}(s^{i+1}_{h\tau}, s^i_{h\tau})(\varphi)
&= \int_0^T\!\Big[(\mathcal{S}^i_{h\tau} + \nabla\cdot\boldsymbol{\sigma}^i_{h\tau},\varphi)
+ (\mathbf{F}^i_{h\tau} + \boldsymbol{\sigma}^i_{h\tau},\nabla\varphi)\Big]\,\dd t.
\end{align*}
By \Cref{mass_balance_property}, $\nabla\cdot\boldsymbol{\sigma}^i_{h\tau} = -\Pi^0_h(\mathcal{S}^i_{h\tau})$ in all $K\in\mathcal{T}_{h\tau}$, so that
\[
\mathcal{S}^i_{h\tau} + \nabla\cdot\boldsymbol{\sigma}^i_{h\tau}
= (\mathbb{I}-\Pi^0_h)\mathcal{S}^i_{h\tau}.
\]
Recalling \eqref{eq:source_flx_Lscheme}, we expand
\[
(\mathbb{I}-\Pi^0_h)\mathcal{S}^i_{h\tau}
= (\mathbb{I}-\Pi^0_h)\partial_t\!\left(L(s^{i+1}_{h\tau}-s^i_{h\tau})\right)
+ (\mathbb{I}-\Pi^0_h)\left(\partial_t b(s^i_{h\tau})-f\right).
\]
Since $s^{i+1}_{h\tau}, s^i_{h\tau}\in\mathcal{V}_{h\tau}\subset\mathcal{P}_1(\mathcal{T}_{h\tau})$, the term $\partial_t(L(s^{i+1}_{h\tau}-s^i_{h\tau}))$ is piecewise constant on $\mathcal{T}_{h\tau}$ and thus, $(\mathbb{I}-\Pi^0_h)\partial_t\!\left(L(s^{i+1}_{h\tau}-s^i_{h\tau})\right) = 0$.
Substituting into the residual identity, we obtain
\begin{align}\label{eq:res-identity}
\mathcal{R}_{\rm lin}(s^{i+1}_{h\tau}, s^i_{h\tau})(\varphi)
&= \int_0^T\!\Big[((\mathbb{I}-\Pi^0_h)(\partial_t b(s^i_{h\tau})-f),\varphi)
+ (\mathbf{F}^i_{h\tau} + \boldsymbol{\sigma}^i_{h\tau},\nabla\varphi)\Big]\,dt.
\end{align}
Applying Cauchy-Schwarz to each term of \eqref{eq:res-identity},
\begin{align*}
\left| \mathcal{R}_{\rm lin}(\varphi)\right|
&\leq \|(\mathbb{I}-\Pi^0_h)(\partial_t b(s^i_{h\tau})-f)\|_{L^2(Q)}\,
\|\varphi\|_{L^2(Q)}+\|\mathbf{F}^i_{h\tau}+\boldsymbol{\sigma}^i_{h\tau}\|_{L^2(Q)}\,
\|\nabla\varphi\|_{L^2(Q)}.
\end{align*}
Since $\varphi(\cdot,t)\in H^1_0(\Omega)$ a.e.\ $t\in [0,T]$ and $\|\nabla \f\|_{L^2(Q)}=1$, using the Poincar\'e-Friedrichs inequality gives 
\begin{align}
    \|\varphi(t)\|_{L^2(Q)}\leq C_{\,\!_{\Omega}}^{\mathrm{P}}\,h_\Omega\, \|\nabla\varphi(t)\|_{L^2(Q)}=C_{\,\!_{\Omega}}^{\mathrm{P}}\, h_\Omega.
\end{align}
Therefore
\begin{align}\label{eq:Hminus1-R}
&\|\mathcal{R}_{\rm lin}(s^{i+1}_{h\tau},s^i_{h\tau})\|_{L^2(0,T;H^{-1}(\Omega))} :=\sup_{\|\nabla \f\|_{L^2(Q)}=1}\left| \mathcal{R}_{\rm lin}(\varphi)\right|\nonumber
\\
&\leq \Big[C_{\,\!_{\Omega}}^{\mathrm{P}}\,h_\Omega\,
\|(\mathbb{I}-\Pi^0_h)(\partial_t b(s^i_{h\tau})-f)\|_{L^2(Q)}
+ \|\mathbf{F}^i_{h\tau}+\boldsymbol{\sigma}^i_{h\tau}\|_{L^2(Q)}\Big]\overset{\eqref{flux_oscillation_estimator}}{=}\;
\eta^i_{\rm osc} \;+\; \eta^i_{\rm F}.
\end{align}
It remains to bound the total error. From the reliability bound 
\eqref{reliability_bound} of \Cref{theo:decomp},
\[
\mathcal{E}(s^i_{h\tau}) 
\leq \mathcal{E}^{\rm H}_{\rm disc}(s^{i+1}_{h\tau}) + L\,\mathcal{E}^i_{\rm lin}.
\]
The first term, defined in \eqref{eq:disc_error}, is bounded by \eqref{eq:Hminus1-R}, giving $\mathcal{E}^{\rm H}_{\rm disc}(s^{i+1}_{h\tau})\leq \eta^i_{\rm disc}$.
Applying \eqref{eq:defHinv_norm} we have $\|\p_t(s^{i+1}_{h\tau}-s^i_{h\tau})\|_{L^2(0,T;H^{-1}(\Om))}\leq C_{\,\!_{\Omega}}^{\mathrm{P}}\, h_\Om\|\p_t(s^{i+1}_{h\tau}-s^i_{h\tau})\|_{L^2(Q)}$. This implies for the linearization error by \eqref{err:linearization} that
\[\mathcal{E}^i_{\rm lin} = \|s^{i+1}_{h\tau}-s^i_{h\tau}\|_{\mathcal{V}} \leq  \eta^i_{\rm lin}.\]
Combining these two bounds, we obtain the estimate \eqref{eq:reliability_estimate}.
\end{proof}

\begin{remark}[Effectivity of the estimator $\eta^i$] The estimator $\eta^i$ is fully and parallelly computable, and does not involve any unknown constants (Poincar\'e constant is generally known, e.g., for convex domains, this is $C_{\,\!_{\Omega}}^{\mathrm{P}}=1/\pi$, \cite{payne1960optimal}). 
However, generally the flux estimator $\eta_{\mathrm{F}}^i$ is much greater than $\eta_{\mathrm{osc}}^i$ (unless the process is multiscale), particularly since $\p_t b(s^i_\htt)|_K=\p_t s^i_\htt|_K \in \calP_0(K)$ if $s^i_\htt\leq u^*$ in $K\in \calT_\htt$ from \Cref{lemma:b_B_construction}. Moreover, the linearization estimator $\eta_{\rm lin}^i$ vanishes as $i\to 0$. Thus, overall the estimators are expected to be quite accurate. This is verified in \Cref{fig:true_error_vs_estimator}, which shows that the effectivity index (error/estimator) is of order 1 in the non-degenerate regions ($s>0$) for the porous medium equation. However, the figure also reveals that proving local efficiency of the estimator is practically impossible at the degenerate regions where $s=0$, since the error is essentially 0. It is also possible to increase accuracy of the estimator further by taking higher degree projections $\Pi^p_h$ by increasing the degree $p$ of $\calQ_\htt$ in \Cref{def:Qhtspace}. 
\end{remark}

\begin{remark}[Global Poincaré constants in \Cref{theo:lsp} and efficiency estimate]\label{remark:efficiency} The use of the global Poincarè constant becomes necessary for the term $\int_0^T ((\mathbb{I}-\Pi_h^0) (\p_t b(s^i_\htt) - f),\f)$ since $\f\in L^2(0,T;H^1_0(\Om))$ only. However, if $\f\in H^1(Q)$, then we could have estimated this term by 
\begin{align*}
    &\int_0^T ((\mathbb{I}-\Pi_h^0) (\p_t b(s^i_\htt) - f),\f)=\sum_{K\in \calT_\htt} \left((\mathbb{I}-\Pi_h^0) (\p_t b(s^i_\htt) - f),\f-\tfrac{1}{|K|}\smallint_K \f \right)_K\\
    &\leq \sum_{K\in \calT_\htt}\|(\mathbb{I}-\Pi_h^0) (\p_t b(s^i_\htt) - f)\|_{L^2(K)}\left\|\f-\tfrac{1}{|K|}\smallint_K \f\right\|_{L^2(K)}\\
    &\leq \sum_{K\in \calT_\htt} \frac{h_K}{\pi} \|(\mathbb{I}-\Pi_h^0) (\p_t b(s^i_\htt) - f)\|_{L^2(K)}\left\|\f\right\|_{H^1(K)} \leq [\hat{\eta}_{\mathrm{osc}}^i] \|\f\|_{H^1(Q)}
\end{align*}
where $\hat{\eta}_{\mathrm{osc}}^i:=\left(\sum_{K\in \calT_\htt} \frac{h_K^2}{\pi^2}\|(\mathbb{I}-\Pi_h^0) (\p_t b(s^i_\htt) - f)\|_{L^2(K)}^2\right)^{\frac{1}{2}} $.
This could be useful in obtaining an global efficiency bound on $\calE(s^i_\htt)$ since $\|\mathcal{R}_{\rm lin}\|_{L^2(0,T;H^{-1}(\Omega))} \geq \|\mathcal{R}_{\rm lin}\|_{H^1(Q)^*}$.
\end{remark}

% Let $\calT_a$ denote the set of elements that have $a\in \mathcal{N}_{h\t}$ as a vertex. 
% We define the space-time patch and time-slice associated with the vertex $a \in\mathcal{N}_{h\tau}$ as 
% \[
% \oma := \mathrm{int}\left(\bigcup_{K\in \calT_{a}} \bar{K}\right), \qquad \omt:= \{(\bm{x},t_a)\in \oma \}, \qquad \tsl{\calT}:= \bigcup_{K\in \calT_a} (\omt\cap K),
% \]
%  i.e., $\oma$ is the union of all simplices sharing the vertex $a$ and $\omt$ is the constant time ($t=t_a$) slice through the patch $\oma$ that includes $a$ (see Figure~\ref{fig:composite_flux}) \KM{[[Reference to figure]}. Since $K\in \calT_{a}$ are simplices, the slice $\omt$ through the vertex $a$ yields $(\omt\cap K)$ that are also simplices in one dimension less. Thus, their collection defines a valid triangulation $\tsl{\calT}$ of $\omt$.
 
%  Let $RT_p(\tsl{K})$ denote the Raviart-Thomas space of order $p$ on the element $\tsl{K}\in \tsl{\calT}$. 
 
 
%  which will be used in the equilibrated flux reconstruction.


% % In this section, we construct a fully computable a posteriori error estimator to guide adaptive mesh refinement in the space-time setting. A uniform refinement of the entire space-time domain is computationally prohibitive; hence, it is essential to identify and refine only those regions that contribute most to the total discretization error. 

% The proposed estimator is based on the equilibrated flux reconstruction developed in \Cref{lem:balance-assembled}, which ensures local mass conservation on each vertex patch and yields a consistent flux source pair $(\boldsymbol{\sigma}_{h\tau}^{\,i}, \ell_{h\tau}^{\,i})$. This reconstruction allows the formulation of reliable and locally efficient residual based indicators that measure the imbalance in the discrete conservation law (cf.~\eqref{eq:reconstruction-balance}). These local indicators serve as the key components of the adaptive refinement strategy, described in Algorithm~\ref{alg:Lscheme-adapt}. 
% At each iteration, the indicators identify the elements where the residuals are largest, and only those elements are refined. In this way, the adaptive procedure focuses computational effort on the most inaccurate space-time regions while keeping the overall cost moderate, resulting in a robust and efficient discretization framework.

% \noindent
% %To implement the adaptive refinement procedure described above, we first define the space-time discretization setting on which the error indicators are evaluated. Let $\mathcal{T}_h$ denote a conforming triangulation of the space-time domain $\Omega_T := \Omega \times (0,T)$. Each element $T \in \mathcal{T}_h$ represents a space-time cell over which the local residuals and flux reconstructions are computed.
% %Let $\mathcal{T}_h$ be a conforming triangulation of the space--time domain.  
% %We denote by $T \in \mathcal{T}_h$ a generic space--time cell.
% %\begin{figure}[h!]
%  % \centering
%  % \begin{subfigure}[t]{0.47\linewidth}
%  %   \centering
%  %   \includegraphics[height=\linewidth, keepaspectratio]{Figures/triangulation_mesh.eps}
%   %  \caption{Uniform triangulation $\mathcal{T}_h^0$.}
%    % \label{fig:Th0}
%   %\end{subfigure}\hfill
%   %\begin{subfigure}[t]{0.47\linewidth}
%   %  \centering
%   %  \includegraphics[height=\linewidth, keepaspectratio]{Figures/Th_local_refine.eps}
%   %  \caption{Locally refined mesh $\mathcal{T}_h^1$, representing regions of high estimated error.}
%   %  \label{fig:Th1}
%   %\end{subfigure}
%   %\caption{[\Cref{sec:apost}] Example of adaptive refinement in the space-time domain. The right graph shows a sample refinement pattern to illustrate the adaptive mechanism; in practice, refinement is guided by local error indicators (see \Cref{thm:apost-reliability-i}).}
%   %\label{fig:adapt_mesh}
% %\end{figure}
% %\noindent
% %Each prism $T$ is further subdivided into simplices $K^a$ by cutting perpendiculars to the time axis through the space-time vertex $a \in \mathcal{V}_h$, where $\mathcal{V}_h$ denotes the set of vertices of the triangulation. We denote by
% %\[
% % \tilde{\boldsymbol{x}} = (\boldsymbol{x},t) \in \mathbb{R}^{d+1}
% %\]
% %a generic space-time point, where $\boldsymbol{x} \in \mathbb{R}^d$ is the spatial component and $t$ the temporal component. The space-time patch associated with a vertex $a \in \mathcal{V}_h$ is
% %\[
% % \tilde{\omega}_a := \bigcup_{K^a} K^a,
% %\]
% %Around each vertex $a \in \mathcal{V}_h$, we define the local patch $\tilde{\omega}_a$ as the union of all elements sharing $a$.
% %\begin{figure}[h!]
% %\centering
% %\includegraphics[width=0.35\linewidth]{Figures/patch_wa.png}
% %\caption{[\Cref{sec:apost}] The vertex-centered patch %$\tilde{\omega}_a$ in the space time mesh.}
% %\label{fig:patch_wa}
% %\end{figure}
% %\noindent
% %These vertex-centred patches serve as the computational domains for the equilibrated flux reconstruction introduced in \Cref{lem:balance-assembled,lemma_welldef}. We introduce the local finite element spaces on each patch $\tilde{\omega}_a$, 
% %\[
% % \mathcal{V}_{h\tau}:=
% % \left\{ v \in H^1(\Omega \times (0,T)) \;\middle|\;
% %    v|_T \in \mathbb{P}_p(T), \quad \forall \quad T \in \mathcal{T}_{h\tau} \right\},
% %\]
% %\[
% % \boldsymbol{Q}_{h\tau}^a :=
% % \left\{ \boldsymbol{v}_h \in H(\mathrm{div},\tilde{\omega}_a;\mathbb{R}^d) 
% % \;\middle|\;
% %    \boldsymbol{v}_h|_{K^a} \in RT_p(K^a),  \quad \forall \quad K^a \in \tilde{\omega}_a \right\}.
% %\]
% \subsubsection{Composite \texorpdfstring{$H(\mathrm{div})$}{H(div)} Flux}\label{sec:Hdiv_flux}

% On each patch $\tilde{\omega}_a$, we introduce the local flux space
% \[
% \boldsymbol{Q}_{h\tau}^a :=
% \left\{ \boldsymbol{v}_h \in H(\mathrm{div},\tilde{\omega}_a;\mathbb{R}^d)
% \;\middle|\;
%    \boldsymbol{v}_h|_{K^a} \in RT_p(K^a), \quad \forall K^a \subset \tilde{\omega}_a \right\},
% \]
% where $RT_p(K^a)$ denotes the Raviart-Thomas space of order $p$ on $K^a$, which will be used in the equilibrated flux reconstruction.

% \begin{definition}[Composite $H(\mathrm{div})$ flux]
% Let $\boldsymbol{a} \in \mathcal{N}_{h\tau}$ be a vertex. In a local coordinate, assume $\boldsymbol{a}=(\boldsymbol{0},0)$, and let ${\boldsymbol{q}}^a_{h\tau} \in {\boldsymbol{Q}}^a_{h\tau}$. We define the composite operator
% \[
%  I^a \boldsymbol{q}^a_{h\tau} : Q \longrightarrow \mathbb{R}^d .
% \]
% For $\boldsymbol{x} \in K^a \subset \tilde{\omega}_a$, let $\boldsymbol{b}$ be the highest (or lowest) vertex of 
% $K^a \subset \{t>0\}$ (or $K^a \subset \{t<0\}$) with respect to the time coordinate.  
% Connect $\boldsymbol{b}$ to $\boldsymbol{x}$ and extend this construction to the entire space-time patch $\tilde{\omega}_a = (\boldsymbol{x}_0,0)$.  
% Observe that $\boldsymbol{x}_0=\boldsymbol{0}$ is possible (see Fig.~\ref{fig:composite_flux}). We then define
% \begin{align}\label{eq:flux_psi}
%  I^a{\boldsymbol{q}}^a_{h\tau}(\boldsymbol{x})
%  := {\boldsymbol{q}}^a_{h\tau}(\boldsymbol{x}_0)\,\phi_a(\boldsymbol{x}),
% \end{align}
% where $\phi_a$ is the hat function associated with $a$. If $\boldsymbol{x}\notin \tilde{\omega}_a$, we set $I^a{\boldsymbol{q}}^a_{h\tau}(\boldsymbol{x}) = \boldsymbol{0}$.
% \begin{figure}[h]
% \centering
% \includegraphics[width=0.3\linewidth]{Figures/space_time_vertex_patch_1.png}
% \includegraphics[width=0.3\linewidth]{Figures/space_time_vertex_patch_2.png}

% \caption{
% Geometric illustration of the space-time patch construction.
% %The vertex $a=(x_0,0)$ lies at $t=0$ and defines a local space-time patch. 
% %For a point $x$ in the patch, the construction follows the segment connecting $a$ to $x$, 
% %while $x_0$ denotes the spatial projection of $a$.
% }
% \label{fig:composite_flux}
% \end{figure}
% \end{definition}
% \begin{lemma}
% Let $\boldsymbol{q}^a_{h\tau} \in \boldsymbol{Q}^a_{h\tau}$. 
% Then the composite flux $I^a \boldsymbol{q}^a_{h\tau}\in \mathcal{P}_{p+2}(\calT_{h\t})$ satisfies
% \begin{align*}
%  I^a \boldsymbol{q}^a_{h\tau}(\cdot,t) &\in H(\mathrm{div};\Omega;\mathbb{R}^d)
%  \quad \text{for all } t \in (0,T).
% \end{align*}
% \end{lemma}
% \begin{proof}
% By \eqref{eq:flux_psi}, $I^a \boldsymbol{q}^a_{h\tau}$ is a polynomial of degree at most $p+2$ on each $K \in \mathcal{T}_{h\tau}$.
% To verify $H(\mathrm{div})$-conformity, it remains to check continuity of the normal component across interior faces. Let $F$ be an interior face shared by two elements. Since $\boldsymbol{q}^a_{h\tau} \in \boldsymbol{Q}^a_{h\tau}$ has continuous normal components and $\phi_a$ is continuous across $F$, the normal component $I^a \boldsymbol{q}^a_{h\tau} \cdot \boldsymbol n$ has no jump across $F$. Using the continuity of both factors, we have
% \[
% \llbracket (I^a \boldsymbol{q}^a_{h\tau}) \cdot \boldsymbol n \rrbracket_F
% =  \llbracket (\boldsymbol{q}^a_{h\tau}\cdot \boldsymbol n)\,\phi_a \rrbracket_F = 0,
% \]
% Hence $I^a \boldsymbol{q}^a_{h\tau} \in H(\mathrm{div};Q;\mathbb{R}^d)$, and the time-slice property follows by restriction to each $t \in (0,T)$.

% \end{proof}

% This result shows that the reconstructed flux is of polynomial degree $p+2$ on each element and, for each $t \in (0,T)$, belongs to the space $H(\mathrm{div}; \Omega;\mathbb{R}^d)$. In particular, it guarantees both polynomial consistency and $H(\mathrm{div})$--conformity of the composite flux construction.
% \noindent
% %We consider the finite element solution and seek $S_{h\tau}$ and $\boldsymbol{F}_{h\tau}$ such that, 
% %We now introduce the equilibrated flux and potential associated with the discrete solution for all $\varphi_{h\tau} \in \mathcal{V}_{h\tau}$, 

% Let $(S_{h\tau}, \boldsymbol{F}_{h\tau})$ denote the finite element solution satisfying
% \begin{align}\label{eq:Galerkin solution}
%  \int_0^T ( S_{h\tau}, \varphi_{h\tau} ) \, dt
%  + \int_0^T ( \boldsymbol{F}_{h\tau}, \nabla \varphi_{h\tau} ) \, dt = 0,
%  \quad \forall \varphi_{h\tau} \in \mathcal{V}_{h\tau}.
% \end{align}
% %Furthermore, let
% %\[
% % \Pi^p_{h\tau} : L^2((0,T)\times \Omega) \longrightarrow \mathbb{P}_p(\mathcal{T}_{h\tau})
% %\]
% %denote the $L^2$ orthogonal projection onto the space of discontinuous piecewise polynomials of degree $p$ over the space-time mesh $\mathcal{T}_{h\tau}$. 
% To obtain guaranteed and locally computable a posteriori error estimates, we now introduce equilibrated reconstructions of the flux and potential associated with this discrete solution.
% \subsubsection{Equilibrated flux and potential}

% For each vertex $a \in \mathcal{N}_{h\tau}$, we define the local space
% \[
%  \hat{\mathcal{H}}^a_{h\tau}
%  := \Big\{ \big(\rho^a_{h\tau}, \boldsymbol{q}^a_{h\tau}\big) 
%     \in \mathbb{P}_{p+1}(\tilde{\omega}_a)\times \boldsymbol{Q}^a_{h\tau}
%     \Big\},
% \]
% such that
% \begin{align}\label{eq:rho_def}
%    \rho^a_{h\tau} - \nabla\cdot \big(I^a \boldsymbol{q}^a_{h\tau}\big)
%    = \Pi^{p+1}_{h\tau}\!\big(\phi_a S_{h\tau}\big)
%      + \nabla \phi_a \cdot \Pi^{p+1}_{h\tau}\boldsymbol{F}_{h\tau}.
% \end{align}
% Then the patchwise equilibrated flux and potential 
% $(\ell^a_{h\tau}, \boldsymbol{\sigma}^a_{h\tau})$ 
% are obtained as the solution of the local minimization problem
% \begin{align}\label{eq:minimization_prob}
%  (\ell^a_{h\tau}, \boldsymbol{\sigma}^a_{h\tau})
%  := \underset{(\rho^a_{h\tau},\boldsymbol{q}^a_{h\tau})\in \hat{\mathcal{H}}^a_{h\tau}}%
%     {\operatorname{arg\,min}}
%  \Big( \| I^a \boldsymbol{q}^a_{h\tau} + \phi_a \boldsymbol{F}_{h\tau}\|^2
%       + h_\Omega \|\rho^a_{h\tau}\|^2_{\tilde{\omega}_a} \Big).
% \end{align}
% Then the global equilibrated flux and potential are defined by
% \begin{align}\label{eq:global_equilibrated_flux_potential}
%      \boldsymbol{\sigma}_{h\tau}
%  := \sum_{a\in \mathcal{N}_{h\tau}} I^a \boldsymbol{\sigma}^a_{h\tau},
%  \qquad
%  \ell_{h\tau}
%  := \sum_{a\in \mathcal{N}_{h\tau}} \ell^a_{h\tau}.
% \end{align}
% It is important to verify that this construction is well-posed, i.e., the equilibrated flux and potential are uniquely defined. This is the content of the following result.
% %%%%%%%%%%%%%%%%%%%%%%%%%%%%%%%%%%%%%%%%%%%%%%%%%%%%%%%%%%%%%%%%%%%%%%%%%%%%%%%%%%%%%%%%%%%%%%%%%%%%%

% \begin{lemma}\label{lemma_welldef}
% The pair $(\boldsymbol{\sigma}_{h\tau}, \ell_{h\tau})$ 
% is well defined.
% \end{lemma}


% \begin{proof}
% For each $a \in \mathcal{N}_{h\tau}$, the constraint \eqref{eq:rho_def} gives
% \begin{align}\label{rho_expression}
%  \rho^a_{h\tau}
%  = \Pi^{p+1}_{h\tau}(\phi_a S_{h\tau})
%    + \nabla \phi_a \cdot \Pi^{p+1}_{h\tau}\boldsymbol{F}_{h\tau}
%    + \nabla \cdot I^a \boldsymbol{q}^a_{h\tau}.
% \end{align}
% Thus, for a given $\boldsymbol{q}^a_{h\tau}$, $\rho^a_{h\tau}$ is uniquely determined. Hence, the minimization problem \eqref{eq:minimization_prob} reduces to a problem in the variable $\boldsymbol{q}^a_{h\tau}$ only.

% Substituting \eqref{rho_expression} into \eqref{eq:minimization_prob}, we obtain the functional
% \[
%  \mathcal{J}(\boldsymbol{q}^a_{h\tau})
%  =
%  h_\Omega \Big\|
%      \nabla\cdot(I^a {\boldsymbol{q}}^a_{h\tau})
%      + \Pi^{p+1}_{h\tau}(\phi_a S_{h\tau})
%      + \nabla \phi_a \cdot \Pi^{p+1}_{h\tau}\boldsymbol{F}_{h\tau}
%  \Big\|^2_{\tilde{\omega}_a}
%  + \big\| I^a \boldsymbol{q}^a_{h\tau} 
%          + \phi_a \boldsymbol{F}_{h\tau} \big\|^2_{\tilde{\omega}_a}.
% \]
% This functional is a sum of squared $L^2$-norms of linear operators applied to $\boldsymbol{q}^a_{h\tau}$, and is therefore a strictly convex quadratic functional on the finite-dimensional space $\boldsymbol{Q}^a_{h\tau}$. Consequently, it admits a unique minimizer $\boldsymbol{\sigma}^a_{h\tau} \in \boldsymbol{Q}^a_{h\tau}$, and the corresponding $\rho^a_{h\tau}$ is uniquely determined by \eqref{rho_expression}. We define
% \[
% \ell^a_{h\tau} := 2h_\Omega \rho^a_{h\tau}.
% \]

% Moreover, the pair $(\boldsymbol{\sigma}^a_{h\tau}, \ell^a_{h\tau})$ can equivalently be characterized as the solution of the following local mixed problem on $\tilde{\omega}_a$: find $(\boldsymbol\sigma^a_{h\tau},\ell^a_{h\tau})\in 
% \boldsymbol Q^a_{h\tau}\times\mathbb P_{p+1}(\tilde{\omega}_a)$ such that, for all $\widehat{\boldsymbol q}^a\in\boldsymbol Q^a_{h\tau}$ and $p^a\in\mathbb P_{p+1}(\tilde{\omega}_a)$,
% \begin{align}
% (\ell^a_{h\tau},\,\nabla\!\cdot(I^a\widehat{\boldsymbol q}^a))_{\tilde{\omega}_a}
% +(I^a\boldsymbol\sigma^a_{h\tau},\,I^a\widehat{\boldsymbol q}^a)_{\tilde{\omega}_a}
% &= -\,(\phi_a\boldsymbol F_{h\tau},\,I^a\widehat{\boldsymbol q}^a)_{\tilde{\omega}_a}, 
% \label{eq:local-mixed-1}\\[2mm]
% (\ell^a_{h\tau},\,p^a)_{\tilde{\omega}_a}
% -(\nabla\!\cdot(I^a\boldsymbol\sigma^a_{h\tau}),\,p^a)_{\tilde{\omega}_a}
% &=(\phi_a S_{h\tau}+\boldsymbol F_{h\tau}\!\cdot\nabla\phi_a,\,p^a)_{\tilde{\omega}_a}.
% \label{eq:local-mixed-2}
% \end{align}
% \noindent
% Therefore, we obtain a unique pair $(\boldsymbol\sigma^a_{h\tau},\ell^a_{h\tau})$ for each patch. The global quantities $(\boldsymbol{\sigma}_{h\tau}, \ell_{h\tau})$ are then defined by \eqref{eq:global_equilibrated_flux_potential}, and uniqueness follows by linearity of the assembly. 

% %%%%%%%%%%%%%%%%%%%%%%%%%%%%%%%%%%%%%%%%%%%%%%%%%%%%%%%%%%%%%%%%%%%%%%%%%%%%%%%%%%%%%%%%%%%%%%%%%%%%%
% \iffalse
% We introduce the local Lagrangian functional
% \begin{align}\label{eq:lagranian}
%  \mathcal{L}(\boldsymbol{q}^a_{h\tau}, \rho^a_{h\tau}; \lambda^a_{h\tau})
%  := \mathcal{J}(\boldsymbol{q}^a_{h\tau}, \rho^a_{h\tau})
%  + \big(\lambda^a_{h\tau}, \,
%    \rho^a_{h\tau} - \nabla\!\cdot(I^a \boldsymbol{q}^a_{h\tau})
%    - \Pi^{p+1}_{h\tau}(\phi_a S_{h\tau})
%    - \nabla \phi_a \cdot \Pi^{p+1}_{h\tau}\boldsymbol{F}_{h\tau}
%  \big)_{\omega_a}
% \end{align}
% where $\lambda^a_{h\tau}\in\mathbb P_{p+1}(\omega_a)$ acts as a Lagrange multiplier enforcing the constraint \eqref{rho_expression}. Let $\widehat{\boldsymbol q}^a\in\boldsymbol Q^a_{h\tau}$ be arbitrary, and we perturb the flux by a small scalar $\varepsilon$ and set
% $\Phi(\varepsilon):=\mathcal L(\boldsymbol q^a_{h\tau}+\varepsilon\widehat{\boldsymbol q}^a,
% \rho^a_{h\tau};\lambda^a_{h\tau})$. In \eqref{eq:lagranian} only two terms depend on $\boldsymbol q^a_{h\tau}$, \\
% \emph{Squared-norm term.}
% Set, $A:=I^a$ and $b:=\phi_a\boldsymbol F_{h\tau}$.
% Then
% \[
% \|A(\boldsymbol q^a_{h\tau}+\varepsilon\widehat{\boldsymbol q}^a)+b\|^2_{\omega_a}
% =\big(A\boldsymbol q^a_{h\tau}+b+\varepsilon A\widehat{\boldsymbol q}^a,\;
% A\boldsymbol q^a_{h\tau}+b+\varepsilon A\widehat{\boldsymbol q}^a\big)_{\omega_a}.
% \]
% Using bilinearity and symmetry of $(\cdot,\cdot)_{\omega_a}$,
% \begin{align*}
% \|A(\boldsymbol q_{h\tau}+\varepsilon\widehat{\boldsymbol q}^a)+b\|^2_{\omega_a}
% &=(u+\varepsilon v,\,u+\varepsilon v)_{\omega_a}
% \\
% &=\|u\|^2_{\omega_a}+2\varepsilon(u,v)_{\omega_a}+\varepsilon^2\|v\|^2_{\omega_a},
% \end{align*}
% where $u:=A\boldsymbol q^a_{h\tau}+b$ and $v:=A\widehat{\boldsymbol q}^a$.
% Hence
% \begin{align}\label{eq:squared_norm_term}
% \frac{d}{d\varepsilon}\Big|_{\varepsilon=0}\|A(\boldsymbol q^a_{h\tau}+\varepsilon\widehat{\boldsymbol q}^a)+b\|^2_{\omega_a}
% =2\,(A\boldsymbol q^a_{h\tau}+b,\;A\widehat{\boldsymbol q}^a)_{\omega_a}
% =2\,(I^a\boldsymbol q^a_{h\tau}+\phi_a\boldsymbol F_{h\tau},\,I^a\widehat{\boldsymbol q}^a)_{\omega_a}.
% \end{align}
% \emph{Constraint term:}
% Using linearity of $I^a$ and $\nabla\!\cdot$,
% \begin{align*}
% (\lambda^a_{h\tau},-\nabla\!\cdot(I^a(\boldsymbol q^a_{h\tau}+\varepsilon\widehat{\boldsymbol q}^a)))_{\omega_a}
% &=(\lambda^a_{h\tau},-\nabla\!\cdot(I^a\boldsymbol q^a_{h\tau}))_{\omega_a}
% -\varepsilon\,(\lambda^a_{h\tau},\nabla\!\cdot(I^a\widehat{\boldsymbol q}^a))_{\omega_a},
% \end{align*}
% thus
% \begin{align}\label{eq:Contraints_term}
% \frac{d}{d\varepsilon}\Big|_{\varepsilon=0}
% (\lambda^a_{h\tau},-\nabla\!\cdot(I^a(\boldsymbol q^a_{h\tau}+\varepsilon\widehat{\boldsymbol q}^a)))_{\omega_a}
% =-(\lambda^a_{h\tau},\,\nabla\!\cdot(I^a\widehat{\boldsymbol q}^a))_{\omega_a}.
% \end{align}
% Collecting both contributions from \eqref{eq:squared_norm_term} and \eqref{eq:Contraints_term}.
% Therefore,
% \begin{align}
% \frac{d}{d\varepsilon}\Big|_{\varepsilon=0}\Phi(\varepsilon)
% =
% 2\,(I^a\boldsymbol q^a_{h\tau}+\phi_a\boldsymbol F_{h\tau},\,I^a\widehat{\boldsymbol q}^a)_{\omega_a}
% -(\lambda^a_{h\tau},\,\nabla\!\cdot(I^a\widehat{\boldsymbol q}^a))_{\omega_a}.
% \end{align}
% At the minimizer $(\boldsymbol\sigma^a_{h\tau},\rho^a_{h\tau},\lambda^a_{h\tau})$ this derivative vanishes for all
% $\widehat{\boldsymbol q}^a\in\boldsymbol Q^a_{h\tau}$:
% \begin{equation*}
% -(\lambda^a_{h\tau},\,\nabla\!\cdot(I^a\widehat{\boldsymbol q}^a))_{\omega_a}
% +2\,(I^a\boldsymbol\sigma^a_{h\tau}+\phi_a\boldsymbol F_{h\tau},\,I^a\widehat{\boldsymbol q}^a)_{\omega_a}=0
% \quad\forall\,\widehat{\boldsymbol q}^a\in\boldsymbol Q^a_{h\tau}
% \end{equation*}
% Equivalently, splitting the second inner product,
% \begin{equation}\label{eq:KKT_q}
% -(\lambda^a_{h\tau},\,\nabla\!\cdot(I^a\widehat{\boldsymbol q}^a))_{\omega_a}
% +2\,(I^a\boldsymbol\sigma^a_{h\tau},\,I^a\widehat{\boldsymbol q}^a)_{\omega_a}
% = -2\,(\phi_a\boldsymbol F_{h\tau},\,I^a\widehat{\boldsymbol q}^a)_{\omega_a}
% \quad\forall\,\widehat{\boldsymbol q}^a.
% \end{equation}
% Let $r^a\in\mathbb P_{p+1}(\omega_a)$ and set
% $\Psi(\varepsilon):=\mathcal L(\boldsymbol\sigma^a_{h\tau},\rho^a_{h\tau}+\varepsilon r^a;\lambda^a_{h\tau})$. In \eqref{eq:lagranian} only two terms depend on $\rho^a_{h\tau}$:
% $h_\Omega\|\rho^a_{h\tau}\|^2_{\omega_a}$ and $(\lambda^a_{h\tau},\rho^a_{h\tau})_{\omega_a}$.
% We handle them one by one,
% \begin{align}
% \frac{d}{d\varepsilon}\Big|_{\varepsilon=0} h_\Omega\|\rho^a_{h\tau}+\varepsilon r^a\|^2_{\omega_a}
% = 2h_\Omega\,(\rho^a_{h\tau},r^a)_{\omega_a},
% \qquad
% \frac{d}{d\varepsilon}\Big|_{\varepsilon=0}(\lambda^a_{h\tau},\rho^a_{h\tau}+\varepsilon r^a)_{\omega_a}
% =(\lambda^a_{h\tau},r^a)_{\omega_a}.
% \end{align}
% Therefore
% \[
% \frac{d}{d\varepsilon}\Big|_{\varepsilon=0}\Psi(\varepsilon)
% =2h_\Omega\,(\rho^a_{h\tau},r^a)_{\omega_a}+(\lambda^a_{h\tau},r^a)_{\omega_a}=0
% \quad\forall\,r^a,
% \]
% and since $r^a$ is arbitrary, we obtain in $\mathbb P_{p+1}(\omega_a)$
% \begin{equation}\label{eq:ell_equals_rho}
% \boxed{\;\lambda^a_{h\tau}=-\,2h_\Omega\,\rho^a_{h\tau}.\;}
% \end{equation}
% Finally, let $p^a \in\mathbb P_{p+1}(\omega_a)$ be arbitrary and perturb the multiplier via $\lambda^a_{h\tau}\mapsto \lambda^a_{h\tau}+\varepsilon p^a$. With
% \[
% \Theta(\varepsilon):=\mathcal L(\boldsymbol\sigma^a_{h\tau},\,\rho^a_{h\tau};\,\lambda^a_{h\tau}+\varepsilon p^a),
% \]
% only the constraint term in \eqref{eq:lagranian} depends on $\lambda^a_{h\tau}$; hence
% \begin{align}
% \frac{d}{d\varepsilon}\Big|_{\varepsilon=0}\Theta(\varepsilon)
% =(p^a,\,
% \rho^a_{h\tau}-\nabla\!\cdot(I^a\boldsymbol\sigma^a_{h\tau})
% -\Pi^{p+1}_{h\tau}(\phi_a S_{h\tau})
% -\nabla\phi_a\cdot \Pi^{p+1}_{h\tau}\boldsymbol F_{h\tau})_{\omega_a}=0.
% \end{align}
% Since $p^a$ is arbitrary,
% \begin{equation}\label{eq:weak_constraint_withPi}
% (\rho^a_{h\tau},p^a)_{\omega_a}
% -(\nabla\!\cdot(I^a\boldsymbol q^a_{h\tau}),p^a)_{\omega_a}
% =(\Pi^{p+1}_{h\tau}(\phi_a S_{h\tau}),p^a)_{\omega_a}
% +(\nabla\phi_a\!\cdot \Pi^{p+1}_{h\tau}\boldsymbol F_{h\tau},p^a)_{\omega_a}
% \end{equation}
% Since, $\Pi^{p+1}_{h\tau}$ is the $L^2(\omega_a)$ orthogonal projector onto $\mathbb P_{p+1}(\omega_a)$, by definition of the $L^2$–projection we have
% \begin{equation}\label{eq:L2-proj}
% \big(g-\Pi^{p+1}_{h\tau}g,\,p\big)_{\omega_a}=0\quad\forall\,p\in\mathbb P_{p+1}(\omega_a)
% \quad\Longleftrightarrow\quad
% \big(\Pi^{p+1}_{h\tau}g,\,p\big)_{\omega_a}=\big(g,\,p\big)_{\omega_a},
% \end{equation}
% \begin{equation}\label{eq:weak_constraint_clean}
% (\rho^a_{h\tau},p^a)_{\omega_a}
% -(\nabla\!\cdot(I^a\boldsymbol q^a_{h\tau}),p^a)_{\omega_a}
% \overset{\eqref{eq:L2-proj}}{=}(\phi_a S_{h\tau},p^a)_{\omega_a}
% +(\boldsymbol F_{h\tau}\!\cdot\nabla\phi_a,\,p^a)_{\omega_a}
% \quad\forall\,p^a \in\mathbb P_{p+1}(\omega_a).
% \end{equation}
% Setting $\boldsymbol\sigma^a_{h\tau}$ to be the unique solution $\boldsymbol q^a_{h\tau}$ of the optimality system \eqref{eq:KKT_q}, \eqref{eq:ell_equals_rho}, \eqref{eq:weak_constraint_clean}, and defining the reconstructed potential on the patch by
% \begin{equation}\label{eq:potential_definition}
% \ell^a_{h\tau}:=-\,\lambda^a_{h\tau}=2h_\Omega\,\rho^a_{h\tau}\in\mathbb P_{p+1}(\omega_a),
% \end{equation}
% \noindent
% Substituting \eqref{eq:potential_definition} into \eqref{eq:KKT_q}, and \eqref{eq:weak_constraint_clean}
% yields the following closed mixed formulation.

% \paragraph{Local mixed problem on $\omega_a$.}
% Find $(\boldsymbol\sigma^a_{h\tau},\ell^a_{h\tau})\in 
% \boldsymbol Q^a_{h\tau}\times\mathbb P_{p+1}(\omega_a)$ such that
% for all $\widehat{\boldsymbol q}^a\in\boldsymbol Q^a_{h\tau}$ 
% and $p^a\in\mathbb P_{p+1}(\omega_a)$,
% \begin{align}
% (\ell^a_{h\tau},\,\nabla\!\cdot(I^a\widehat{\boldsymbol q}^a))_{\omega_a}
% +(I^a\boldsymbol\sigma^a_{h\tau},\,I^a\widehat{\boldsymbol q}^a)_{\omega_a}
% &= -\,(\phi_a\boldsymbol F_{h\tau},\,I^a\widehat{\boldsymbol q}^a)_{\omega_a}, 
% \label{eq:local-mixed-1}\\[2mm]
% (\ell^a_{h\tau},\,p^a)_{\omega_a}
% -(\nabla\!\cdot(I^a\boldsymbol\sigma^a_{h\tau}),\,p^a)_{\omega_a}
% &=(\phi_a S_{h\tau}+\boldsymbol F_{h\tau}\!\cdot\nabla\phi_a,\,p^a)_{\omega_a}.
% \label{eq:local-mixed-2}
% \end{align}
% \noindent
% we obtain a unique pair $(\boldsymbol\sigma^a_{h\tau},\ell^a_{h\tau})$ for each patch. The global objects are then defined by \eqref{eq:global_equilibrated_flux_potential}, and uniqueness follows by linearity of the assembly. 
% \fi
% \end{proof}
% 
\subsubsection{Space-time adaptive algorithm}
We outline an adaptive solution strategy based on the estimators derived in \Cref{thm:apost-reliability-i}, and the $\texttt{SOLVE} \;\longrightarrow\; \texttt{ESTIMATE} \;\longrightarrow\; \texttt{MARK} \;\longrightarrow\; \texttt{REFINE}$ principle \cite{carstensen2014axioms}.
On a given mesh $\mathcal{T}_{h\tau}^\ell$ ($\ell\in \N$), the nonlinear problem is solved by the L-scheme linearization (Problem \ref{pb:2}). Once the linearization error $\eta^i_{\rm lin}$   is below a prescribed tolerance, the equilibrated flux reconstruction $\boldsymbol{\sigma}^i_{h\tau}$ (see \Cref{def:equilibrated}) is performed, yielding discretization estimator $\eta^i_{\mathrm{disc}}$ given in \eqref{flux_oscillation_estimator}. If the linearization error $\eta^i_{\rm lin}$ is smaller than a fraction $\d_{\mathrm{lin}}\ll 1$ of the discretization error $\eta^i_{\mathrm{disc}}$, then a refinement step is done.
The element-wise error indicators $\eta^i_K$ for each element $K \in \mathcal{T}_{h\tau}$ are computed, based on \Cref{thm:apost-reliability-i}, comprising of flux, data oscillation, and initial condition contributions,
\begin{align}
    \eta^i_K := \left(2\left([\eta^i_{\mathrm{osc},K}]^2 + [\eta^i_{\mathrm{F},K}]^2\right) + L^2 \int_{K\cap (\{0\}\times \Om)} |s^i_{\htt}(0)-s_0|^2 \right)^{\frac{1}{2}} \text{ such that } \sum_{K\in \calT_\htt} [\eta^i_K]^2\geq [\eta^i_{\mathrm{disc}}]^2. \label{eq:etaK}
\end{align}
The elements that contribute the most to the total error are marked for refinement using the D\"{o}rfler criterion. The marked elements are refined using vertex bisection. After projecting to the new mesh $\calT^{\ell+1}_\htt$, the linear iterations are restarted. The resulting adaptive loop is summarized in Algorithm~\ref{alg:Lscheme-adapt}.

To evaluate the performance of the adaptive strategy, along with the adaptive algorithm, we also investigate a \emph{uniform refinement} strategy. In this case, a uniform refinement is done by a factor $\mathrm{r}_\mathrm{ref}=2$ if the adaptive criterion in Algorithm~\ref{alg:Lscheme-adapt},
\[\eta_{\mathrm{lin}}^i\le \delta_{\mathrm{lin}}\eta^i_{\mathrm{disc}},\]
is met. As evaluation criteria, we consider both the number of iterations, as well as the degree of freedom for each mesh. More precisely, for any mesh-level $\ell\in \N$, let $\overline{i}(\ell)\in \N$ be the iteration at which the adaptive criterion is met both for the adaptive and the uniform strategies. Letting $|\mathrm{DOF}|_\ell$ denote the degrees of freedom of the corresponding finite element space on $\calT^\ell_\htt$, to assess the adaptive strategy, we consider the total cost, respectively the total number of iterations. For the $i^{\mathrm{th}}$ iteration at the $\ell^{\rm th} \text{level}$ we define
\begin{equation}
  \label{eq:cost}\mathrm{Cost}(i,\ell)
:= \sum_{\ell'=0}^{\ell-1}\sum_{i'=1}^{\overline{i}(\ell')}|\mathrm{DOFs}|_{\ell'} + \sum_{i'=1}^i |\mathrm{DOFs}|_{\ell},\qquad \mathrm{IterNo}(i,\ell):=\sum_{\ell'=0}^{\ell-1} \overline{i}(\ell') + i.
\end{equation}
%The cost represents the combined degree of freedom used over all iterations and mesh-levels.

\textcolor{black}{Here, $\mathrm{Cost}(i,\ell)$ represents the accumulated degrees of freedom over all iterations at previous mesh levels $\ell'=0,\ldots,\ell-1$ until the stopping criterion is met, plus those at the current level $\ell$ up to iteration $i$. Similarly, $\mathrm{IterNo}(i,\ell)$ is the total number of iterations performed across all mesh levels through iteration $i$ at level $\ell$.}
 

\begin{algorithm}[h!]
\caption{\vspace{0.4em}Iterative scheme with equilibrated--flux a posteriori adaptivity\vspace{0.5em}}
\label{alg:Lscheme-adapt}
\begin{minipage}{\linewidth}
\begin{algorithmic}[1]
\Require initial mesh $\mathcal{T}_{h\tau}^0$, initial guess $s^0$, and parameters: $L>1$, tolerance $\texttt{Tol}\ll 1$, maximum iteration $i_{\max}$, maximum level $\ell_{\max}$, D\"orfler parameter 
$0<\texttt{$\theta$}<1$, refinement factor $\texttt{$\mathrm{r}_\mathrm{ref}$}>1$, balance factor
$0<\texttt{$\delta_{\mathrm{lin}}$}<1$, %$\varepsilon_{\mathrm{sol}}>0$,
and an entry threshold 
$0<\texttt{$\varepsilon_{\mathrm{adapt}}$}\le 0.1$.
%\Ensure iterates $(s^{i})_{i\ge0}$ and a sequence of adapted meshes.
\State \textbf{set} $\texttt{Error}=1,\;\;$ $i=0,\;\;$, $\ell=0$
\While{$i<\texttt{$i_{\texttt{max}}$}\;\;$ or $\;\;\texttt{Error}\ge\texttt{Tol}\;\;$ or  $\;\;\ell<\ell_{\max}$\;\;} \Comment{nonlinear (L scheme) loop}
  \State Solve linearized problem for $s^{i+1}_\htt$ on $\mathcal T_{h\tau}^\ell$ \hfill \Comment{(\texttt{SpaceTime} solve)(see~\eqref{eq:linearization_gen})}
  \State Compute $\eta_{\mathrm{lin}}^i$ \Comment{linearization error}

  %=================SKIP THIS STEP============
  %===========================================
  \If{$\eta^i_{\mathrm{lin}}<\e_{\mathrm{adapt}}$}\hfill  \Comment{enter adaptive (refinement) loop}
      \State Reconstruct $\boldsymbol{\sigma}^{\,i}_{h\tau}$ 
      \hfill (local patch problem, see \Cref{def:equilibrated})
     \State Compute $\eta^i_K$, $\eta^i_{\mathrm{disc}}$, and $\eta^i\mapsto \texttt{Error}$
      \hfill (see~\Cref{thm:apost-reliability-i},\eqref{eq:etaK})
     \If{$\eta_{\mathrm{lin}}^i\le \delta_{\mathrm{lin}}\cdot\eta^i_{\mathrm{disc}}$} 
     \Comment{Mesh adaptivity gate}
      \State \textbf{Dörfler mark:} identify minimal $\mathcal{M}\subset\mathcal{T}_{h\tau}^\ell$ with
       $\sum_{K\in\mathcal M}(\eta^i_K)^2 \ge \theta^2\sum_{K\in\mathcal{T}_{h\tau}^\ell}(\eta^i_K)^2$
       %=======================================
         %===================================
        \State \textbf{Refine} marked elements 
               $K\in\mathcal{M}$ by setting 
               $h_K^{\mathrm{new}} = h_K / \texttt{$\mathrm{r}_\mathrm{ref}$}$
        \State \textbf{Update mesh} $\mathcal T_{h\tau}^{\ell+1}\gets\textsc{adaptmesh}(\mathcal T_{h\tau}^\ell,\{h_K^{\mathrm{new}}\}_{\mathcal{M}})$  \Comment{generate new mesh}
         \State \textbf{Post-refinement:} $\ell\gets (\ell+1)$. Project $s^{i+1}_{h\tau}, s^{i}_{h\tau}$ to the new $\calV_\htt^{\ell}$ on $\calT_{h\tau}^\ell$
        %======================================
     \EndIf
  \EndIf
  \State Update $s^{i}\gets s^{i+1}$, $i\gets i+1$  \Comment{next iteration}
  \EndWhile
\end{algorithmic}
\end{minipage}
\end{algorithm}
%====================================================
\noindent
In the following section, we compare the performance of these refinement strategies with respect to $\mathrm{Cost}(i,\ell)$ and $\mathrm{IterNo}(i,\ell)$ across several numerical test cases, demonstrating their robustness and efficiency. For cases when the exact solution $u_{\mathrm{exact}}$ is known, we would define the element-wise error and local effectivity index of the estimator as
\begin{align}\left\{\begin{aligned}
        &\|e\|_K^2:= \|b(s^i_\htt)-u_{\mathrm{exact}}\|_K^2 + \|\nabla (B(s^i_\htt)-\Phi(u_{\mathrm{exact}}))\|_K^2, \\
    &(\text{Effectivity index})_K:= \|e\|_K/\eta_K^i,
\end{aligned}
\right. \qquad \text{ for all } K\in \calT_\htt.\label{eq:effectivity}
\end{align}

\section{Numerical Study of the Adaptive Space-Time Scheme}\label{sec:Numerical_results}

This section illustrates the performance of the space-time linearization scheme combined with the adaptive algorithm of Section~\ref{sec:apost}. As a baseline for comparison, we consider both the L-scheme and a regularized Newton scheme, accompanied by adaptive or uniform refinements.

\subsubsection{A modified Newton ($N_{\mathrm{reg}}$) scheme}
For the numerical tests, we also consider a linear iterative scheme that can be seen as a regularized Newton ($N_{\mathrm{reg}}$) scheme. This provides a baseline for comparisons with the L-scheme. The iteration in the regularized Newton scheme can be interpreted as the solution to the following problem.
\begin{problem}
Let $\epsilon\geq0$ be fixed. For $i\in \N$ and $s^i\in (\calS\cap \calV)$, find $s^{i+1}\in \calV$ solving
\begin{align}\label{eq:M-scheme}
\int_0^T \!\left\langle \partial_t \left( L^i_b(s^{i+1} - s^i) + b(s^i) \right), \varphi \right\rangle \,  
+ \int_0^T \!\left( \nabla \left( L^i_B(s^{i+1} - s^i) + B(s^i) \right), \nabla \varphi \right) 
= \int_0^T \!( f, \varphi ),
\end{align}
\end{problem}
with $L^i_b:= b'(s^i) + \varepsilon$ and $L^i_B:=B'(s^i) +\varepsilon$.

Observe that the $\varepsilon=0$ case corresponds to the Newton scheme. However, below we keep $\e>0$ since, for the degenerate test cases investigated, the $\e=0$ option always diverges. This is either due to $b'$ or $B'$ being 0, which makes the time-derivative or space-derivative terms vanish, and therefore the Jacobian becomes singular.
Note that this regularization is similar to the $M$-scheme analyzed in \cite{mitra2019modified,javed2025robust}

The adaptive estimator $\eta^i$ can also be computed following the analysis in \Cref{sec:a-posteriori}. More precisely, one can change the $L$-scheme in \eqref{eq:source_flx} by considering
\begin{align}\label{eq:ShFh_M}
\Shti=
\p_t(L^i_b(s^{i+1} - s^i) + b(s^i) )-f, \qquad
\Fhti=\nabla (L^i_B (s^{i+1} - s^i) + B(s^i))
\end{align}
taking $L$ constant in \eqref{eq:source_flx_Lscheme}. Note that the proofs of \Cref{theo:converegence,theo:decomp,theo:lsp,thm:apost-reliability-i} do not extend to this scheme, so the convergence of the scheme and the reliability of the estimator cannot be guaranteed. However, it will serve as an important point of comparison.

In the following, we denote by \textbf{Adapt-L} and \textbf{Adapt-$N_{\mathrm{reg}}$} the adaptive mesh refinement strategies for the L and regularized Newton schemes, respectively; by \textbf{Uni-L} and \textbf{Uni-$N_{\mathrm{reg}}$} the uniform mesh refinement for L and regularized Newton schemes; and by \textbf{L (fixed mesh)} and \textbf{$N_{\mathrm{reg}}$ (fixed mesh)}, the L and regularized Newton schemes on a fixed mesh. 

% \subsection{Notations for numerical experiments}\label{sec:notations_numerical_scheme}
% We summarize the notation used for the numerical schemes in \Cref{tab:notatin_numerical_scheme}.
% \begin{table}[h!]
% \begin{center}
% \begin{tabular}{cp{12cm}}
% \textbf{Schemes} & \textbf{Description} \\[0.5em]
% Adapt-L & Adaptive mesh refinement with the $L$-scheme \\[0.2em]
% Adapt-M & Adaptive mesh refinement with the regularized Newton ($M$-) scheme \\[0.2em]
% Uni-L   & Uniform mesh refinement with the $L$-scheme \\[0.2em]
% Uni-M   & Uniform mesh refinement with the regularized Newton ($M$-) scheme \\[0.2em]
% L       & $L$-scheme on a fixed (non-adaptive) mesh \\[0.2em]
% \end{tabular}
% \caption{[\Cref{sec:notations_numerical_scheme}] Notation of the numerical schemes, including adaptive and uniform refinement with the $L$-scheme and the regularized Newton ($M$-) scheme.}
% \label{tab:notatin_numerical_scheme}
% \end{center}
% \end{table}

\subsubsection{Numerical setup}\label{subsec:num-setup}
All numerical experiments are implemented in 
\texttt{FreeFem++} \cite{hecht2012new}.
We investigate the problems in the space-time domain $Q=\Omega\times(0,T)$, where $\Omega=[-8,8]^d$ with $d=1,2$ being the space dimension. The discretization parameters are listed in ~\Cref{tab:mesh_params_all}.
Unless stated otherwise, the adaptive algorithm (\Cref{alg:Lscheme-adapt}) uses the parameters $L = 1.0$, $\theta = 0.8$, $\varepsilon_{\mathrm{adapt}} = 0.1$, $r_{\mathrm{ref}} = 4$, $\delta_{\mathrm{lin}} = 0.2$, and $\texttt{Tol}=10^{-8}$. The maximum number of iterations is chosen as $i_{\max}\in\{50,100\}$ depending on the problem. Similarly, the maximum refinement level $\ell_{\max}\in\{8,\ldots,15\}$ varies across the considered test cases. For the regularized Newton scheme, the regularization parameter is fixed to $\varepsilon = 0.1$ throughout all simulations in this section.



% Adaptive strategies (Adapt-L and Adapt-M) refine a coarse mesh locally using the a posteriori estimator and Dörfler marking (Algorithm~\ref{alg:Lscheme-adapt}). The $L$-scheme (L) is computed on a fixed uniform mesh, while Uni-L and Uni-M use global refinement at each level with $h_{\ell+1}=h_\ell/2$, $\Delta t_{\ell+1}=\Delta t_\ell/2$.
% All simulations are performed on the space--time domain and consider two model problems: the porous medium equation (PME) and a biofilm growth model.
%The same profile is used for the biofilm model to ensure consistency across test cases.

\iffalse
\begin{table}[h!]
\centering
% ------- 1D TABLE -------
\begin{minipage}{0.95\textwidth}
\centering
\begin{tabular}{lccccccc}
\toprule
\textbf{Run} & \textbf{Mesh Type} & $n_{\mathrm{space}}$ & $n_{\mathrm{time}}$
& $m$ & $x_{\min}$ & $x_{\max}$ & Time Interval \\
\midrule
1D-PME / 1D-Biofilm & Uniform  & 160 & 20 & 2 & $-8$ & $8$ & $[0,1]$ \\
1D-PME / 1D-Biofilm & Adaptive & 16  & 2  & 2 & $-8$ & $8$ & $[0,1]$ \\
\midrule
1D-PME              & Uniform  & 160 & 20 & 4 & $-8$ & $8$ & $[0,1]$ \\
1D-PME              & Adaptive & 16  & 2  & 4 & $-8$ & $8$ & $[0,1]$ \\
\midrule
1D-PME              & Uniform  & 160 & 20 & 6 & $-8$ & $8$ & $[0,1]$ \\
1D-PME              & Adaptive & 16  & 2  & 6 & $-8$ & $8$ & $[0,1]$ \\
\bottomrule
\end{tabular}
\end{minipage}

\vspace{1em}

% ------- 2D TABLE -------
\begin{minipage}{0.95\textwidth}
\centering
\begin{tabular}{lcccccccc}
\toprule
\textbf{Run} & \textbf{Mesh Type} & $n_{\mathrm{spaceX}}$ & $n_{\mathrm{spaceY}}$
& $n_{\mathrm{time}}$ & $m$ & $x_{\min}/y_{\min}$ & $x_{\max}/y_{\max}$ & Time \\
\midrule
2D-PME      & Uniform  & 32 & 32 & 4 & 2 & $-8$ & $8$ & $[0,1]$ \\
2D-Biofilm  & Uniform  & 32 & 32 & 4 & 2 & $-8$ & $8$ & $[0,1]$ \\
2D-PME      & Adaptive & 16 & 16 & 4 & 2 & $-8$ & $8$ & $[0,1]$ \\
2D-Biofilm  & Adaptive & 16 & 16 & 2 & 2 & $-8$ & $8$ & $[0,1]$ \\
\midrule
2D-PME      & Uniform  & 32 & 32 & 4 & 4 & $-8$ & $8$ & $[0,1]$ \\
2D-PME      & Adaptive & 16 & 16 & 4 & 4 & $-8$ & $8$ & $[0,1]$ \\
\midrule
2D-PME      & Uniform  & 32 & 32 & 4 & 6 & $-8$ & $8$ & $[0,1]$ \\
2D-PME      & Adaptive & 16 & 16 & 4 & 6 & $-8$ & $8$ & $[0,1]$ \\
\bottomrule
\end{tabular}
\end{minipage}
\caption{Meshes and model parameters used in all 1D and 2D numerical experiments.}
\label{tab:mesh_params_all}
\end{table}
\fi
%==============================================================
\begin{table}[h!]
\centering
\caption{ Mesh and model parameters for space-time experiments on the domain $[-8,8]$ (or $[-8,8]^2$ in 2D) and time interval $[0,1]$.}

\resizebox{\textwidth}{!}{%
\begin{tabular}{llccc}
\toprule
 & \textbf{Mesh} & $n_{\text{space}}$ & $n_{\text{time}}$ & $m$ \\
\midrule

\multirow{3}{*}{1D}
& Adaptive ref. & 16 & 2 
& PME: $2,4,6$; Others: $2$ \\

& Uniform ref. & 16 & 2 
& PME: $2,4,6$; Others: $2$ \\

& Fixed mesh 
& 160 (PME), 80 Others 
& 20 (PME), 10 Others 
& PME: $2,4,6$; Others: $2$\\

\midrule

\multirow{2}{*}{2D}
& Adaptive ref. 
& $16 \times 16$ 
& 4 
& PME: $2,4,6$; Others: $2$ \\

& Fixed mesh 
& $32 \times 32$ 
& 4 
& PME: $2,4,6$; Others: $2$ \\

\bottomrule
\end{tabular}%
}

\label{tab:mesh_params_all}
\end{table}
%========================================

\iffalse
\begin{table}[h!]
\centering
\caption{Numerical parameters for 1D PME on an adaptive mesh.}
\label{tab:1DPME_16_T4}
\begin{minipage}{0.95\textwidth}
\centering
\begin{tabular}{lccccccc}
\toprule
\textbf{Run} & \textbf{Mesh} & $n_{\mathrm{space}}$ & $n_{\mathrm{time}}$
& $m$ & Domain & Time Interval \\
\midrule
1D-PME & Adaptive & 16 & 4 & 6 & $[-8,8]$ & $[0,4]$ \\
\bottomrule
\end{tabular}
\end{minipage}
\end{table}
\fi
\vspace{2mm}
In $d$ spatial dimensions ($d=1$ or $d=2$), we employ the Barenblatt solution \cite{vazquez2007porous} at $t=0$ as our initial condition for all test cases. For $m>1$ and for $\nu = \frac{1}{(m-1 + \frac{2}{d})}$, this solution is given by
\begin{align}\label{eq:Barenblatt}
  u_{\,\!_{\rm BB}}(\bm{x},t) = (1 + t)^{-\nu} \left[\max\left(\gamma - \frac{\nu(m-1)|\bm{x}|^2}{2dm(t+1)^{2\nu/d}}, 0\right)\right]^{1/(m-1)}.
\end{align}
Homogeneous Dirichlet boundary conditions are imposed on $\Gamma_T=\partial\Omega\times(0,T]$.
For the porous medium equation with exponents $m$ (see \eqref{eq:PME}), $ u_{\,\!_{\rm BB}}$ is also the exact solution. This enables us to compute the exact error for this case. 

\subsection{Porous medium equation (PME)}\label{sec:porous_medium_equation}
We consider the porous medium equation
\begin{equation}\label{eq:PME}
\partial_t u = \Delta u^m \quad \text{ for }\quad  m>1
\end{equation}
which corresponds to $\Phi(u)=u^m$. It shows degeneracy at $u = 0$ since $\Phi'$ vanishes. Following the $b$--$B$ decomposition of \Cref{lemma:b_B_construction} (introduced in \cite{javed2025robust}), we reformulate the PME as $\partial_t b(s) - \Delta B(s) = 0$. Let $u^*>0$ satisfy $\Phi'(u^*)=1$, that is, $u^*=\left(\frac{1}{m}\right)^{\frac{1}{m-1}}$. Then the functions $b$ and $B$ admit the explicit representations
\begin{align}\label{eq:bB_PME}
b(s)=
\begin{cases}
0, & s<0,\\[4pt]
s, & 0\le s\le u^*,\\[4pt]
\bigl(s-u^*+(u^*)^m\bigr)^{1/m}, & s>u^*,
\end{cases}
\qquad
B(s)=
\begin{cases}
s, & s<0,\\[4pt]
s^m, & 0\le s\le u^*,\\[4pt]
s-u^*+(u^*)^m, & s>u^*.
\end{cases}
\end{align}
 Since $u_{\,\!_{\rm BB}}$, the exact solution of the porous medium equation is known, we can determine $s$ exactly, and use it to assess the quality of the scheme.
 %with $\Phi(u)=u^m$, this choice allows us to validate the numerical implementation.
\subsubsection{1D PME}\label{sec:1DPME_1DST}

%\subsection{Numerical setup}\label{subsec:num-setup}
%We investigate the above problem in the space-time domain $\Omega\times(0,T)$, where $\Omega\subset\mathbb{R}^d$ with $d=1,2$ is a bounded spatial domain (e.g., $\Omega=[x_{\min},x_{\max}]$ in one dimension or  $\Omega = [x_{\min},x_{\max}]\times[y_{\min},y_{\max}]$ in two dimensions), with parameters listed in ~\cref{tab:mesh_params_all}. 
%Uniform runs are typically performed on fine meshes; however, in some experiments (e.g., \Cref{fig:uniform_adaptive_meshes_halved}), we also consider uniform refinements starting from a coarse mesh, where the mesh size and time step are successively halved at each refinement level. Whereas adaptive runs start from coarse meshes and activate the vertex–patch reconstructions when the linearization error is sufficiently small (see Algorithm~\ref{alg:Lscheme-adapt}). 
%========================================

\begin{figure}[h]
\centering
% --- Initial 1D space--time mesh ---
%\begin{subfigure}{0.49\textwidth}
 % \centering
 % \includegraphics[width=\linewidth]{Figures/Initial_mesh.png.jpeg}
  %\caption{Initial uniform space-time mesh in $(x,t)$.}
  %\label{fig:init_mesh_1d}
%\end{subfigure}
%\hfill
% --- Final 1D space--time mesh ---
%\begin{subfigure}{0.49\textwidth}
%  \centering
%  \includegraphics[width=\linewidth]{Figures/Final_mesh_1D_M6_10_refinemnts.jpeg}
%  \caption{Final adapted space-time mesh in $(x,t)$.}
%  \label{fig:final_mesh_1d}
%\end{subfigure}

%\vspace{1em}

% --- Solution on initial 1D mesh ---
%\begin{subfigure}{0.49\textwidth}
%  \centering
%  \includegraphics[width=1.05\linewidth]{Figures/Initail_Solution_M6_test.jpeg}
%  \caption{Initial condition on space-time mesh.}
%  \label{fig:sol_init_1d}
%\end{subfigure}
%\hfill
% --- Solution on adapted 1D mesh ---
%\begin{subfigure}{0.49\textwidth}
%  \centering
%  \includegraphics[width=1.05\linewidth]{Figures/Solution_Final_Mesh_M6_10_refinements.jpeg}
%  \caption{Solution on the adapted space-time mesh.}
%  \label{fig:sol_final_1d}
%\end{subfigure}
%\hfill
%\begin{subfigure}{0.65\textwidth}
%  \centering
  \includegraphics[width=0.75\linewidth]{Figures/PME_u_3D.eps}
 % \caption{[\Cref{sec:1DPME_1DST}] 3D view of the numerical solution $u = b(s)$ on the final adapted space-time mesh.}
 % \label{fig:sol_3d_PME}
%\caption{1D PME in space and 1D in time: mesh refinement and corresponding numerical solutions. \textbf{Top:} initial and final space-time meshes in $(x,t)$. \textbf{Bottom:} (left) initial condition on space-time mesh. (right) solution after mesh adaptivity. }
\caption{[Section~\ref{sec:1DPME_1DST}, $m=6$, $\gamma=1.0$, $T=4$] Three-dimensional view of the numerical solution $u = b(s)$ of the porous medium equation on the final adapted space-time mesh.}
\label{fig:mesh_refinement_summary_1d}
\end{figure}
Figure~\ref{fig:mesh_refinement_summary_1d} illustrates the behaviour of the solution for the 1D PME evaluated using the adaptive refinement strategy. It is seen that the estimator $\eta^i_K$ of \eqref{eq:etaK} drives refinement in regions where the spatial gradients and temporal variations are high, i.e., at the free boundaries and close to $t=0$. 
%The three-dimensional view of the numerical solution on the adapted mesh confirms that the free boundary, where the degeneracy of the equation produces sharp fronts, dominates the residual contributions and is automatically detected by the equilibrated flux estimator, while regions where the solution remains smooth or close to zero are left essentially unrefined.
This is further confirmed in Figure~\ref{fig:true_error_vs_estimator}, which compares the element-wise true error (Left) and the effectivity indices (Right) introduced in \eqref{eq:effectivity} over successive adaptive refinement levels. It is seen that errors are concentrated near the free boundary with much lower concentration at the degenerate $s=0$ region. The refinement happens accordingly, localized to the free-boundary and the initial condition regions. The errors decrease systematically as the mesh is refined. The effectivity indices are of the order of 1 in the region $\{s>0\}$ including the free-boundaries. However, in the region $s=0$ they become very small since the errors $\|e\|_K$ are practically zero. 


%, as shown in~\Cref{Figure:1DPME_m=2_dof,Figure:1DPME_m=4_dof,Figure:1DPME_m=6_dof}.
%===================================================
%\begin{figure}[h!]
 % \centering
 % \begin{subfigure}[t]{0.33\textwidth}
 %       \includegraphics[width=\textwidth]{Figures/region_error_comparison.eps}
 %         \end{subfigure}
%\begin{subfigure}[t]{0.33\textwidth}
%  \centering
%  \includegraphics[width=\textwidth]{Figures/Iterations_vs_RefinementLevel.eps}
%\end{subfigure}
%\hfill
%\begin{subfigure}[t]{0.32\textwidth}
%  \centering
%  \includegraphics[width=\textwidth]{Figures/refined_region_error_reduction.eps}
%\end{subfigure}
%\caption{[\Cref{tab:1DPME_16_T4,sec:1DPME_1DST}] Performance of the adaptive space-time refinement strategy.
 %\textbf{Left} Evolution of the a posteriori error estimator $\eta(K)$, evaluated separately on two regions of the mesh: the \emph{refined region} (elements selected for refinement) and the \emph{coarse region} (non-selected elements). Both quantities are shown before and after each refinement step. \textbf{Middle} Number of solver iterations required between successive adaptive refinement levels. \textbf{Right} Reduction factor of the a posteriori error in the refined region.}
%\label{fig:region-error-redistribution}
%\end{figure}
\noindent
%Figure \ref{fig:region-error-redistribution} (Left) shows the evolution of the a posteriori error indicator $\eta(K)$ evaluated separately on the elements selected for refinement and on the remaining coarse elements, both before and after mesh adaptation. The solid blue curve corresponds to the error in the selected refinement region, evaluated on the mesh prior to refinement. These elements are identified as carrying a relatively large portion of the total error and are therefore marked for refinement. The solid red curve shows the error in the non-selected region, also evaluated before refinement. At this stage, both regions are subsets of the same current mesh and are therefore evaluated at identical resolution. The impact of the adaptive strategy becomes apparent only after refinement, as illustrated by the dashed curves, where the estimator is recomputed on the updated mesh. The dashed green curve represents the error in the refined region after local mesh refinement. In the selected refinement region, the error is evaluated over the same physical region, projected onto the adapted mesh. A noticeable reduction of the error is observed, demonstrating the effectiveness of the local refinement in the targeted high-error region. The dashed red curve corresponds to the coarse region after refinement. Although these elements are not directly selected for refinement, the mesh in this region may change a bit due to mesh conformity, smoothing, and indirect coupling with refined neighbouring elements. Consequently, a moderate reduction in error can also be observed. However, this reduction remains smaller than that achieved in the refined region, confirming that the refinement strategy acts primarily where intended.
%===================================================================

%\begin{figure}[h!]
%\centering
%\begin{subfigure}[t]{0.4\textwidth}
 % \centering
  %\includegraphics[width=\textwidth]{Figures/Iterations_vs_RefinementLevel.eps}
  %\caption{Number of iterations between successive adaptive refinement levels.}
%  \label{fig:iter_per_level}
%\end{subfigure}
%\hfill
%\begin{subfigure}[t]{0.4\textwidth}
 % \centering
 % \includegraphics[width=\textwidth]{Figures/refined_region_error_reduction.eps}
 % \caption{Reduction factor of the a posteriori error in the region selected for refinement.}
  %\label{fig:error_reduction}
%\end{subfigure}
%\caption{[\Cref{tab:1DPME_16_T4}] Performance of the adaptive space-time refinement strategy.
%\textbf{Left}: number of solver iterations required between successive adaptive refinement levels. \textbf{Right}: reduction factor of the a posteriori error in the refined region, defined as
%$\eta_{\mathrm{refined,after}} / \eta_{\mathrm{refined,before}}$.}
%\label{fig:refined-reduction}
%\end{figure}

%======================================================
%\begin{figure}[h]
 % \centering
  %\includegraphics[width=0.75\textwidth]{Figures/refined_region_error_reduction.eps}
 % \caption{Reduction factor of the a posteriori error in the refined region.
%  The quantity shown is
 % \(
  %  \eta_{\mathrm{refined,after}} / \eta_{\mathrm{refined,before}}
  %\),
 % where the error is evaluated over the same physical region before and after mesh refinement.
 % Values significantly below one indicate effective local error reduction due to refinement.}
 % \label{fig:refined-reduction}
%\end{figure}
\noindent
%Figure~\ref{fig:region-error-redistribution} (Middle) shows the number of iterations performed between successive adaptive refinement levels. The initial increase in the number of iterations is due to early refinements uncovering sharp and degenerate solution features that are poorly resolved on coarse meshes. Once these features are adequately captured by the adaptive mesh, the nonlinear solver benefits from better initial guesses,  leading to a reduced number of iterations in later refinement stages. Figure~\ref{fig:region-error-redistribution} (right) shows the reduction factor of the a posteriori error in the region selected for refinement. Although the mesh changes after refinement, the refined region is transferred to the new mesh via projection of the characteristic function, so the error is evaluated over the same physical region in space-time. The consistently small values of the ratio confirm that local mesh refinement effectively reduces the error in the targeted region.

%===================================================

\iffalse
\begin{figure}[h]
\centering
\begin{subfigure}{0.42\textwidth}
  \centering
  \includegraphics[width=\textwidth]{Figures/RefinemnetRegion_1.png}
  \caption{Marked elements — level 0}
\end{subfigure}
\hfill
\begin{subfigure}{0.42\textwidth}
  \centering
  \includegraphics[width=\textwidth]{Figures/AfterRefinemnetRegion_1.png}
  \caption{Mesh after refinement — level 1}
\end{subfigure}
\hfill
\begin{subfigure}{0.42\textwidth}
  \centering
  \includegraphics[width=\textwidth]{Figures/RefinemnetRegion_2.png}
  \caption{Marked elements — level 1}
\end{subfigure}
\hfill
\begin{subfigure}{0.42\textwidth}
  \centering
  \includegraphics[width=\textwidth]{Figures/AfterRefinemnetRegion_2.png}
  \caption{Mesh after refinement — level 2}
\end{subfigure}
\hfill
\begin{subfigure}{0.42\textwidth}
  \centering
  \includegraphics[width=\textwidth]{Figures/RefinemnetRegion_3.png}
  \caption{Marked elements — level 2}
\end{subfigure}
\hfill
\begin{subfigure}{0.42\textwidth}
  \centering
  \includegraphics[width=\textwidth]{Figures/AfterRefinemnetRegion_3.png}
  \caption{Mesh after refinement — level 3}
\end{subfigure}
\hfill
\begin{subfigure}{0.42\textwidth}
  \centering
  \includegraphics[width=\textwidth]{Figures/RefinemnetRegion_4.png}
  \caption{Marked elements — level 3}
\end{subfigure}
\hfill
\begin{subfigure}{0.42\textwidth}
  \centering
  \includegraphics[width=\textwidth]{Figures/AfterRefinemnetRegion_4.png}
  \caption{Mesh after refinement — level 4}
\end{subfigure}
\hfill
\begin{subfigure}{0.42\textwidth}
  \centering
  \includegraphics[width=\textwidth]{Figures/RefinemnetRegion_5.png}
  \caption{Marked elements — level 4}
\end{subfigure}
\hfill
\begin{subfigure}{0.42\textwidth}
  \centering
  \includegraphics[width=\textwidth]{Figures/AfterRefinemnetRegion_5.png}
  \caption{Mesh after refinement — level 5}
\end{subfigure}
\hfill
\begin{subfigure}{0.42\textwidth}
  \centering
  \includegraphics[width=\textwidth]{Figures/RefinemnetRegion_6.png}
  \caption{Marked elements — level 5}
\end{subfigure}
\hfill
\begin{subfigure}{0.42\textwidth}
  \centering
  \includegraphics[width=\textwidth]{Figures/AfterRefinemnetRegion_6.png}
  \caption{Mesh after refinement — level 6}
\end{subfigure}
\hfill
\begin{subfigure}{0.42\textwidth}
  \centering
  \includegraphics[width=\textwidth]{Figures/RefinemnetRegion_7.png}
  \caption{Marked elements — level 6}
\end{subfigure}
\hfill
\begin{subfigure}{0.42\textwidth}
  \centering
  \includegraphics[width=\textwidth]{Figures/AfterRefinemnetRegion_7.png}
  \caption{Mesh after refinement — level 7}
\end{subfigure}
%\hfill
%\begin{subfigure}{0.42\textwidth}
%  \centering
%  \includegraphics[width=\textwidth]{Figures/RefinemnetRegion_8.png}
%  \caption{Marked elements — level 7}
%\end{subfigure}
%\hfill
%\begin{subfigure}{0.42\textwidth}
%  \centering
%  \includegraphics[width=\textwidth]{Figures/AfterRefinemnetRegion_8.png}
%  \caption{Mesh after refinement — level 8}
%\end{subfigure}
%\hfill
%\begin{subfigure}{0.42\textwidth}
%  \centering
%  \includegraphics[width=\textwidth]{Figures/RefinemnetRegion_9.png}
%  \caption{Marked elements — level 8}
%\end{subfigure}
%\hfill
%\begin{subfigure}{0.42\textwidth}
%  \centering
%  \includegraphics[width=\textwidth]{Figures/AfterRefinemnetRegion_9.png}
 % \caption{Mesh after refinement — level 9}
%\end{subfigure}
%\hfill
%\begin{subfigure}{0.42\textwidth}
%  \centering
 % \includegraphics[width=\textwidth]{Figures/RefinemnetRegion_10.png}
  %\caption{Marked elements — level 9}
%\end{subfigure}
%\hfill
%\begin{subfigure}{0.42\textwidth}
%  \centering
%  \includegraphics[width=\textwidth]{Figures/AfterRefinemnetRegion_10.png}
%  \caption{Mesh after refinement — level 10}
%\end{subfigure}

\caption{[\Cref{tab:mesh_params_all,,sec:1DPME_1DST}, $T=4$, $m=6$] 
Adaptive refinement process over successive levels driven by the a posteriori local indicators $\eta^i_K$ defined in \eqref{eq:etaK}. \textbf{Left column:} current mesh with elements marked for refinement (blue) satisfying the D\"{o}rfler bulk criterion with $\theta=0.8$, activated when $\eta^i_{\mathrm{lin}} \leq \delta_{\mathrm{lin}}\, \eta^i_{\mathrm{disc}}$ with $\delta_{\mathrm{lin}}=0.2$. \textbf{Right column:} resulting mesh after refinement at the next level using refinement factor $4$. Elements not selected (red) have smaller local indicators $\eta^i_K$ and may be marked in later levels.}
\label{fig:refinement-regions}
\end{figure}
\noindent
\fi
%================================================================
\begin{center}
\setlength{\tabcolsep}{2pt}
\setlength{\LTpre}{0pt}
\setlength{\LTpost}{0pt}
\begin{longtable}{cc}

% ---- header row (column titles) ----------------------
 \multicolumn{1}{c}{\textbf{$\log_{10}(\|e\|_K)$ -- Element-wise true error}} &
\multicolumn{1}{c}{\textbf{$\log_{10}(\textit{Eff. Index})$ -- Effectivity index}}\\[4pt]

% -------------------------------------------------------
% Row 0
% -------------------------------------------------------

\includegraphics[width=0.49\textwidth]{Figures/true_errror_1.PNG} &\includegraphics[width=0.49\textwidth]{Figures/logEffIdx_level_1.png} \\
{\small $\log_{10}(\|e\|_K)$ --- level 0} 
&{\small $\textit{Eff. Index} :={\|e\|_K}/\eta_K $,\; level 0} %
\\[6pt]
% -------------------------------------------------------
% Row 1
% -------------------------------------------------------
\includegraphics[width=0.49\textwidth]{Figures/true_errror_2.PNG}
& \includegraphics[width=0.49\textwidth]{Figures/logEffIdx_level_2.png}\\
{\small $\log_{10}(\|e\|_K)$ --- level 1}
&{\small $\textit{Eff. Index} :={\|e\|_K}/\eta_K $,\; level 1}\\[6pt]

% -------------------------------------------------------
% Row 2
% ------------------------------------------------------
\includegraphics[width=0.49\textwidth]{Figures/true_errror_3.PNG} &\includegraphics[width=0.49\textwidth]{Figures/logEffIdx_level_3.png} \\
{\small $\log_{10}(\|e\|_K)$ --- level 2} &
{\small $\textit{Eff. Index} :={\|e\|_K}/\eta_K $,\; level 2}\\[6pt]

% -------------------------------------------------------
% Row 3
% -------------------------------------------------------

\includegraphics[width=0.49\textwidth]{Figures/true_errror_4.PNG} 
& \includegraphics[width=0.49\textwidth]{Figures/logEffIdx_level_4.png}\\
{\small  $\log_{10}(\|e\|_K)$ --- level 3} 
&{\small $\textit{Eff. Index} :={\|e\|_K}/\eta_K $,\; level 3}
\\[6pt]
% -------------------------------------------------------
% Row 4
% -------------------------------------------------------
\includegraphics[width=0.49\textwidth]{Figures/true_errror_5.PNG} &
\includegraphics[width=0.49\textwidth]{Figures/logEffIdx_level_5.png} \\
{\small $\log_{10}(\|e\|_K)$ --- level 4} &
{\small $\textit{Eff. Index} :={\|e\|_K}/\eta_K $,\; level 4}
\\[6pt]

% -------------------------------------------------------
% Row 5
% -------------------------------------------------------
% \includegraphics[width=0.49\textwidth]{Figures/logEffIdx_level_6.png} &
% \includegraphics[width=0.49\textwidth]{Figures/true_errror_6.PNG} \\
% {\small (xi) $\textit{Eff. Index} :={\|e\|_K}/\eta_K $,\; level 5} &
% {\small (xii) $\log_{10}(\|e\|_K)$ --- level 5} \\[6pt]

% -------------------------------------------------------
% Row 6
% -------------------------------------------------------
\includegraphics[width=0.49\textwidth]{Figures/true_errror_7.PNG} &
\includegraphics[width=0.49\textwidth]{Figures/logEffIdx_level_7.png} 
 \\
{\small $\log_{10}(\|e\|_K)$ --- level 6} 
&{\small $\textit{Eff. Index} :={\|e\|_K}/\eta_K $,\; level 6}\\[6pt]

% -------------------------------------------------------
% Row 7
% -------------------------------------------------------
% \includegraphics[width=0.49\textwidth]{Figures/logEffIdx_level_8.png} &
% \includegraphics[width=0.49\textwidth]{Figures/true_errror_8.PNG} \\
% {\small (xv) $\textit{Eff. Index} :={\|e\|_K}/\eta_K $,\; level 7} &
% {\small (xvi) $\log_{10}(\|e\|_K)$ --- level 7} \\[6pt]

% % -------------------------------------------------------
% % Row 8
% % -------------------------------------------------------
% \includegraphics[width=0.49\textwidth]{Figures/logEffIdx_level_9.png} &
% \includegraphics[width=0.49\textwidth]{Figures/true_errror_9.PNG} \\
% {\small (xvii) $\textit{Eff. Index} :={\|e\|_K}/\eta_K $,\; level 8} &
% {\small (xviii) $\log_{10}(\|e\|_K)$ --- level 8} \\[6pt]

% -------------------------------------------------------
% Row 9
% -------------------------------------------------------
\includegraphics[width=0.49\textwidth]{Figures/true_errror_10.PNG} &
\includegraphics[width=0.49\textwidth]{Figures/logEffIdx_level_10.png} \\
{\small $\log_{10}(\|e\|_K)$ --- level 9} &
{\small $\textit{Eff. Index} :={\|e\|_K}/\eta_K $,\; level 9}
\\[6pt]
\end{longtable}
\captionof{figure}{[Section~\ref{sec:1DPME_1DST}, $m=6$, $\gamma=1.0$, $T=4$] \textbf{Left column:} 
    Element-wise error $\|e\|_K$ defined in \eqref{eq:effectivity} in $\log_{10}$-scale over multiple levels of refinement. \textbf{Right column:} Element-wise effectivity indices in $\log_{10}$-scale. }

\label{fig:true_error_vs_estimator}
\end{center}

\Cref{Figure:1DPME_m=2_dof} shows the error decays with the cost $\mathrm{Cost}(i,\ell)$ defined in \eqref{eq:cost} comparing the adaptive, uniform, and fixed mesh strategies for the $L$-scheme with $m=2$. Both adaptive and uniform strategies reach lower error levels, and the adaptive strategy performs and scales more favorably, albeit slightly for $m=2$ (better performance is observed for higher $m$ values depicted later). \Cref{Figure:1DPME_m=2_iterationno} shows the error decay in different metrics with respect to total iteration number from \eqref{eq:cost}. The adaptive strategy requires more iterations due to the projection process after each refinement step which introduces additional errors. We will see similar behaviour later on. \Cref{Figure:1DPME_Newton_m=2_dof} shows the error decay for the Adapt-$N_{\mathrm{reg}}$, Uni-$N_{\mathrm{reg}}$, and $N_{\mathrm{reg}}$ (fixed mesh) schemes. Similar to \Cref{fig:iter-combined}, these schemes exhibit significantly more oscillations compared to the \(L\)-schemes, even in the $L^2$-norm. Even for uniform refinements, the oscillations persist, as opposed to the \(L\)-schemes, see \Cref{fig:compasion-Uni-Adapt} (Left).  Thus, as concluded also in \cite{mitra2023guaranteed}, the faster convergence of the linearization scheme does not necessarily translate to faster overall convergence in terms of total error. Henceforth, we will focus only on the metric $\|b(s_{h\tau}^{(i)}) - u_{\text{exact}}\|_{L^2(Q)}$ since \Cref{theo:converegence} predicts this metric to vanish as the iterations progress, and also since \Cref{fig:iter-combined,Figure:1DPME_m=2_dof,Figure:1DPME_m=2_iterationno,Figure:1DPME_Newton_m=2_dof} show the errors to oscillate less in this measure.  

\begin{figure}[h]
    \centering
    \begin{subfigure}[b]{0.49\textwidth}
        \includegraphics[width=\textwidth]
{Figures/EnergyError_vs_DoFs_adaptive_vs_uniform_BH1_dx_Sh_m2_Newcode.eps}       %{Figures/EnergyError_vs_DoFs_adaptive_vs_uniform_BH1_dx_Sh_m2.eps}
    \end{subfigure}
    \hfill
    \begin{subfigure}[b]{0.49\textwidth}
        \includegraphics[width=\textwidth]
{Figures/EnergyError_vs_DoFs_adaptive_vs_uniform_bbL2_dx_Sh_m2_Newcode.eps}
%{Figures/EnergyError_vs_DoFs_adaptive_vs_uniform_bbL2_dx_Sh_m2.eps}
    \end{subfigure}
    \centering
    \caption{[\Cref{tab:mesh_params_all}, Section \ref{sec:1DPME_1DST}, $m=2$] Comparison of errors versus $\mathrm{Cost}(i,\ell)$ \eqref{eq:cost} for fixed, uniform, and adaptive $L$-schemes. 
\textbf{Left}: Error in terms of $\|\nabla\!\left(B(s_{h\tau}^{(i)}) - B(s)\right)\|_{L^2(Q)}$. 
\textbf{Right}: Error in terms of $\|b(s_{h\tau}^{(i)}) - u_{\text{exact}}\|_{L^2(Q)}$. 
%Linearization error: $\eta^i_{\mathrm{lin}}=\| \nabla (s_h^{(i)} - s_h^{(i-1)}) \bigr\|_{L^2(Q)}$.
}
\label{Figure:1DPME_m=2_dof}
\end{figure}

%======================================================
\begin{figure}[h]
    \centering
    \begin{subfigure}[b]{0.49\textwidth}
        \includegraphics[width=\textwidth]       {Figures/EnergyError_vs_Iteration_adaptive_vs_uniform_BH1_dx_Sh_m2_Newcode.eps}       
        %{Figures/EnergyError_vs_Iteration_adaptive_vs_uniform_BH1_dx_Sh_m2.eps}
    \end{subfigure}
    \hfill
    \begin{subfigure}[b]{0.49\textwidth}
        \includegraphics[width=\textwidth]
{Figures/EnergyError_vs_Iteration_adaptive_vs_uniform_bbL2_dx_Sh_m2_Newcode.eps}
%{Figures/EnergyError_vs_Iteration_adaptive_vs_uniform_bbL2_dx_Sh_m2.eps}

    \end{subfigure}
    \centering
     \caption{[\Cref{tab:mesh_params_all}, Section \ref{sec:1DPME_1DST}, $m=2$] Comparison of errors versus $\mathrm{IterNo}(i,\ell)$ \eqref{eq:cost} for fixed, uniform and adaptive $L$-schemes.
\textbf{Left}: Error in terms of $\|\nabla\!\left(B(s_{h\tau}^{(i)}) - B(s)\right)\|_{L^2(Q)}$. 
\textbf{Right}: Error in terms of $\|b(s_{h\tau}^{(i)}) - u_{\text{exact}}\|_{L^2(Q)}$. 
%Linearization error: $\eta^i_{\mathrm{lin}}=\| \nabla (s_h^{(i)} - s_h^{(i-1)}) \bigr\|_{L^2(Q)}$.
}
\label{Figure:1DPME_m=2_iterationno}
\end{figure}

%===============================================
%====================================================
\iffalse
\begin{figure}[h]
    \centering
    \begin{subfigure}[b]{0.49\textwidth}
        \includegraphics[width=\textwidth]
{Figures/EnergyError_vs_DoFs_adaptive_vs_uniform_BH1_dx_Sh_new.eps}       %
    \end{subfigure}
    \hfill
    \begin{subfigure}[b]{0.49\textwidth}
        \includegraphics[width=\textwidth]
{Figures/EnergyError_vs_DoFs_adaptive_vs_uniform_bbL2_dx_Sh_m2_new.eps}
    \end{subfigure}
    \centering
    \caption{[\Cref{tab:mesh_params_all}, Section \ref{sec:1DPME_1DST}, $m=2$, $h_{k}^2$] Comparison of errors versus $\mathrm{Cost}(i,\ell)$ for fixed, uniform and adaptive space-time meshes.
\textbf{Right}: Error in terms of $\|b(s_{h\tau}^{(i)}) - u_{\text{exact}}\|_{L^2(Q)}$. 
%Linearization error: $\eta^i_{\mathrm{lin}}=\| \nabla (s_h^{(i)} - s_h^{(i-1)}) \bigr\|_{L^2(Q)}$.
}
\end{figure}
%===============================================
\begin{figure}[h]
    \centering
    \begin{subfigure}[b]{0.49\textwidth}
        \includegraphics[width=\textwidth]       {Figures/EnergyError_vs_Iteration_adaptive_vs_uniform_BH1_dx_Sh_new.eps}       
    \end{subfigure}
    \hfill
    \begin{subfigure}[b]{0.49\textwidth}
      \includegraphics[width=\textwidth]{Figures/EnergyError_vs_Iteration_adaptive_vs_uniform_bbL2_dx_Sh_m2_ne.eps}
    \end{subfigure}
    \centering    \caption{[\Cref{tab:mesh_params_all}, Section \ref{sec:1DPME_1DST}, $m=2$, $h_{k}^2$] Comparison of energy errors versus $\mathrm{Cost}(i,\ell)$ for fixed, uniform and adaptive space-time meshes.
\textbf{Left}:  $\|\nabla_x\!\left(B(s_{h\tau}^{(i)}) - B(s)\right)\|_{L^2(Q)}$. 
\textbf{Right}:  $\|b(s_{h\tau}^{(i)}) - u_{\text{exact}}\|_{L^2(Q)}$. 
Linearization error: $\eta^i_{\mathrm{lin}}=\| \nabla (s_h^{(i)} - s_h^{(i-1)}) \bigr\|_{L^2(Q)}$.}
\end{figure}
\fi
%=========================================
%\begin{figure}[h]
 %  \centering        \includegraphics[width=0.5\textwidth]   {Figures/EnergyError_vs_DoFs_adaptive_vs_uniform_BL2.eps}
 %  \end{figure}
%===============================================   
   
%\begin{itemize}
 %  \item \textbf{Linearization Error:} 
 %  \[ 
 %  \bigl\|(s_h^{(i+1)} - s_h^{(i)}) \bigr\|_{L^2(Q)},
%   \qquad
%\]
%\end{itemize}
%\begin{figure}[h]
 %   \centering
  %  \begin{subfigure}[b]{0.49\textwidth}
   %     \includegraphics[width=\textwidth]{Figures/EnergyError_vs_DoFs_adaptive_vs_uniform_BH1_Sh.eps}
    %\end{subfigure}
    %\hfill
    %\begin{subfigure}[b]{0.49\textwidth}
     %   \includegraphics[width=\textwidth]{Figures/EnergyError_vs_DoFs_adaptive_vs_uniform_BL2_Sh.eps}
    %\end{subfigure}
    %\caption{Comparison between the uniform and adaptive space--time meshes.
  %  The adaptive mesh concentrates refinement in regions of high error,
   % leading to improved accuracy with less degrees of freedom.
    %Linearization error: $\| s_h^{(i+1)} - s_h^{(i)} \|_{L^2(Q)}$.}
  %  \label{fig:uniform_vs_adaptive}
%\end{figure}
%\begin{figure}[h]
  %  \centering        \includegraphics[width=0.7\textwidth]{Figures/EnergyError_vs_DoFs_adaptive_vs_uniform_bbL2_Sh.eps}
  %  \end{figure}
 %=========================================================
\begin{figure}[h]
    \centering
    \begin{subfigure}[b]{0.49\textwidth}
        \includegraphics[width=\textwidth]{Figures/EnergyError_vs_DoFs_adaptive_vs_newton_M2_BH1_Level.eps}
        %{Figures/EnergyError_vs_DoFs_adaptive_vs_newton_M2_BH1.eps}
    \end{subfigure}
    \hfill
    \begin{subfigure}[b]{0.49\textwidth}
        \includegraphics[width=\textwidth]
{Figures/EnergyError_vs_DoFs_adaptive_vs_newton_M2_bbL2_Level.eps}
%{Figures/EnergyError_vs_DoFs_adaptive_vs_newton_M2_bbL2.eps}
    \end{subfigure}
    \centering
    \caption{[\Cref{tab:mesh_params_all}, Section \ref{sec:1DPME_1DST}, $m=2$, $\delta_{\mathrm{lin}}=1.0$] Comparison of errors versus $\mathrm{Cost}(i,\ell)$ for fixed, uniform and adaptive $N_{\mathrm{reg}}$-schemes.
\textbf{Left}: Error in terms of  $\|\nabla_x\!\left(B(s_{h\tau}^{(i)}) - B(s)\right)\|_{L^2(Q)}$. 
\textbf{Right}: Error in terms of  $\|b(s_{h\tau}^{(i)}) - u_{\text{exact}}\|_{L^2(Q)}$. 
%Linearization error: $\eta^i_{\mathrm{lin}}=\| \nabla (s_h^{(i)} - s_h^{(i-1)}) \bigr\|_{L^2(Q)}$.
}
\label{Figure:1DPME_Newton_m=2_dof}
\end{figure}

%%================================================================
\iffalse
\begin{figure}[h]
    \centering
    \begin{subfigure}[b]{0.49\textwidth}
        \includegraphics[width=\textwidth]       {Figures/EnergyError_vs_Iteration_adaptive_vs_uniform_BH1_dx_Sh_m2_Test_Newton.eps}       
    \end{subfigure}
    \hfill
    \begin{subfigure}[b]{0.49\textwidth}
        \includegraphics[width=\textwidth]
{Figures/EnergyError_vs_Iteration_adaptive_vs_uniform_bbL2_dx_Sh_m2_TEST_Newton.eps}
    \end{subfigure}
    \centering
     \caption{[\Cref{tab:mesh_params_all}, Section \ref{sec:1DPME_1DST}, $m=2$] Comparison of errors versus the number of iterations for uniform and adaptive space-time meshes. 
\textbf{Left}: the energy error $\|\nabla_x\!\left(B(s_{h\tau}^{(i)}) - B(s)\right)\|_{L^2(Q)}$. 
\textbf{Right}: the energy error $\|b(s_{h\tau}^{(i)}) - u_{\text{exact}}\|_{L^2(Q)}$. 
Linearization error: $\| \nabla (s_h^{(i)} - s_h^{(i-1)}) \bigr\|_{L^2(Q)}$}
\label{Figure:1DPME_m=2_iterationno_Newton}
\end{figure}

\fi

%%===============================================================
%======================================================
%\begin{table}[h!]
%\centering
%\caption{Numerical parameters for 1D PME on adaptive and uniform meshes.}
%\label{tab:1DPME_uniform_mesh_ref_vs_adap}
%\begin{minipage}{0.95\textwidth}
%\centering
%\begin{tabular}{lccccccc}
%\toprule
%\textbf{Run} & \textbf{Mesh} & %$n_{\mathrm{space}}$ & $n_{\mathrm{time}}$
%& $m$ & Domain & Time Interval \\
%\midrule
%1D-PME & Adaptive & 16 & 2 & 2 & $[-8,8]$ & $[0,1]$ \\
%1D-PME & Uniform  & 16 & 2 & 2 & $[-8,8]$ & $[0,1]$ \\
%\bottomrule
%\end{tabular}
%\end{minipage}
%\end{table}
\begin{figure}[h]
    \centering
   \includegraphics[width=0.47\linewidth]{Figures/EnergyError_vs_iteration_uniform_bbL2_M2_fixnormb.eps}   
    \includegraphics[width=0.49\linewidth]{Figures/error_vs_dofs_bbL2_M2_fixbnorm_new.eps}
    \caption{[\Cref{tab:mesh_params_all}, Section \ref{sec:1DPME_1DST}, $m=2$] \textbf{Left}: Comparison of errors versus the number of iterations under successive uniform mesh refinement. Each marker corresponds to a refinement step where the mesh size is halved.
\textbf{Right}: Comparison of errors versus the degrees of freedom ($|\mathrm{DOF}|_\ell$) for uniform and adaptive space-time meshes $\calT^\ell_\htt$ as the iterations progress. 
}
    \label{fig:compasion-Uni-Adapt}
\end{figure}

\Cref{fig:compasion-Uni-Adapt} (Right) compares the error decay with respect to the degrees of freedom associated with each mesh level $\ell$ visited, for both the uniform and adaptive strategies. Again, the adaptive strategy shows better performance, demonstrating almost first-order convergence. \Cref{Figure:1DPME_m=4_dof,Figure:1DPME_m=6_dof}
show the decay of errors with respect to costs (Left) and iteration number (Right) for porous medium exponents $m=4$ and $6$ respectively. The gap between adaptive and uniform strategies is even more pronounced in these cases, with the adaptive scheme incurring costs one order of magnitude less. This is due to the sharp interfaces being more localized for higher $m$ values, resulting in the adaptive scheme being more effective.

%%%%%%%%%%%%%%%%%%%%%%%%%%%%%%%%%%%%%%%%%%%%%%%%%%%%%%%%%%%%%



%%%%%%%%%%%%%%%%%%%%%%%%%%%%%%%%%%%%%%%%%%%%%%%%%%%%%%%%%%%%%
%=============================================
%\subsubsubsection*{Results for $m=4$}\label{subsec:m4}
\begin{figure}[h]
    \centering
    \begin{subfigure}[b]{0.49\textwidth}
        \includegraphics[width=\textwidth]
{Figures/EnergyError_vs_DoFs_adaptive_vs_uniform_bbL2_dx_Sh_m4.eps}
        %{Figures/EnergyError_vs_DoFs_adaptive_vs_uniform_BH1_dx_Sh_m4.eps}
    \end{subfigure}
   \hfill
   \begin{subfigure}[b]{0.49\textwidth}
\includegraphics[width=\textwidth]
{Figures/EnergyError_vs_Iteration_adaptive_vs_uniform_bbL2_dx_Sh_m4.eps}
%{Figures/EnergyError_vs_DoFs_adaptive_vs_uniform_bbL2_dx_Sh_m4.eps}
   \end{subfigure}
    \centering
    \caption{[\Cref{tab:mesh_params_all}, Section \ref{sec:1DPME_1DST}, $m=4$]  Comparison of  error $\|b(s_{h\tau}^{(i)}) - u_{\mathrm{exact}}\|_{L^2(Q)}$ 
for fixed, uniform, and adaptive $L$-schemes. 
\textbf{Left:} error versus $\mathrm{Cost}(i,\ell)$. 
\textbf{Right:} error versus $\mathrm{IterNo}(i,\ell)$. 
%Linearization error: $\eta^i_{\mathrm{lin}} = 
%\|\nabla(s_h^{(i)} - s_h^{(i-1)})\|_{L^2(Q)}$.
}
\label{Figure:1DPME_m=4_dof}

\end{figure}

\begin{figure}[h]
    \centering
    \begin{subfigure}[b]{0.49\textwidth}
        \includegraphics[width=\textwidth]{Figures/EnergyError_vs_DoFs_adaptive_vs_uniform_bbL2_dx_Sh_m6.eps}
    \end{subfigure}
    \hfill
    \begin{subfigure}[b]{0.49\textwidth}
        \includegraphics[width=\textwidth]
{Figures/EnergyError_vs_Iteration_adaptive_vs_uniform_bbL2_dx_Sh_m6.eps}
    \end{subfigure}
 \caption{[\Cref{tab:mesh_params_all}, Section \ref{sec:1DPME_1DST}, $m=6$] Comparison of  error $\|b(s_{h\tau}^{(i)}) - u_{\mathrm{exact}}\|_{L^2(Q)}$ 
for fixed, uniform, and adaptive $L$-schemes. 
\textbf{Left:} error versus $\mathrm{Cost}(i,\ell)$. 
\textbf{Right:} error versus $\mathrm{IterNo}(i,\ell)$. 
%Linearization error: $\eta^i_{\mathrm{lin}} = 
%\|\nabla(s_h^{(i)} - s_h^{(i-1)})\|_{L^2(Q)}$..
}
\label{Figure:1DPME_m=6_dof}
\end{figure}
%=============================================
\subsubsection{2D PME}\label{2DPME_2DST}
%\begin{figure}[h]
%  \centering \includegraphics[width=0.23\linewidth]{Figures/Initial_mesh_M2.png}
%  \caption{[\Cref{tab:mesh_params_all,2DPME_2DST}] Uniform space-time mesh in $(x,y,t)$ (same for all $m$).}
%\end{figure}
%================================================
%\begin{figure}[h]
 %  \centering        \includegraphics[width=0.5\textwidth]
%{Figures/EnergyError_vs_DoFs_adaptive_vs_uniform_BBL2_dx_Sh_m6.eps}
 %  \end{figure}
%=================================================

%==================================
% \begin{figure}[h]
% % ==================== m = 6 ====================
% \begin{subfigure}{0.34\textwidth}
%   \centering
%   \includegraphics[width=1.02\linewidth]{Figures/Solution_initial_mesh_test.eps}
%   \caption{Initial condition on space-time mesh in $(x,y,t)$ for $m=6$.}
% \end{subfigure}
% \hfill
% \begin{subfigure}{0.29\textwidth}
%   \centering
%   \includegraphics[width=\linewidth]{Figures/Final_Adap_Mesh.eps}
%   \caption{Final adapted space-time mesh in $(x,y,t)$ for $m=6$.}
% \end{subfigure}
% \hfill
% \begin{subfigure}{0.34\textwidth}
%   \centering
%   \includegraphics[width=1.09\linewidth]{Figures/Solution_Mesh_final_M6_test.eps}
%   \caption{Solution on adapted space-time mesh in $(x,y,t)$ for $m=6$.}
% \end{subfigure}
% \caption{[\Cref{tab:mesh_params_all,2DPME_2DST}] Mesh refinement and numerical solution in the space-time domain $(x,y,t)$ for $m=6$. From left to right: initial condition on the uniform space-time mesh, final adapted mesh, and solution on the adapted mesh.}
% \label{fig:mesh_refinement_m2_m4_m6}
% \end{figure}
Next, we look at numerical results for the $L$-scheme in two space dimensions, i.e., $d=2$. Starting with parameters from \Cref{tab:mesh_params_all} and the Barenblatt initial condition \eqref{eq:Barenblatt},
\Cref{fig:cube_mesh_refinement_3d} (a) shows the solution $u=b(s)$ plotted over the final adapted mesh. 
The time-stacked slices in \Cref{fig:cube_mesh_refinement_3d} (b) show that refinement remains focused around the free boundary at all time levels. The spatial slices along the $x$-axis show that the mesh is fine only in a thin region where the solution has steep changes. The efficiency of the adaptive approach is further confirmed in \Cref{fig:2DPME_DOF_INTERATIONNO_M6}, where we compare the error against the cost and total number of iterations for fixed and adaptive strategies.  To achieve a prescribed accuracy, the adaptive method requires significantly fewer DoFs than both the fixed-mesh and the uniform refinement strategies. It also reaches much lower error levels. Note that uniform refinement becomes prohibitively more expensive for higher dimensional computations. Overall, we find the performance of the adaptive space-time scheme to be excellent, particularly for higher $m$ values and in higher space dimensions.


%======================================================
\iffalse

%===================================
\subsubsubsection{Mesh refinement and numerical solution (case $m=2$)}

\begin{figure}[h]
\centering
% --- Initial mesh ---
\begin{subfigure}{0.45\textwidth}
  \centering
  \includegraphics[width=\linewidth]{Figures/Initial_mesh_M2.png}
  \caption{Initial uniform space-time mesh in $(x,y,t)$.}
\end{subfigure}
\hfill
% --- Final mesh ---
\begin{subfigure}{0.45\textwidth}
  \centering
  \includegraphics[width=\linewidth]{Figures/Final_adapted_mesh_M2.png}
  \caption{Final adapted space-time mesh in $(x,y,t)$.}
\end{subfigure}

\vspace{1em}

% --- Initial solution ---
\begin{subfigure}{0.45\textwidth}
  \centering
  \includegraphics[width=\linewidth]{Figures/Initial_solution_initial_mesh_M2.png}
  \caption{Solution on initial space-time mesh in $(x,y,t)$.}
\end{subfigure}
\hfill
% --- Final refined solution ---
\begin{subfigure}{0.45\textwidth}
  \centering
  \includegraphics[width=\linewidth]{Figures/Final_solution_Final_mesh_M2.png}
  \caption{Solution on adapted space-time mesh in $(x,y,t)$.}
\end{subfigure}

\caption{Mesh refinement and corresponding numerical solutions. 
\textbf{Top}: initial and final meshes. \textbf{Bottom}: solution before and after mesh adaptivity.}
\end{figure}


%===================================
\subsubsubsection{Mesh refinement and numerical solution (case $m=4$)}

\begin{figure}[h]
\centering
% --- Initial mesh ---
\begin{subfigure}{0.45\textwidth}
  \centering
  \includegraphics[width=\linewidth]{Figures/Initial_mesh_M4.png}
  \caption{Initial uniform space-time mesh in $(x,y,t)$.}
\end{subfigure}
\hfill
% --- Final mesh ---
\begin{subfigure}{0.45\textwidth}
  \centering
  \includegraphics[width=\linewidth]{Figures/Final_adapted_mesh_M4.png}
  \caption{Final adapted space-time mesh in $(x,y,t)$.}
\end{subfigure}

\vspace{1em}

% --- Initial solution ---
\begin{subfigure}{0.45\textwidth}
  \centering
  \includegraphics[width=\linewidth]{Figures/Initial_solution_initial_mesh_M4.png}
  \caption{Solution on initial space-time mesh in $(x,y,t)$.}
\end{subfigure}
\hfill
% --- Final refined solution ---
\begin{subfigure}{0.45\textwidth}
  \centering
  \includegraphics[width=\linewidth]{Figures/Final_solution_Final_mesh_M4.png}
  \caption{Solution on adapted space-time mesh in $(x,y,t)$.}
\end{subfigure}

\caption{Mesh refinement and corresponding numerical solutions. 
\textbf{Top}: initial and final meshes. \textbf{Bottom}: solution before and after mesh adaptivity.}
\end{figure}

%==============================================
\subsubsubsection{Mesh refinement and numerical solution (case $m=6$)}

\begin{figure}[h]
\centering
% --- Initial mesh ---
\begin{subfigure}{0.45\textwidth}
  \centering
  \includegraphics[width=\linewidth]{Figures/Initial_Mesh_M6.eps}
  \caption{Initial uniform space-time mesh in $(x,y,t)$.}
  \label{fig:init_mesh}
\end{subfigure}
\hfill
% --- Final mesh ---
\begin{subfigure}{0.45\textwidth}
  \centering
  \includegraphics[width=\linewidth]{Figures/Final_Adap_Mesh.eps}
  \caption{Final adapted space-time mesh in $(x,y,t)$.}
  \label{fig:final_mesh}
\end{subfigure}

\vspace{1em}

% --- Initial solution ---
\begin{subfigure}{0.45\textwidth}
  \centering
  \includegraphics[width=\linewidth]{Figures/Solution_initial_mesh.eps}
  \caption{Solution on initial space-time mesh in $(x,y,t)$.}
  \label{fig:sol_init}
\end{subfigure}
\hfill
% --- Final refined solution ---
\begin{subfigure}{0.45\textwidth}
  \centering
  \includegraphics[width=\linewidth]{Figures/Solution_Mesh_final_M6.eps}
  \caption{Solution on adapted space-time mesh in $(x,y,t)$.}
  \label{fig:sol_final}
\end{subfigure}

\caption{Mesh refinement and corresponding numerical solutions. 
\textbf{Top}: initial and final meshes. \textbf{Bottom}: solution before and after mesh adaptivity.}
\label{fig:mesh_refinement_summary}
\end{figure}
\fi
%================================================
\iffalse
\begin{table}[h]
\centering
\begin{minipage}{0.95\textwidth}
\centering
\caption{Numerical parameters for the 2D PME on an adaptive $16\times16$ mesh, time interval $[0,1]$.}
\label{tab:2DPME_16_T1}
\begin{tabular}{lcccccccc}
\toprule
\textbf{Run} & \textbf{Mesh Type} & $n_{\mathrm{spaceX}}$ & $n_{\mathrm{spaceY}}$ & $n_{\mathrm{time}}$
& $m$ & Domain & Time Interval \\
\midrule
2D-PME & Adaptive & $16$ & $16$ & $4$ & $6$ & $[-8,8]^2$& $[0,1]$ \\
\bottomrule
\end{tabular}
\end{minipage}
\end{table}

\begin{figure}[h]
    \centering  
    \begin{subfigure}[b]{0.49\textwidth}
\includegraphics[width=\textwidth]{Figures/3D_PME_Visualization_M6.png}
 \end{subfigure}
     \hfill
    \begin{subfigure}[b]{0.49\textwidth}
  \includegraphics[width=\textwidth]
{Figures/3D_mesh_slices_diff_Z.png}
    \end{subfigure}
    \caption{[\Cref{tab:mesh_params_all,2DPME_2DST}, $m=6$] \textbf{Left} Adapted space-time mesh. \textbf{Right} 3D visualization of mesh slices stacked along the Z (time) axis.}
    \label{fig:3d_mesh_slices}
\end{figure}
\fi
%======================================================
%\begin{table}[h!]
%\centering
%\begin{minipage}{0.95\textwidth}
%\centering
%\caption{Numerical parameters for the 2D PME on an adaptive space-time mesh.}
%\label{tab:2DPME_16_mesh}
%\begin{tabular}{lcccccccc}
%\toprule
%\textbf{Run} & \textbf{Mesh} & $n_{\mathrm{space}}$ &  $n_{\mathrm{time}}$
%& $m$ & Domain & Time Interval \\
%\midrule
%2D-PME & Adaptive & $16\times16$ & $4$ & $6$ & $[-8,8]^2$& $[0,8]$ \\
%\bottomrule
%\end{tabular}
%\end{minipage}
%\end{table}
\begin{figure}[h]
\centering
\begin{subfigure}[t]{0.34\textwidth}
    \centering
    \includegraphics[width=\linewidth]{Figures/Final_Solution_Mesh_M6_dt_2.png}
    \caption{}
\end{subfigure}
\hfill
\begin{subfigure}[t]{0.32\textwidth}
    \centering
    \includegraphics[width=\linewidth]{Figures/3D_TIME_mesh_slices_Nx_16_Nz_4_1.png}
    \caption{}
\end{subfigure}
\hfill
\begin{subfigure}[t]{0.32\textwidth}
    \centering
    \includegraphics[width=\linewidth]{Figures/X_slices_ladder_fixed_orientation.png}
    \caption{}
    %X_slices_Nx_16_Nz_4.png
\end{subfigure}
\caption{[\Cref{tab:mesh_params_all,2DPME_2DST}, $m=6$, $T=8$] Numerical results for 2D PME. \textbf{(a)} Numerical solution $u=b(s)$ on the final adapted mesh in the $(x,y,t)$-domain. \textbf{(b)} Mesh slices stacked along the $t$-axis, illustrating the adaptive 
refinement across time levels. \textbf{(c)} Spatial mesh slices stacked along the $x$-axis.}
\label{fig:cube_mesh_refinement_3d}
\end{figure}

%===========================================
%\begin{table}[h!]
%\centering
%\begin{minipage}{0.95\textwidth}
%\centering
%\caption{Numerical parameters for the 2D PME on an adaptive space-time mesh.}
%\label{tab:2DPME_32_mesh}
%\begin{tabular}{lcccccccc}
%\toprule
%\textbf{Run} & \textbf{Mesh} & $n_{\mathrm{space}}$ %& $n_{\mathrm{time}}$
%& $m$ & Domain & Time Interval \\
%\midrule
%2D-PME & Adaptive & $32\times32$ & $4$ & $6$ & $[-8,8]^2$& $[0,8]$ \\
%\bottomrule
%\end{tabular}
%\end{minipage}
%\end{table}
%==============================================
%==============================
\iffalse
\begin{figure}[h]
    \centering  
    \begin{subfigure}[b]{0.32\textwidth}
\includegraphics[width=\textwidth]{Figures/FINAL_ADAPTED_MESH_nX_34.png}
 \end{subfigure}
     \hfill
    \begin{subfigure}[b]{0.33\textwidth}
  \includegraphics[width=\textwidth]
{Figures/3D_TIME_mesh_slices_Nx_34_Nz_4.png}
    \end{subfigure}
    \hfill
  \begin{subfigure}[b]{0.33\textwidth}
\includegraphics[width=\linewidth]{Figures/Slices_at_diff_X_Nx_32_Nz_4 (2).png}
\end{subfigure}
    \caption{[\Cref{tab:mesh_params_all,2DPME_2DST}, $m=6$] \textbf{Left} Adapted space-time mesh. \textbf{Middle} 3D visualization of mesh slices stacked along the Z (time) axis. \textbf{Right} 3D visualization of spatial slices stacked along the x-axis. $n_{\mathrm{space}}= 32\times 32$. Time interval = $[0,8]$ }
    \label{fig:spacetime_slices_32}
\end{figure}
\fi
%==============================
% \subsubsubsection{Adaptive vs Uniform: Error Behavior}\label{sub:2D_PME_AdapVsUni}
% %\subsubsubsection*{Result for $m=2$}
% \iffalse
% \begin{figure}[h]
%     \centering
%     \begin{subfigure}[b]{0.49\textwidth}
%         \includegraphics[width=\textwidth]{Figures/EnergyError_vs_DoFs_adaptive_vs_uniform_3D_bbL2_M2.eps}
%     \end{subfigure}
%      \hfill
%     \begin{subfigure}[b]{0.49\textwidth}
%         \includegraphics[width=\textwidth]
% {Figures/EnergyError_vs_Iteration_adaptive_vs_uniform_3D_bbL2_M2.eps}
%     \end{subfigure}
%     \centering
%     \caption{
% [\Cref{tab:mesh_params_all}, Section~\ref{sub:2D_PME_AdapVsUni}, $m=2$] Comparison between the uniform and adaptive space-time meshes. \textbf{Left:} Energy error versus Cost. \textbf{Right:} Energy error versus iteration number.
% The linearization error is measured by $
% \left\| \nabla\!\left( s_h^{(i)} - s_h^{(i-1)} \right) \right\|_{L^2(Q)}$.}
% \label{fig:2DPME_DOF_INTERATIONNO_M2}
% \end{figure}

% \noindent
% %============================================================================

% %\subsubsubsection*{Results for $m=4$}\label{subsec:2Dm4}
% \begin{figure}[h]
%     \centering
%     \begin{subfigure}[b]{0.49\textwidth}
%         \includegraphics[width=\textwidth]{Figures/EnergyError_vs_DoFs_adaptive_vs_uniform_3D_bbL2_m4.eps}
%     \end{subfigure}
%      \hfill
%     \begin{subfigure}[b]{0.49\textwidth}
%         \includegraphics[width=\textwidth]
% {Figures/EnergyError_vs_Iteration_adaptive_vs_uniform_3D_bbL2_m4}
%     \end{subfigure}
% \centering
% \caption{[\Cref{tab:mesh_params_all}, Section~\ref{sub:2D_PME_AdapVsUni}, $m=4$]
% Comparison between the uniform and adaptive space-time meshes. \textbf{Left:} Energy error versus Cost.
% \textbf{Right:} Energy error versus iteration number.
% The linearization error is measured by $
% \left\| \nabla\!\left( s_h^{(i)} - s_h^{(i-1)} \right) \right\|_{L^2(Q)}.$}
% \label{fig:2DPME_DOF_INTERATIONNO_M4}
% \end{figure}
% \fi
% %=========================================
% %\subsubsubsection*{Results for $m=6$}\label{subsec:2Dm6}
\begin{figure}[h]
    \centering
    \begin{subfigure}[b]{0.49\textwidth}
        \includegraphics[width=\textwidth]{Figures/EnergyError_vs_DoFs_3D_PME_M6.eps}
    \end{subfigure}
     \hfill
    \begin{subfigure}[b]{0.49\textwidth}
        \includegraphics[width=\textwidth]
{Figures/EnergyError_vs_Iteration_3cD_PME_M6.eps}
    \end{subfigure}
    \centering
    \caption{[\Cref{tab:mesh_params_all}, Section~\ref{2DPME_2DST}, $m=6$]
Comparison of error $\|b(s_{h\tau}^{(i)}) - u_{\mathrm{exact}}\|_{L^2(Q)}$ for fixed and adaptive $L$-schemes. \textbf{Left:} error versus $\mathrm{Cost}(i,\ell)$. \textbf{Right:} error versus $\mathrm{IterNo}(i,\ell)$.
}
\label{fig:2DPME_DOF_INTERATIONNO_M6}

\end{figure}
%===========================================
\subsection{Biofilm growth model}\label{sec:Biofilm}
We consider the biofilm equation \cite{eberl2001new}
\begin{equation*}
\frac{\partial u}{\partial t} = \Delta \Phi(u) \qquad \text{ where } \qquad \Phi(u):=\left(\frac{u}{1-u}\right)^m \qquad \text{ with exponent }\; m>1.
\end{equation*}
We choose $m=2$ for this section. 
For the initial condition, $\g=0.8$ is additionally chosen in \eqref{eq:Barenblatt} so that $\max_\Om \Phi'(u_0)=200$. For these parameters, $u^* = 0.229083$ in \Cref{lemma:b_B_construction} is computed, satisfying $\Phi'(u^*)=1$. Then, the $b$, $B$ functions in \Cref{lemma:b_B_construction} become
\[
b(s)=
\begin{cases}
0, & s<0,\\
s, & 0\le s\le u^*,\\
\dfrac{(s-u^*+\Phi(u^*))^{1/m}}{1+(s-u^*+\Phi(u^*))^{1/m}}, & s>u^*,
\end{cases}
\quad
B(s)=
\begin{cases}
s, & s<0,\\
\dfrac{s^m}{(1-s)^m}, & 0\le s\le u^*,\\
s-u^*+\Phi(u^*), & s>u^*.
\end{cases}
\]

%\begin{figure}[h]
%\centering
% --- Final 1D space--time mesh ---
%\begin{subfigure}{0.48\textwidth}
%  \centering
%  \includegraphics[width=\linewidth]{Figures/Adapted_mesh_1d_biofilm_new.jpeg}
%  \caption{Final adapted space-time mesh in $(x,t)$.}
%\end{subfigure}
%\hfill
% --- Solution on adapted 1D mesh ---
% \begin{subfigure}{0.5\textwidth}
%   \centering
%   \includegraphics[width=1.1\linewidth]{Figures/Final_adapted_solution_1d_biofilm.jpeg}
%   %\caption{Solution on the adapted space-time mesh.}
% \end{subfigure}
% \caption{[\Cref{sec:Biofilm}, $m=2$, $\gamma=0.8$, $T=4$] Numerical solution $u = b(s)$ on the final adapted space-time mesh.}
% \end{figure}


\begin{figure}[h]
\centering

% --- u = b(s) ---
\begin{subfigure}{0.42\textwidth}
  \centering
  \includegraphics[width=1.05\linewidth]{Figures/Biofilm_Solution1D_u_3D_imposed.png}
\end{subfigure}
\hfill
% --- s ---
\begin{subfigure}{0.49\textwidth}
\centering
\includegraphics[width=1.05\linewidth]{Figures/Solution1D_B_3D.eps}
\end{subfigure}
\caption{[\Cref{sec:Biofilm}, $d=1$, $m=2$, $\gamma=0.8$, $T=4$] 
Three-dimensional space-time profiles of the numerical solution of the 1D biofilm model on the final adapted space-time mesh. \textbf{Left:} $u = b(s)$, \textbf{Right:} $w = B(s)$.}
\label{fig:biofilm_1d_3D}
\end{figure}

\Cref{fig:biofilm_1d_3D} shows the numerical solution $(u,w)$ of the biofilm problem computed over the final adapted mesh. It is seen that most of the refinement happens close to $t=0$ due to steep $w$ gradients. The results in Figure~\ref{fig:adaptive_vs_uniform} illustrate the performance of the three strategies: fixed mesh (L), uniform refinement (Uni-L), and adaptive refinement (Adapt-L), in terms of the total error estimated by $\eta^i$ (since the exact solution is not known).
The fixed-mesh strategy shows a steady, smooth decrease in the total error throughout the computation. 
%This decrease reflects the progressive convergence of the linearization iterations on the fixed mesh. 
However, once the linearization converges, the total error reaches a plateau determined by the discretization error, the best approximation the fixed mesh can deliver. 
%No further reduction is possible without mesh refinement, confirming that mesh refinement is essential for achieving higher accuracy.
For Uni-L, the total error decreases smoothly between refinements. However, at each global refinement step (marked by a star), the solution is interpolated onto the new, finer mesh, temporarily increasing the linearization error. The iterations then converge again on the refined mesh. This produces the visible spikes in the curve at each refinement step. This is substantiated by \Cref{fig:adaptive_vs_uniform_L2norm} which shows error decay when the linearization estimator $\eta_{\mathrm{lin}}^i$ is expressed by $\|s^{i+1}-s^i\|_{L^2(Q)}$ which is less effected by the projection operation. Being devoid of the $\|\nabla (s^{i+1}-s^i)\|_{L^2(Q)}$ and $\|\p_t (s^{i+1}-s^i)\|_{L^2(Q)}$ components causes the total error to decay monotonically in \Cref{fig:adaptive_vs_uniform_L2norm}. For the Adapt-L strategy, a similar behavior is observed. Local refinements are triggered more frequently, as indicated by the circles, see \Cref{fig:adaptive_vs_uniform}. The projection operation causes the spikes at each refinement step. In terms of accuracy for a given computational cost, the Uni-L and Adapt-L strategies are quite similar for this test case.




%===========================================================
\iffalse
\begin{figure}[h]
    \centering
    \begin{subfigure}[b]{0.49\textwidth}
        \includegraphics[width=\textwidth]{Figures/Energy_vs_DoFs_all.eps}
    \end{subfigure}
    \hfill
    \begin{subfigure}[b]{0.49\textwidth}
        \includegraphics[width=\textwidth]
{Figures/Energy_vs_iter_all.eps}
    \end{subfigure}
    \centering
\caption{[\Cref{tab:mesh_params_all,sec:Biofilm}, $m=2$, $\gamma=0.8$, 1D] 
Comparison of total error $\eta^i = \eta^i_{\mathrm{disc}} + L\,\eta^i_{\mathrm{lin}}$ 
for fixed, uniform, and adaptive space-time meshes. 
\textbf{Left:} $\eta^i$ versus $\mathrm{Cost}(i,\ell)$. 
\textbf{Right:} $\eta^i$ versus iteration number. 
Linearization error: $\eta^i_{\mathrm{lin}} = 
\left\|\nabla(s_h^{(i)} - s_h^{(i-1)})\right\|_{L^2(Q)}$.}
\end{figure}
\fi

\begin{figure}[h]
    \centering
    \includegraphics[width=0.48\textwidth]{Figures/ErrTot_vs_DoFs_all_Biofilm_H1_norm.eps}
    \hfill
    \includegraphics[width=0.48\textwidth]{Figures/ErrTot_vs_Iteration_all_Biofilm_H1_norm.eps}
\caption{[\Cref{tab:mesh_params_all,sec:Biofilm}, $d=1$] 
Comparison of total error estimator $\eta^i = \eta^i_{\mathrm{disc}} + L\,\eta^i_{\mathrm{lin}}$ 
for fixed, uniform, and adaptive $L$-schemes. 
\textbf{Left:} $\eta^i$ versus $\mathrm{Cost}(i,\ell)$. 
\textbf{Right:} $\eta^i$ versus $\mathrm{IterNo}(i,\ell)$. }
\label{fig:adaptive_vs_uniform}
\end{figure}




% Each local refinement adds degrees of freedom only 
% where $\eta_{\text{disc}}$ is largest, causing a temporary spike in the total error comparable to that of global refinement. However, because refinements are triggered more often, the total error decreases more rapidly overall, each spike is followed by a faster return to a lower error level. The adaptive strategy therefore reaches a lower final error at comparable total computational cost, confirming that targeting refinement where it is most needed leads to more efficient use of computational resources. In practice, if the computation is stopped at any intermediate stage, the adaptive solution consistently delivers lower total error than both the uniform and fixed mesh strategies, the solution is better resolved at every stage, not only at the end.

%=============================================================

\Cref{fig:adaptive_vs_uniform_Mschemes} shows the error decay for the $N_{\mathrm{reg}}$-schemes. A similar error level is reached as compared to the $L$-schemes, both for the adaptive and uniform strategies. More oscillations are present, similar to \Cref{Figure:1DPME_Newton_m=2_dof}. The spikes are seen again at each refinement step, which get removed if the linearization estimator does not use $L^2$-norms of derivatives. \Cref{fig:adaptive_Lschemes_Biofilm_2D} shows how L-schemes compare when the same experiment is run in two space dimensions ($d=2$). The Adapt-L scheme significantly outperforms the Uni-L scheme since extremely fine resolution is required to capture the very high error close to $t=0$. In fact, in this case, even the fixed mesh scheme shows better performance.



\begin{figure}[h]
    \centering
    \includegraphics[width=0.48\textwidth]{Figures/ErrTot_vs_DoFs_all_Biofilm_L2_Lscheme_norm.eps}
    \hfill
    \includegraphics[width=0.48\textwidth]{Figures/ErrTot_vs_Iteration_all_Biofilm_L2_Lscheme_norm.eps}
\caption{[\Cref{tab:mesh_params_all,sec:Biofilm}, $d=1$] 
Comparison of total error estimator $\eta^i = \eta^i_{\mathrm{disc}} + L\,\eta^i_{\mathrm{lin}}$ 
for fixed, uniform, and adaptive $L$-schemes when linearization error is estimated by $\left\|s_h^{(i)} - s_h^{(i-1)}\right\|_{L^2(Q)}$. 
\textbf{Left:} $\eta^i$ versus $\mathrm{Cost}(i,\ell)$. 
\textbf{Right:} $\eta^i$ versus $\mathrm{IterNo}(i,\ell)$.}
\label{fig:adaptive_vs_uniform_L2norm}
\end{figure}


\begin{figure}[h]
\centering
\includegraphics[width=0.49\textwidth]{Figures/Newton_Biofilm_ErrTot_vs_DoFs_H1norm.eps}
\hfill
\includegraphics[width=0.49\textwidth]{Figures/Newton_Biofilm_ErrTot_vs_Iteration_H1norm.eps}
\caption{
[\Cref{tab:mesh_params_all,sec:Biofilm}, $d=1$] 
Comparison of total error estimator $\eta^i = \eta^i_{\mathrm{disc}} + L\,\eta^i_{\mathrm{lin}}$ 
for fixed, uniform, and adaptive $N_{\mathrm{reg}}$-schemes. 
\textbf{Left:} $\eta^i$ versus $\mathrm{Cost}(i,\ell)$. 
\textbf{Right:} $\eta^i$ versus $\mathrm{IterNo}(i,\ell)$.}
\label{fig:adaptive_vs_uniform_Mschemes}
\end{figure}


%============================================================= for n_time=2 (the comented figures) for both uniform and adaptive spatial mesh==================================
\begin{figure}[h]
    \centering
    \begin{subfigure}[b]{0.5\textwidth}
        \includegraphics[width=1.01\textwidth]
{Figures/Biofilm3D_ErrorTotal_vs_DOFs_H1_3curves_test.eps}
    \end{subfigure}
    \hfill
    \begin{subfigure}[b]{0.49\textwidth}
        \includegraphics[width=1.01\textwidth]
{Figures/Biofilm3D_ErrorTotal_vs_Iteration_H1_3curves.eps}

\end{subfigure}
\centering
\caption{[\Cref{tab:mesh_params_all,sec:Biofilm}, $d=2$] Comparison of total error estimator $\eta^i = \eta^i_{\mathrm{disc}} + L\,\eta^i_{\mathrm{lin}}$ 
for fixed, uniform, and adaptive $L$-schemes. 
\textbf{Left:} $\eta^i$ versus $\mathrm{Cost}(i,\ell)$. 
\textbf{Right:} $\eta^i$ versus $\mathrm{IterNo}(i,\ell)$.}
\label{fig:adaptive_Lschemes_Biofilm_2D}
\end{figure}


% \begin{figure}[h]
% \centering
% % --- Final 1D space--time mesh ---
% \begin{subfigure}[t]{0.48\textwidth}
%   \centering  \includegraphics[width=\linewidth]{Figures/final_solution_2d_biofilm.png}
% \end{subfigure}
% \hfill
% % --- Solution on adapted 1D mesh ---
% \begin{subfigure}[t]{0.45\textwidth}
%   \centering
%   \includegraphics[width=1.1\linewidth]{Figures/3D_mesh_slices_Biofilm}
% \end{subfigure}
% \caption{[\Cref{tab:mesh_params_all,sec:Biofilm}, $m=2$] 2D biofilm model in space-time. \textbf{Left:} numerical solution $u = b(s)$ on the final adapted mesh in the $(x,y)$-plane. \textbf{Right:} mesh slices stacked along the $t$-axis, illustrating the adaptive refinement across time levels. \KM{[[Should be removed in my opinion. Does not add much]]}}
% \end{figure}
% % \subsubsubsection{Adaptive vs Uniform: Error Behavior}
% % %\subsubsubsection*{Result for $m=2$}


\subsection{A double degenerate toy-model}\label{sec:toy}
Finally, we investigate a double degenerate toy-model from \cite{javed2025robust} where $\Phi$ becomes multivalue at $\vs=1$, and there is a degenerate elliptic region (fast-diffusion)  $\{u=\om\}$ of non-zero measure:
\begin{align}\label{eq:toy}
\frac{\partial u}{\partial t} = \Delta w+ \frac{1}{2}\, u , \quad  w \in \Phi( u) \quad \text{ where } \quad \Phi(u) =
\begin{cases}
0 & \text{if } u \leq 0, \\
1-\sqrt{1-u^2} & \text{if } 0 \leq u < 1,\\
 [1, \infty] & \text{if } u = 1.
\end{cases}
\end{align}
The initial condition \eqref{eq:Barenblatt} with $m=2$ and$\g=1.0$ as in the PME case.

For this problem, using \eqref{eq:bBexpression}, the functions  $b,\, B$ can be expressed explicitly. For $u^* = \frac{1}{\sqrt{2}}$,
\[
b(s)=
\begin{cases}
s, & s \le u^*,\\[4pt]
\sqrt{1 - (\sqrt{2} - s)^2}, & u^* < s < \sqrt{2},\\[4pt]
1, &  \text{otherwise},
\end{cases}
\qquad
B(s)=
\begin{cases}
0, & s < 0,\\[4pt]
1 - \sqrt{1 - s^2}, & 0 \le s \le u^*,\\[4pt]
s + 1 - \sqrt{2}, &  \text{otherwise}.
\end{cases}
\]
\Cref{fig:sol_toy_model} shows the computed solution of the toy-model after 8 levels of mesh-adaptation for the one-dimensional space case. The degenerate region $\{u=\om\}$ appears for $t>0$ and grows due to the source term. Mesh adaptation is focused on the free-boundaries and domain boundaries after it intersects the free-boundaries. Unlike previous cases, refinement happens more for later time-points since the gradients near the free-boundaries become steeper. 

\Cref{fig:adaptive_L_toy_1D} shows the decay of the estimator $\eta^i$ for the different refinement strategies of the L-scheme. Both uniform and adaptive strategies work well with a steady decay of error except for spikes when the refinement happens, due to the projection operation. In general, lower error levels are reached with fewer DoFs with the refinement strategies, with Adapt-L strategy performing marginally better, even in one space dimension.

\begin{figure}[h]
\centering
% --- Final 1D space--time mesh ---

\hfill
% --- Solution on adapted 1D mesh ---
% \begin{subfigure}{0.5\textwidth}
%   \centering
%   \includegraphics[width=1.1\linewidth]{Figures/Final_adapted_solution_Lscheme_toymodel_source}
%   \caption{Solution on the adapted space-time mesh.}
% \end{subfigure}
% \hfill
  \begin{subfigure}[b]{0.49\textwidth}
        \centering
        \includegraphics[width=\textwidth]{Figures/ToyModel_u_3D_without_source.png}
 %       \caption{Solution $u = b(s)$ with source term.}
        \label{fig:toy_u_t_4}
        \end{subfigure}       
        \hfill
        \begin{subfigure}{0.49\textwidth}
  \centering
  \includegraphics[width=\linewidth]{Figures/Final_adapted_mesh_Lscheme_toymodel_source}
  \vspace{4em}
%  \caption{Final adapted space-time mesh in $(x,t)$.}
\end{subfigure}
\caption{[\Cref{sec:toy}, $d=1$, $m=2$, $\gamma=1.0$, $T=4$] \textbf{Left:} The numerical solution $u = b(s)$ of the toy-model over the final adapted mesh. \textbf{Left:} Final adapted space-time mesh.}\label{fig:sol_toy_model}
\end{figure}




\begin{figure}[h]
\centering
\includegraphics[width=0.49\textwidth]%{Figures/ErrTot_vs_DoFs_Toymodel_Lscheme_H1_source.eps}
{Figures/ErrTot_vs_DoFs_Toymodel_Lscheme_test_source.eps}
\hfill
\includegraphics[width=0.49\textwidth]%{Figures/ErrTot_vs_Iteration_Toymodel_Lscheme_H1_source.eps}
{Figures/ErrTot_vs_Iteration_Toymodel_Lscheme_test_source.eps}
\caption{[\Cref{tab:mesh_params_all,sec:toy}, $d=1$] Comparison of total error estimator $\eta^i = \eta^i_{\mathrm{disc}} + L\,\eta^i_{\mathrm{lin}}$ 
for fixed, uniform, and adaptive $L$-schemes. 
\textbf{Left:} $\eta^i$ versus $\mathrm{Cost}(i,\ell)$. 
\textbf{Right:} $\eta^i$ versus $\mathrm{IterNo}(i,\ell)$.}
\label{fig:adaptive_L_toy_1D}
\end{figure}

\iffalse
\begin{figure}[h]
    \centering

    % --- With source ---
    \begin{subfigure}[b]{0.5\textwidth}
        \centering
        \includegraphics[width=\textwidth]{Figures/ToyModel_u_3D_with_source.eps}
        \caption{Solution $u = b(s)$ with source term.}
        \label{fig:toy_u}
    \end{subfigure}
    \hfill
    % --- Without source ---
    \begin{subfigure}[b]{0.49\textwidth}
        \centering
        \includegraphics[width=1.1\textwidth]{Figures/ToyModel_u_3D_without_source.eps}
        \caption{Solution $u = b(s)$ without source term.}
        \label{fig:toy_s}
    \end{subfigure}
\caption{[\Cref{sec:Toymodel_1d}, $m=2$, $\g=1.0$] Comparison of numerical solutions of the toy-model in the space--time domain $(x,t)$. 
Left: solution $u = b(s)$ with source term. 
Right: solution $u = b(s)$ without source term. }    
    \label{fig:toy_solution}
\end{figure}
\fi
%==============================================================

\section{Conclusion}\label{sec:conclusion}
In this work, we developed and analyzed an adaptive space-time framework for a class of nonlinear and degenerate parabolic equations. The approach is motivated by the necessity for local time-step refinement to resolve degeneracies in both the time derivative and the 
diffusion operator.

We employed a reformulation of the problem, based on \cite{javed2025robust}, that splits the complex nonlinearity into two simpler nonlinearities. A iterative linearization scheme (L-scheme) was then proposed to solve the nonlinear problem which requires solving only heat equations. It was then proven that this scheme is well-posed and unconditionally convergent irrespective of the initial guess, and even in double-degenerate regimes. The convergence is linear in $L^2$ if the problem is non-degenerate. The convergence properties were confirmed numerically, with linear decay of the linearization error 
observed in the non-degenerate and single degenerate cases. The linearization further allows the total error corresponding to an approximate solution, measured through the dual norm of the residual, to be decomposed into two components: linearization  and discretization errors.  The linearization error is measured by the difference of consecutive iterates, whereas, the discretization error corresponds to the iterative step which is simply the heat equation. 
Evaluation of these error components are performed only on fixed norms, independent of mesh, and free from any dependence on the nonlinear function. This mirrors the orthogonal decomposition of error  that was proven in \cite{mitra2023guaranteed} for elliptic problems.

Next, we considered a space-time conforming finite element method to  discretize the L-scheme, providing us with approximate numerical solutions.
A central contribution of this work is the development of guaranteed, fully- and parallely-computable a posteriori error estimators based on locally equilibrated flux reconstructions on space-time vertex patches. The quality of the estimators are independent of the nonlinearities. Moreover, they are straightforward to compute for both the L-scheme, and a modified Newton scheme that was considered for comparison. Numerical experiments show that they robustly capture the error distribution even for degenerate problems. 

Based on the a posteriori estimator, a fully adaptive solution strategy was considered where the iterations are continued until the linearization error becomes smaller than a fraction of the discretization error. Upon meeting this criteria, a refinement step is performed, based on the distribution of the discretization estimator. The effectiveness of the adaptive strategy was demonstrated through numerical experiments for the porous medium equation, the biofilm growth model, and a double degenerate toy-model. The refinements automatically concentrate along propagating free boundaries and regions of fast dynamics. The adaptive approach achieves the same overall accuracy with lower total computational costs compared to uniform refinements and fixed meshes. The adaptive scheme particularly shines in higher-dimensional spaces and for harsher nonlinearities.

% Future work will focus on extending the present framework to more complex coupled systems and on investigating parallel and large-scale implementations of adaptive space-time refinement strategies.


% The discretization estimator $\eta^i_{\mathrm{disc}}$ and the computable upper bound $\eta^i_{\mathrm{lin}}$ on the linearization 
% error $\mathcal{E}^i_{\mathrm{lin}}$ are both assembled from the discrete iterates $s^i_{h\tau}$ and $s^{i+1}_{h\tau}$ alone, and together provide a rigorous upper bound on the total error $\mathcal{E}(s^i) \leq \eta^i_{\mathrm{disc}} + L\,\eta^i_{\mathrm{lin}}$. This separation between discretization error $\mathcal{E}^{\mathrm{H}}_{\mathrm{disc}}$ and linearization error $\mathcal{E}^i_{\mathrm{lin}}$ is crucial in the adaptive setting: mesh refinement is activated only when $\eta^i_{\mathrm{lin}} \leq \delta_{\mathrm{lin}}\, \eta^i_{\mathrm{disc}}$.





%A central contribution of this work is the development of guaranteed a posteriori error estimators based on locally equilibrated flux reconstructions on space-time vertex patches. These estimators are fully computable and allow for a clear separation between linearization and discretization errors. This separation is crucial in the adaptive setting, as it enables mesh refinement to be activated only when the discretization error dominates, thereby avoiding unnecessary refinement during the iterative linearization process. The effectiveness of the proposed adaptive strategy was demonstrated through numerical experiments comparing adaptive and uniform space-time meshes. In particular, the evolution of the discretization error with respect to the number of degrees of freedom showed that the adaptive approach achieves a given accuracy with significantly fewer unknowns than uniform refinement. Moreover, the observed iteration histories confirm that adaptivity does not compromise the convergence behavior of the underlying iterative scheme. The adaptive meshes were shown to automatically concentrate resolution in regions exhibiting strong space-time variations, such as steep fronts and localized nonlinear effects.


% Overall, the results demonstrate that the proposed adaptive space-time framework provides a reliable and efficient tool for the numerical approximation of nonlinear degenerate parabolic equations. By combining a convergent linearization strategy with guaranteed a posteriori error control, the method achieves accurate solutions at a substantially reduced $\mathrm{Cost}(i,\ell)$ compared to non-adaptive approaches.


\textbf{Acknowledgement} AJ acknowledge the financial support of the  HEC grant: 1(2)/HRD/OSS-III/BATCH-3/2022/HEC/384. The work of ISP was supported by the Research Foundation - Flanders (FWO), project G0A9A25N and the German Research Foundation (DFG) through the SFB 1313, project number 327154368.

%==============================================
%\bibliographystyle{plain}
%\bibliography{bibliography}

\begin{thebibliography}{99}
\bibitem{brenner2017improving}
K. Brenner and C. Canc{\`e}s, \emph{Improving {N}ewton's {Method Performance} by {Parametrization}: {The Case} of the {R}ichards {Equation}}, SIAM Journal on Numerical Analysis, \textbf{55} (4), 1760--1785, (2017), SIAM. \href{https://doi.org/10.1137/16M1083414}
{DOI:10.1137/16M1083414}

\bibitem{evans2022partial}
L. C. Evans, \emph{Partial differential equations}, \textbf{19}, (2022), American mathematical society.

\bibitem{MR2249024}
V. Thom{\'e}e, \emph{Galerkin Finite Element Methods for Parabolic Problems}, \textbf{25}, (2006), Springer, Berlin, Springer Series in Computational Mathematics, 2 ed.. \href{https://doi.org/10.1007/3-540-33122-0}{DOI:10.1007/3-540-33122-0}

\bibitem{vazquez2007porous}
J.L. V{\'a}zquez, \emph{{The Porous Medium Equation: Mathematical Theory}}, (2007), Oxford University Press. \href{https://doi.org/10.1093/acprof:oso/9780198569039.001.0001}{DOI:10.1093/acprof:oso/9780198569039.001.0001}

\bibitem{Malgo}
N. Vohra and M. Peszynska, \emph{Robust conservative scheme and nonlinear solver for phase transitions in heterogeneous permafrost}, Journal of Computational and Applied Mathematics, \textbf{442}, 115719, (2024). \href{https://doi.org/10.1016/j.cam.2023.115719}
{DOI:10.1016/j.cam.2023.115719}

\bibitem{pop2011regularization}
I.S. Pop and B. Schweizer, \emph{Regularization schemes for degenerate {R}ichards equations and outflow conditions}, Mathematical Models and Methods in Applied Sciences, \textbf{21} (08), 1685--1712, (2011), World Scientific. \href{https://doi.org/10.1142/S0218202511005532}{DOI:10.1142/S0218202511005532}

\bibitem{ZouExistence}
W. Zou and W. Wang, \emph{Existence of solutions for some doubly degenerate parabolic
equations with natural growth terms}, Nonlinear Analysis. Theory, Methods \& Applications, \textbf{125}, 150--166, (2015). \href{https://doi.org/10.1016/j.na.2015.05.007}
{DOI:10.1016/j.na.2015.05.007}


\bibitem{javed2025robust}
A. Javed, K. Mitra, and I. S. Pop, \emph{Robust, fast, and adaptive splitting schemes for nonlinear doubly-degenerate diffusion equations}, Math Comp,  Accepted, (2025). \href{https://arxiv.org/abs/2508.07420}{arXiv:2508.07420}


\bibitem{neumuller2013space}
M. Neum{\"u}ller, \emph{Space-time methods: fast solvers and applications}, (2013), Graz University of Technology, Graz, Austria. \href{https://doi.org/10.3217/aw6jx-1fz52}{DOI:10.3217/aw6jx-1fz52}

\bibitem{eriksson1991adaptive}
K. Eriksson and C. Johnson, \emph{{Adaptive Finite Element Methods} for {Parabolic Problems} {I}: {A Linear Model Problem}}, SIAM Journal on Numerical Analysis, \textbf{28} (1), 43--77, (1991). \href{https://doi.org/10.1137/0728003}{DOI:10.1137/0728003}

\bibitem{schmich2008adaptivity}
M. Schmich and B. Vexler, \emph{Adaptivity with Dynamic Meshes for Space-Time Finite Element Discretizations of Parabolic Equations}, SIAM Journal on Scientific Computing, \textbf{30} (1), 369--393, (2008), SIAM. \href{https://doi.org/10.1137/060670468}{DOI:10.1137/060670468}


\bibitem{simon1986compact}
J. Simon, \emph{Compact sets in the space {$L^p(0, T; B)$}}, Annali di Matematica pura ed applicata, \textbf{146} (1), 65--96, (1986), Springer. \href{https://doi.org/10.1007/BF01762360}{DOI:10.1007/BF01762360}

\bibitem{Meidner2012}
D. Meidner and B. Vexler, \emph{Adaptive space-time finite element methods for parabolic optimization problems}, Constrained Optimization and Optimal Control for Partial Differential Equations, 319--348, (2012), Springer Basel, Basel. \href{https://doi.org/10.1007/978-3-0348-0133-1\_18}{DOI:10.1007/978-3-0348-0133-1\_18}

\bibitem{ErnSmearsVohralik2017}
A. Ern, I. Smears, and M. Vohral\'{\i}k, \emph{Guaranteed, Locally Space-Time Efficient, and Polynomial-Degree Robust a Posteriori Error Estimates for High-Order Discretizations of Parabolic Problems}, SIAM Journal on Numerical Analysis, \textbf{55} (6), 2811-2834, (2017). \href{https://doi.org/10.1137/16M1097626}{DOI:10.1137/16M1097626}

\bibitem{gong2012space}
W. Gong, M. Hinze, and Z. J. Zhou, \emph{Space-time finite element approximation of parabolic optimal control problems.}, Journal of Numerical Mathematics, \textbf{20} (2), (2012). \href{https://doi.org/10.1515/jnum-2012-0005}{DOI:10.1515/jnum-2012-0005}

\bibitem{LangerSchafelner+2022+247+266}
U. Langer and A. Schafelner, \emph{Adaptive space--time finite element methods for parabolic optimal control problems}, Journal of Numerical Mathematics, \textbf{30} (4), 247--266, (2022). \href{https://doi.org/10.1515/jnma-2021-0059}{DOI:10.1515/jnma-2021-0059}

\bibitem{Dyja2018}
R. Dyja, B. Ganapathysubramanian, and K. G. van der Zee, \emph{{Parallel-In-Space-Time}, {Adaptive} {Finite} {Element} {Framework} for {Nonlinear} {Parabolic} {Equations}}, SIAM Journal on Scientific Computing, \textbf{40} (3), C283-C304, (2018). \href{https://doi.org/10.1137/16M108985X}{DOI:10.1137/16M108985X}

\bibitem{congreve2026residual}
S. Congreve and H.-G. Shin, \emph{Residual-based a posteriori error estimation for a space-time discontinuous {Galerkin formulation of the Richards} equation}, ESAIM: Mathematical Modelling and Numerical Analysis, \textbf{60} (1), 25--50, (2026), EDP Sciences. \href{ https://doi.org/10.1051/m2an/2025094}{DOI:10.1051/m2an/2025094}


\bibitem{dolejvsi2024error}
V. Dolej{\v{s}}{\'\i}, H-G. Shin, and M. Vlas{\'a}k, \emph{Error Estimates and Adaptivity of the Space-Time Discontinuous {Galerkin} Method for Solving the {Richards} Equation}, Journal of Scientific Computing, \textbf{101} (1), 11, (2024), Springer. \href{ https://doi.org/10.1007/s10915-024-02650-x}{DOI:10.1007/s10915-024-02650-x}


\bibitem{list2016study}
F. List and F. A. Radu, \emph{A study on iterative methods for solving {R}ichards' equation}, Computational Geosciences, \textbf{20}, 341--353, Springer. \href{https://doi.org/10.1007/s10596-016-9566-3}{DOI:10.1007/s10596-016-9566-3}

\bibitem{pop2004mixed}
I.S. Pop, F. Radu, and P. Knabner, \emph{Mixed finite elements for the {Richards'} equation: linearization procedure}, Journal of Computational and Applied Mathematics, \textbf{168} (1), 365-373, (2004). \href{https://doi.org/10.1016/j.cam.2003.04.008}{DOI:10.1016/j.cam.2003.04.008}


\bibitem{van2002mathematical}
M. C. M. Van Loosdrecht, J. J. Heijnen, H. Eberl, J. Kreft, and C. Picioreanu, \emph{Mathematical modelling of biofilm structures}, Antonie van Leeuwenhoek, \textbf{81}, 245--256, Springer. \href{https://doi.org/10.1023/A:1020527020464}{DOI:10.1023/A:1020527020464}

\bibitem{steinbach2015space}
O. Steinbach, \emph{{Space-Time Finite Element Methods} for {Parabolic Problems}}, Computational methods in applied mathematics, \textbf{15} (4), 551--566, (2015), de Gruyter. \href{https://doi.org/10.1515/cmam-2015-0026}{DOI:10.1515/cmam-2015-0026}

\bibitem{payne1960optimal}
L.E. Payne and H.F. Weinberger, \emph{An optimal {Poincar\'e} inequality for convex domains}, Archive for Rational Mechanics and Analysis, \textbf{5} (1), 286--292, (1960), Springer. \href{https://doi.org/10.1007/BF00252910}{DOI:10.1007/BF00252910}

\bibitem{boffi2013mixed}
D. Boffi, F. Brezzi, and M. Fortin, \emph{{Mixed Finite Element Methods} and {Applications}}, \textbf{44}, (2013), Springer. \href{https://doi.org/10.1007/978-3-642-36519-5}{DOI:10.1007/978-3-642-36519-5}

\bibitem{carstensen2014axioms}
C. Carstensen, M. Feischl, M. Page, and D. Praetorius, \emph{Axioms of adaptivity}, Computers \& Mathematics with Applications, \textbf{67} (6), 1195-1253, (2014). \href{https://doi.org/10.1016/j.camwa.2013.12.003}{DOI:10.1016/j.camwa.2013.12.003}

\bibitem{mitra2019modified}
K. Mitra and I.S. Pop, \emph{A modified {L}-scheme to solve nonlinear diffusion problems}, Computers \& Mathematics with Applications, \textbf{77} (6), 1722--1738, (2019), Elsevier. \href{https://doi.org/10.1016/j.camwa.2018.09.042}{DOI:10.1016/j.camwa.2018.09.042}

\bibitem{DOLEJSI2019276}
V. Dolej\v{s}\'{\i} and M. Kuraz
and P.  Solin, \emph{Adaptive higher-order space-time discontinuous {G}alerkin 
method for the computer simulation of variably-saturated 
porous media flows}, Applied Mathematical Modelling, \textbf{72}, 276--305, (2019). \href{https://doi.org/10.1016/j.apm.2019.02.037}{DOI:10.1016/j.apm.2019.02.037}

\bibitem{jayadharan:hal-03355088}
M. Jayadharan, M. Kern, M. Vohral\'ik, and I. Yotov, \emph{A Space-Time Multiscale Mortar Mixed Finite Element Method for Parabolic Equations}, SIAM Journal on Numerical Analysis, \textbf{61} (2), 675-706, (2023). \href{https://doi.org/10.1137/21M1447945}{DOI:10.1137/21M1447945}

\bibitem{schwab2009space}
C. Schwab and R. Stevenson, \emph{Space-time adaptive wavelet methods for parabolic evolution problems}, Mathematics of Computation, \textbf{78} (267), 1293--1318, (2009). \href{https://doi.org/10.1090/S0025-5718-08-02205-9}{DOI:10.1090/S0025-5718-08-02205-9}

\bibitem{GanderNeumuller2016}
M. J. Gander and M. Neum\"uller, \emph{{Analysis} of a {New Space-Time Parallel Multigrid Algorithm} for {Parabolic Problems}}, SIAM Journal on Scientific Computing, \textbf{38} (4), A2173-A2208, (2016). \href{https://doi.org/10.1137/15M1046605}{DOI:10.1137/15M1046605}


\bibitem{Andreev2013}
R. Andreev, \emph{Stability of sparse space--time finite element discretizations of linear parabolic evolution equations}, IMA Journal of Numerical Analysis, \textbf{33} (1), 242-260, (2013). \href{https://doi.org/10.1093/imanum/drs014}
{DOI:10.1093/imanum/drs014}

\bibitem{StevensonWesterdiep2021}
R. Stevenson and J. Westerdiep, \emph{Stability of {Galerkin} discretizations of a mixed space--time variational formulation of parabolic evolution equations}, IMA Journal of Numerical Analysis, \textbf{41} (1), 28-47, (2021). \href{https://doi.org/10.1093/imanum/drz069}{DOI:10.1093/imanum/drz069}


\bibitem{GunzburgerKunoth2011}
M. D. Gunzburger and A. Kunoth, \emph{Space-Time Adaptive Wavelet Methods for Optimal Control Problems Constrained by Parabolic Evolution Equations}, SIAM Journal on Control and Optimization, \textbf{49} (3), 1150-1170, (2011). \href{https://doi.org/10.1137/100806382}
{DOI:10.1137/100806382}

\bibitem{LANGER2016342}
U. Langer, S. E. Moore, and M. Neum\"uller, \emph{Space--time isogeometric analysis of parabolic evolution problems}, Computer Methods in Applied Mechanics and Engineering, \textbf{306}, 342-363, (2016). \href{https://doi.org/10.1016/j.cma.2016.03.042}
{DOI:10.1016/j.cma.2016.03.042}

\bibitem{neumuller2011refinement}
M. Neum{\"u}ller and O. Steinbach, \emph{Refinement of flexible space--time finite element meshes and discontinuous {Galerkin} methods}, Computing and visualization in science, \textbf{14} (5), 189--205, (2011), Springer. \href{https://doi.org/10.1007/s00791-012-0174-z}
{DOI:10.1007/s00791-012-0174-z}

\bibitem{devaud2018space}
D. Devaud and C. Schwab, \emph{Space--time $hp$-approximation of parabolic equations}, Calcolo, \textbf{55} (3), 35, (2018), Springer. \href{https://doi.org/10.1007/s10092-018-0275-2}
{DOI:10.1007/s10092-018-0275-2}


\bibitem{rekatsinas2019optimal}
N. Rekatsinas and R. Stevenson, \emph{An optimal adaptive tensor product wavelet solver of a space-time {FOSLS} formulation of parabolic evolution problems}, Advances in Computational Mathematics, \textbf{45} (2), 1031--1066, (2019), Springer. \href{https://doi.org/10.1007/s10444-018-9644-2}
{DOI:10.1007/s10444-018-9644-2}


\bibitem{SteinbachZank2020}
O. Steinbach and M. Zank, \emph{Coercive space-time finite element methods for initial boundary value problems}, Electronic Transactions on Numerical Analysis, \textbf{52}, 154--194, (2020). \href{https://doi.org/10.1553/etna_vol52s154}
{DOI:10.1553/etna\_vol52s154}


\bibitem{FUHRER202127}
T. F\"uhrer and M. Karkulik, \emph{Space--time least-squares finite elements for parabolic equations}, Computers \& Mathematics with Applications, \textbf{92}, 27-36, (2021). \href{https://doi.org/10.1016/j.camwa.2021.03.004}
{DOI:10.1016/j.camwa.2021.03.004}

\bibitem{BernardiElAlaouiMghazli2014}
C. Bernardi and L. El Alaoui
and Z. Mghazli, \emph{A posteriori analysis of a space and time 
discretization of a nonlinear model for the 
flow in partially saturated porous media}, IMA Journal of Numerical Analysis, \textbf{34} (3), 1002--1036, (2014). \href{https://doi.org/10.1093/imanum/drt014}
{DOI:10.1093/imanum/drt014}


\bibitem{BARON2017104}
V. Baron, Y. Coudi\`ere, and P. Sochala, \emph{Adaptive multistep time discretization and linearization based on a posteriori error estimates for the {R}ichards equation}, Applied Numerical Mathematics, \textbf{112}, 104-125, (2017). \href{https://doi.org/10.1016/j.apnum.2016.10.005}
{DOI:10.1016/j.apnum.2016.10.005}


\bibitem{CLEMENT2021103897}
J.-B. Cl\'ement, F. Golay, M. Ersoy, and D. Sous, \emph{An adaptive strategy for discontinuous {G}alerkin simulations of {R}ichards' equation: Application to multi-materials dam wetting}, Advances in Water Resources, \textbf{151}, 103897, (2021). \href{https://doi.org/10.1016/j.advwatres.2021.103897}
{DOI:10.1016/j.advwatres.2021.103897}

\bibitem{mitra2024posteriori}
K. Mitra and M. Vohral\'ik, \emph{A posteriori error estimates for the 
{R}ichards equation}, Mathematics of Computation, \textbf{93} (347), 1053--1096, (2024). \href{https://doi.org/10.1090/mcom/3932}
{DOI:10.1090/mcom/3932}


\bibitem{DolejsiRoskovec2021}
V. Dolej\v{s}\'{\i}, F. Roskovec, and M. Vlas\'{a}k, \emph{A Posteriori Error Estimates for Higher Order Space-Time {Galerkin} Discretizations of Nonlinear Parabolic Problems}, SIAM Journal on Numerical Analysis, \textbf{59} (3), 1486-1509, (2021). \href{https://doi.org/10.1137/18M117594X}
{DOI:10.1137/18M117594X}


\bibitem{Eriksson1995}
K. Eriksson and C. Johnson, \emph{{Adaptive Finite Element Methods} for {Parabolic Problems} {II}: {Optimal Error Estimates} in {$L_\infty L_2$} and {$L_\infty L_\infty$}}, SIAM Journal on Numerical Analysis, \textbf{32} (3), 706-740, (1995). \href{https://doi.org/10.1137/0732033}
{DOI:10.1137/0732033}

\bibitem{mitra2023guaranteed}
K. Mitra and M. Vohral\'ik, \emph{Guaranteed, locally efficient, and robust a posteriori estimates for nonlinear elliptic problems in iteration-dependent norms. An orthogonal decomposition result based on iterative linearization}, Preprint hal-04156711v2, (2023).



\bibitem{ErnSmearsVohralik2019}
A. Ern, I. Smears, and M. Vohral\'ik, \emph{Equilibrated flux a posteriori error estimates in ${L}^{2}({H}^{1})$-norms for high-order discretizations of parabolic problems}, IMA Journal of Numerical Analysis, \textbf{39} (3), 1158-1179, (2019). \href{https://doi.org/10.1093/imanum/dry035}
{DOI:10.1093/imanum/dry035}

\bibitem{eberl2001new}
H. J. Eberl, D. F. Parker, and M. C. M. {Van Loosdrecht}, \emph{A New Deterministic {Spatio-Temporal} {Continuum} {Model} for {Biofilm} {Development}}, Computational and Mathematical Methods in Medicine, \textbf{3} (3), 429794, (2001). \href{https://doi.org/10.1080/10273660108833072}{DOI:10.1080/10273660108833072}


\bibitem{RattiVerani2019}
L. Ratti and M. Verani, \emph{A posteriori error estimates for the 
monodomain model in cardiac 
electrophysiology}, Calcolo, \textbf{56} (3), 25, (2019). \href{https://doi.org/10.1007/s10092-019-0327-2}
{DOI:10.1007/s10092-019-0327-2}


\bibitem{GeorgoulisLakkisWihler2021}
E. H. Georgoulis, O. Lakkis, and T. P. Wihler, \emph{A posteriori error bounds for 
fully-discrete $hp$-discontinuous 
{G}alerkin timestepping methods for 
parabolic problems}, Numerische Mathematik, \textbf{148} (2), 363--386, (2021). \href{https://doi.org/10.1007/s00211-021-01187-7}
{DOI:10.1007/s00211-021-01187-7}

\bibitem{AntoniettiDede2025}
P. F. Antonietti, L. Ded\`e, G. Loli, M. Montardini, G. Sangalli, and P. Tesini, \emph{Space--time {Isogeometric Analysis} of cardiac electrophysiology}, Computer Methods in Applied Mechanics and Engineering, \textbf{441}, 117957, (2025). \href{https://doi.org/10.1016/j.cma.2025.117957}{DOI:10.1016/j.cma.2025.117957}


\bibitem{berthelin2008model}
F. Berthelin, P. Degond, M. Delitala, and M. Rascle, \emph{{A model} for the {Formation} and {Evolution} of {Traffic Jams}}, Archive for Rational Mechanics and Analysis, \textbf{187} (2), 185--220, (2008), Springer. \href{ https://doi.org/10.1007/s00205-007-0061-9}
{DOI:10.1007/s00205-007-0061-9}


\bibitem{kavcur2006method}
J. Ka{\v{c}}ur, \emph{{Method of Rothe in evolution equations}}, Equadiff 6: Proceedings of the International Conference on Differential Equations and their Applications held in Brno, Czechoslovakia, Aug. 26--30, 1985, 23--34, (2006), Springer. \href{https://doi.org/10.1007/BFb0076049}{DOI:10.1007/BFb0076049}

\bibitem{SmeetsEtAl2024}
R K H Smeets, K Mitra, I S Pop, and S Sonner, \emph{Robust time-discretization and linearization schemes for singular and degenerate evolution systems modelling biofilm growth}, IMA Journal of Numerical Analysis, draf077, (2025). \href{https://doi.org/10.1093/imanum/draf077}{DOI:10.1093/imanum/draf077}


\bibitem{bergamaschi1999mixed}
L. Bergamaschi and M. Putti, \emph{Mixed finite elements and Newton-type linearizations for the solution of {R}ichards' equation}, International journal for numerical methods in engineering, \textbf{45} (8), 1025--1046, (1999), Wiley Online Library. \href{https://doi.org/10.1002/(SICI)1097-0207(19990720)45:8<1025::AID-NME615>3.0.CO;2-G}
{DOI:10.1002/(SICI)1097-0207(19990720)45:8<1025::AID-NME615>3.0.CO;2-G}


\bibitem{lehmann1998comparison}
F. Lehmann and P. Ackerer, \emph{Comparison of iterative methods for improved solutions of the fluid flow equation in partially saturated porous media}, Transport in Porous Media, \textbf{31}, 275--292, (1998), Springer. \href{https://doi.org/10.1023/A:1006555107450}{DOI:10.1023/A:1006555107450}

\bibitem{radu2006newton}
F.A. Radu, I.S. Pop, and P. Knabner, \emph{Newton-type methods for the mixed finite element discretization of some degenerate parabolic equations}, Numerical Mathematics and Advanced Applications: Proceedings of ENUMATH 2005, 1192--1200, (2006), Springer. \href{https://doi.org/10.1007/978-3-540-34288-5_120}{DOI:10.1007/978-3-540-34288-5120}

\bibitem{fevotte2024adaptive}
F. F{\'e}votte, A. Rappaport, and M. Vohral{\'\i}k, \emph{Adaptive regularization for the Richards equation}, Computational Geosciences, \textbf{28} (6), 1371--1388, (2024), Springer. \href{https://doi.org/10.1007/s10596-024-10309-7}{DOI:10.1007/s10596-024-10309-7}

\bibitem{casulli2010nested}
V. Casulli and P. Zanolli, \emph{A Nested Newton-type algorithm for finite volume methods solving Richards' equation in mixed form}, SIAM Journal on Scientific Computing, \textbf{32} (4), 2255--2273, (2010), SIAM. \href{https://doi.org/10.1137/100786320}{DOI:10.1137/100786320}

\bibitem{celia1990general}
M.A. Celia, E.T. Bouloutas, and R.L. Zarba, \emph{A general mass-conservative numerical solution for the unsaturated flow equation}, Water resources research, \textbf{26} (7), 1483--1496, (1990), Wiley Online Library. \href{https://doi.org/10.1029/WR026i007p01483}{DOI:10.1029/WR026i007p01483}


\bibitem{stokke2023adaptive}
J.S. Stokke, K. Mitra, E. Storvik, J.W. Both, and F.A. Radu, \emph{An adaptive solution strategy for {R}ichards' equation}, Computers \& Mathematics with Applications, \textbf{152}, 155-167, (2023). \href{https://doi.org/10.1016/j.camwa.2023.10.020}
{DOI:10.1016/j.camwa.2023.10.020}

\bibitem{kohle2025robust}
RN K{\"o}hle, KTW Menting, K Mitra, and JHM Boonkkamp, \emph{Robust and fast iterative method for the elliptic Monge-Amp\`ere equation}, arXiv preprint arXiv:2509.00794, (2025). \href{
https://doi.org/10.48550/arXiv.2509.00794}{DOI:10.48550/arXiv.2509.00794}

\bibitem{Simplex_sp}
M. von Danwitz, V. Karyofylli, N. Hosters, and M. Behr, \emph{Simplex space-time meshes in compressible flow simulations}, International Journal for Numerical Methods in Fluids, \textbf{91} (1), 29-48, (2019). \href{https://doi.org/10.1002/fld.4743}{DOI:10.1002/fld.4743}

\bibitem{hecht2012new}
F. Hecht, \emph{New development in {FreeFem++}}, Journal of Numerical Mathematics, \textbf{20} (3-4), 1--14, (2012). \href{https://doi.org/10.1515/jnum-2012-0013}
{DOI:10.1515/jnum-2012-0013}

\bibitem{harnist2023robust}
A. Harnist, K. Mitra, A. Rappaport, and M. Vohral{\'i}k, \emph{Robust energy a posteriori estimates for nonlinear elliptic problems}, (2023). working paper or preprint

\bibitem{beranek2025quasioptimal}
N. Beranek, R. Smeets, and R. Stevenson, \emph{Quasi-optimal time-space discretizations for a class of nonlinear parabolic {PDEs}}, SIAM Journal on Numerical Analysis, \textbf{64} (3), 804--827, (2026). \href{https://doi.org/10.1137/25M1797156}{DOI:10.1137/25M1797156}

\bibitem{alt1983quasilinear}
H.W. Alt and S. Luckhaus, \emph{Quasilinear elliptic-parabolic differential equations}, Mathematische Zeitschrift, \textbf{183} (3), 311--341, (1983). \href{https://doi.org/10.1007/BF01176474}{DOI:10.1007/BF01176474}

\bibitem{ANDREIANOV20173633}
B. Andreianov, C. Canc{\`e}s, and A. Moussa, \emph{A nonlinear time compactness result and applications to discretization of degenerate parabolic--elliptic PDEs}, Journal of Functional Analysis, \textbf{273} (12), 3633-3670, (2017). \href{https://doi.org/10.1016/j.jfa.2017.08.010}
{DOI:10.1016/j.jfa.2017.08.010}


\bibitem{CarilloEntropy}
J. Carrillo, \emph{Entropy Solutions for Nonlinear Degenerate Problems}, Archive for Rational Mechanics and Analysis, \textbf{147} (4), 269-361, (1999). \href{https://doi.org/10.1007/s002050050152}
{DOI:10.1007/s002050050152}

\bibitem{CARRILLO199993}
J. Carrillo and P. Wittbold, \emph{Uniqueness of Renormalized Solutions of Degenerate Elliptic--Parabolic Problems}, Journal of Differential Equations, \textbf{156} (1), 93-121, (1999). \href{https://doi.org/10.1006/jdeq.1998.3597}
{DOI:10.1006/jdeq.1998.3597}

\bibitem{DroniouEymardTalbot}
J. Droniou, R. Eymard, and K.S. Talbot, \emph{Convergence in {$C([0,T];L^2(\Omega))$} of weak solutions to
perturbed doubly degenerate parabolic equations}, Journal of Differential Equations, \textbf{260} (11), 7821--7860, (2016). \href{https://doi.org/10.1016/j.jde.2016.02.004}
{DOI:10.1016/j.jde.2016.02.004}


\bibitem{LuJaeger}
Y. Lu and W. J\"ager, \emph{On solutions to nonlinear reaction-diffusion-convection
equations with degenerate diffusion}, J. Differential Equations, \textbf{170}, 1--21, (2001).
\href{https://doi.org/10.1006/jdeq.2000.3800}
{DOI:10.1006/jdeq.2000.3800}


\bibitem{Jan2012}
J. \v{C}esenek and M. Feistauer, \emph{Theory of the {Space-Time} {Discontinuous} {Galerkin} {Method for Nonstationary Parabolic Problems with Nonlinear Convection and Diffusion}}, SIAM Journal on Numerical Analysis, \textbf{50} (3), 1181-1206, (2012). \href{https://doi.org/10.1137/110828903}{DOI:10.1137/110828903}

\end{thebibliography}

\end{document}
